% !TEX program = pdflatex
\documentclass[12pt,a4paper]{report}

% =========================
% CẤU HÌNH CƠ BẢN CHO PDFLATEX
% =========================
\usepackage[utf8]{inputenc}
\usepackage[T5]{fontenc}
\usepackage[vietnamese]{babel}
\usepackage[a4paper,margin=2.5cm]{geometry}

\usepackage{amsmath,amssymb}
\usepackage{graphicx}
\usepackage{array}
\usepackage{longtable}
\usepackage{booktabs}
\usepackage{multirow}
\usepackage{enumitem}
\usepackage{xcolor}
\usepackage{hyperref}
\usepackage{tikz}
\usepackage[most]{tcolorbox}
\usetikzlibrary{arrows.meta,positioning,calc,fit}

% =========================
% VIỆT HÓA MỘT SỐ TÊN MỤC
% =========================
\renewcommand{\contentsname}{Mục lục}
\renewcommand{\chaptername}{Chương}
\renewcommand{\bibname}{Tài liệu tham khảo}
\renewcommand{\figurename}{Hình}
\renewcommand{\tablename}{Bảng}

% =========================
% FIX TOC: increase space for 3-digit page numbers
% =========================
\makeatletter
\renewcommand{\@pnumwidth}{2em}  % default 1.55em, increase for 3-digit pages
\renewcommand{\@tocrmarg}{3em}   % right margin
\makeatother

% =========================
% ĐỊNH DẠNG HYPERLINK
% =========================
\hypersetup{
    colorlinks=true,
    linkcolor=blue!60!black,
    urlcolor=blue!60!black,
    citecolor=blue!60!black
}

% =========================
% HỘP GHI CHÚ
% =========================
\newtcolorbox{ghinho}{
    colback=blue!4,
    colframe=blue!60!black,
    title=Ghi nhớ,
    fonttitle=\bfseries,
    arc=2mm
}

\newtcolorbox{vidu}{
    colback=green!4,
    colframe=green!50!black,
    title=Ví dụ minh họa,
    fonttitle=\bfseries,
    arc=2mm
}

\newtcolorbox{luuy}{
    colback=orange!6,
    colframe=orange!70!black,
    title=Lưu ý,
    fonttitle=\bfseries,
    arc=2mm
}

\newtcolorbox{canhbao}{
    colback=red!4,
    colframe=red!60!black,
    title=Cảnh báo,
    fonttitle=\bfseries,
    arc=2mm
}

% =========================
% THÔNG TIN TÀI LIỆU
% =========================
\title{
    \textbf{Trí tuệ Nhân tạo - Học Máy - Học sâu}\\
    \textbf{AI - ML - DL}\\
    \large Artificial Intelligence \& Machine Learning
}
\author{Biên soạn tổng hợp}
\date{\today}

\begin{document}

\maketitle
\tableofcontents
\newpage

% ==========================================================
% PHẦN I: MẠNG NƠ-RON NHÂN TẠO CƠ BẢN
% ==========================================================
\part{Mạng Nơ-ron Nhân tạo Cơ bản}

\chapter{ANN1: Neural Networks Introduction -- Giới thiệu mạng nơ-ron}
\input{chapters/ann1.tex}

\chapter{ANN3: Forward Propagation -- Lan truyền tiến}


% ==========================================================
\section*{Learning Objectives -- Mục tiêu bài học}
\addcontentsline{toc}{section}{Mục tiêu bài học}

Bài ANN3 tập trung vào quá trình huấn luyện mạng nơ-ron nhân tạo. Sau khi học xong, người học cần nắm được:

\begin{enumerate}[label=\arabic*.]
    \item Vì sao mạng nơ-ron ban đầu chưa có khả năng suy luận.
    \item Ý nghĩa của việc khởi tạo ngẫu nhiên trọng số và bias.
    \item Quá trình lan truyền tiến từ input đến output.
    \item Ý tưởng tổng quát của lan truyền ngược sai số, tiếng Anh là \textit{backpropagation}.
    \item Vai trò của dữ liệu, nhãn, output dự đoán và sai số.
    \item Cách biểu diễn nhãn phân loại bằng one-hot vector.
    \item Một số cách đo sai số giữa hai vector: chuẩn $L_2$, góc giữa hai vector và cross entropy.
    \item Vì sao huấn luyện mạng nơ-ron là quá trình lặp nhiều lần và có thể tốn nhiều thời gian.
\end{enumerate}

% ==========================================================
\section{Supervised Learning Recap -- Nhắc lại bài toán học có giám sát}

\subsection{Mạng nơ-ron cần học ánh xạ}

Trong các bài trước, ta đã biết mạng nơ-ron nhân tạo được dùng để xấp xỉ một hàm ánh xạ:

\[
    f: X \rightarrow Y
\]

Trong đó:

\begin{itemize}
    \item $X$ là không gian dữ liệu đầu vào;
    \item $Y$ là không gian kết quả đầu ra;
    \item $f$ là quan hệ mà mô hình cần học.
\end{itemize}

Ví dụ:

\[
    \text{Ảnh con chó} \rightarrow \text{Nhãn ``chó''}
\]

hoặc:

\[
    \text{Giá cổ phiếu 5 ngày gần nhất} \rightarrow \text{Giá cổ phiếu ngày tiếp theo}
\]

\begin{ghinho}
Ban đầu, mạng nơ-ron chưa biết hàm $f$. Quá trình huấn luyện chính là quá trình điều chỉnh các tham số bên trong mạng để mạng học được ánh xạ từ input sang output.
\end{ghinho}

\subsection{Dữ liệu trong học có giám sát}

Trong học có giám sát, tiếng Anh là \textit{supervised learning}, ta cung cấp cho hệ thống nhiều cặp dữ liệu:

\[
    (x_i, y_i)
\]

Trong đó:

\begin{itemize}
    \item $x_i$ là dữ liệu đầu vào, còn gọi là input hoặc data;
    \item $y_i$ là nhãn đúng, còn gọi là label hoặc target.
\end{itemize}

Tập dữ liệu huấn luyện có dạng:

\[
    D = \{(x_1,y_1),(x_2,y_2),\ldots,(x_N,y_N)\}
\]

Mục tiêu của mô hình là tạo ra dự đoán:

\[
    \hat{y}_i = f(x_i)
\]

sao cho $\hat{y}_i$ càng gần $y_i$ càng tốt.

\begin{vidu}
Bài toán phân loại ảnh gồm 4 lớp: chó, mèo, diều và xe hơi.

\begin{center}
\begin{tabular}{lll}
\toprule
Mẫu dữ liệu & Input & Label đúng \\
\midrule
1 & Ảnh con chó & Chó \\
2 & Ảnh con mèo & Mèo \\
3 & Ảnh con diều & Diều \\
4 & Ảnh xe hơi & Xe hơi \\
\bottomrule
\end{tabular}
\end{center}

Mạng nơ-ron phải học từ nhiều ảnh đã được gán nhãn. Sau khi huấn luyện, khi gặp ảnh mới, mạng dự đoán ảnh thuộc lớp nào.
\end{vidu}

% ==========================================================
\section{Initial Network Parameters -- Thông số ban đầu của mạng nơ-ron}

\subsection{Tham số của từng nơ-ron}

Một nơ-ron nhân tạo nhận nhiều đầu vào:

\[
    x_1, x_2, \ldots, x_n
\]

Mỗi đầu vào có một trọng số tương ứng:

\[
    w_1, w_2, \ldots, w_n
\]

Ngoài ra, nơ-ron còn có bias $b$. Giá trị trước hàm kích hoạt là:

\[
    z = \sum_{i=1}^{n} w_i x_i + b
\]

Output của nơ-ron là:

\[
    a = \varphi(z)
\]

Trong đó $\varphi$ là hàm kích hoạt.

\subsection{Mạng có rất nhiều tham số}

Khi nhiều nơ-ron được kết nối thành một mạng, mỗi nơ-ron có bộ trọng số và bias riêng. Tổng số tham số trong mạng có thể rất lớn.

Ví dụ, nếu một nơ-ron ở lớp hiện tại nhận kết nối từ 5 nơ-ron ở lớp trước, thì nó có:

\[
    5 \text{ trọng số} + 1 \text{ bias} = 6 \text{ tham số}
\]

Nếu mạng có hàng triệu kết nối, tổng số tham số có thể lên đến vài chục triệu hoặc vài trăm triệu.

\begin{ghinho}
Huấn luyện mạng nơ-ron thực chất là tìm bộ giá trị phù hợp cho tất cả các trọng số và bias trong mạng.
\end{ghinho}

\subsection{Khởi tạo ngẫu nhiên}

Khi bắt đầu huấn luyện, mạng chưa biết gì về bài toán. Vì vậy, các trọng số và bias thường được khởi tạo bằng các giá trị ngẫu nhiên.

Ký hiệu bộ tham số của mạng là:

\[
    \theta = \{W,b\}
\]

Ban đầu:

\[
    \theta = \theta_0
\]

Trong đó $\theta_0$ là bộ tham số khởi tạo ngẫu nhiên.

\begin{luuy}
Trong thực tế, việc khởi tạo không phải là ngẫu nhiên hoàn toàn tùy tiện. Có nhiều kỹ thuật khởi tạo đặc biệt để giúp quá trình huấn luyện ổn định hơn. Tuy nhiên, ở mức nhập môn, ta có thể hiểu đơn giản rằng mạng bắt đầu từ một trạng thái ngẫu nhiên.
\end{luuy}

\subsection{Vì sao output ban đầu thường sai?}

Do các trọng số ban đầu là ngẫu nhiên, nên khi đưa input vào mạng, output cũng thường là một kết quả không có ý nghĩa.

\begin{vidu}
Giả sử đưa ảnh con chó vào mạng.

\[
    x = \text{ảnh con chó}
\]

Label đúng là:

\[
    y = \text{chó}
\]

Nhưng vì mạng mới khởi tạo ngẫu nhiên, nó có thể dự đoán:

\[
    \hat{y} = \text{máy bay}
\]

Kết quả này sai vì mạng chưa được huấn luyện.
\end{vidu}

% ==========================================================
\section{Forward Propagation -- Lan truyền tiến}

\subsection{Ý tưởng của lan truyền tiến}

Lan truyền tiến, tiếng Anh là \textit{forward propagation}, là quá trình đưa dữ liệu từ lớp input đi qua các lớp ẩn và cuối cùng sinh ra output.

\[
    \text{Input} \rightarrow \text{Lớp ẩn 1} \rightarrow \text{Lớp ẩn 2} \rightarrow \cdots \rightarrow \text{Output}
\]

Ở mỗi nơ-ron, ta tính:

\[
    z = \sum_{i=1}^{n} w_i x_i + b
\]

sau đó đưa qua hàm kích hoạt:

\[
    a = \varphi(z)
\]

Output của lớp trước trở thành input của lớp sau.

\subsection{Sơ đồ lan truyền tiến}

\begin{center}
\begin{tikzpicture}[x=2.2cm,y=1.0cm,>=Stealth,
    neuron/.style={circle,draw,minimum size=8mm},
    every node/.style={font=\small}
]

% Input layer
\node[neuron] (i1) at (0,1) {$x_1$};
\node[neuron] (i2) at (0,0) {$x_2$};
\node[neuron] (i3) at (0,-1) {$x_3$};

% Hidden layer 1
\node[neuron] (h11) at (1.7,1.2) {$a_1^{(1)}$};
\node[neuron] (h12) at (1.7,0.4) {$a_2^{(1)}$};
\node[neuron] (h13) at (1.7,-0.4) {$a_3^{(1)}$};
\node[neuron] (h14) at (1.7,-1.2) {$a_4^{(1)}$};

% Hidden layer 2
\node[neuron] (h21) at (3.4,0.8) {$a_1^{(2)}$};
\node[neuron] (h22) at (3.4,0) {$a_2^{(2)}$};
\node[neuron] (h23) at (3.4,-0.8) {$a_3^{(2)}$};

% Output layer
\node[neuron] (o1) at (5.1,0.4) {$\hat{y}_1$};
\node[neuron] (o2) at (5.1,-0.4) {$\hat{y}_2$};

% Connections
\foreach \i in {i1,i2,i3}
    \foreach \h in {h11,h12,h13,h14}
        \draw[->] (\i) -- (\h);

\foreach \h in {h11,h12,h13,h14}
    \foreach \k in {h21,h22,h23}
        \draw[->] (\h) -- (\k);

\foreach \k in {h21,h22,h23}
    \foreach \o in {o1,o2}
        \draw[->] (\k) -- (\o);

% Labels
\node at (0,2) {\textbf{Input}};
\node at (1.7,2) {\textbf{Lớp ẩn 1}};
\node at (3.4,2) {\textbf{Lớp ẩn 2}};
\node at (5.1,2) {\textbf{Output}};

\end{tikzpicture}
\end{center}

\subsection{Lan truyền tiến qua nhiều lớp}

Giả sử mạng có $L$ lớp. Với lớp thứ $l$, ta có:

\[
    z^{(l)} = W^{(l)}a^{(l-1)} + b^{(l)}
\]

\[
    a^{(l)} = \varphi \left(z^{(l)}\right)
\]

Trong đó:

\begin{itemize}
    \item $a^{(0)} = x$ là input ban đầu;
    \item $W^{(l)}$ là ma trận trọng số của lớp $l$;
    \item $b^{(l)}$ là vector bias của lớp $l$;
    \item $a^{(l)}$ là output của lớp $l$;
    \item output cuối cùng của mạng là $\hat{y}=a^{(L)}$.
\end{itemize}

\begin{vidu}
Xét một mạng đơn giản:

\[
    x \rightarrow \text{lớp ẩn} \rightarrow \hat{y}
\]

Input:

\[
    x = 
    \begin{bmatrix}
    1\\
    2
    \end{bmatrix}
\]

Giả sử lớp ẩn có 2 nơ-ron, dùng hàm kích hoạt tuyến tính để tính đơn giản:

\[
    W^{(1)}=
    \begin{bmatrix}
    1 & 0\\
    0.5 & 1
    \end{bmatrix},
    \qquad
    b^{(1)}=
    \begin{bmatrix}
    0\\
    1
    \end{bmatrix}
\]

Ta tính:

\[
    z^{(1)} = W^{(1)}x+b^{(1)}
\]

\[
    z^{(1)}
    =
    \begin{bmatrix}
    1 & 0\\
    0.5 & 1
    \end{bmatrix}
    \begin{bmatrix}
    1\\
    2
    \end{bmatrix}
    +
    \begin{bmatrix}
    0\\
    1
    \end{bmatrix}
    =
    \begin{bmatrix}
    1\\
    3.5
    \end{bmatrix}
\]

Nếu hàm kích hoạt là tuyến tính, output lớp ẩn là:

\[
    a^{(1)}=
    \begin{bmatrix}
    1\\
    3.5
    \end{bmatrix}
\]
\end{vidu}

% ==========================================================
\section{Error \& Training Concept -- Sai số và ý tưởng huấn luyện}

\subsection{So sánh output dự đoán với output mong muốn}

Sau khi lan truyền tiến, mạng tạo ra output dự đoán:

\[
    \hat{y}
\]

Trong học có giám sát, ta biết output đúng:

\[
    y
\]

Sai số được hiểu là sự khác biệt giữa $\hat{y}$ và $y$.

\[
    \text{Sai số} = \text{khác biệt giữa output dự đoán và label đúng}
\]

\begin{vidu}
Nếu ảnh đầu vào là ảnh con chó, label đúng là:

\[
    y = \text{chó}
\]

Nhưng mạng dự đoán:

\[
    \hat{y} = \text{máy bay}
\]

Thì mạng đã sai. Ta cần đo mức độ sai đó bằng một hàm mất mát.
\end{vidu}

\subsection{Hàm mất mát}

Hàm mất mát, tiếng Anh là \textit{loss function}, là hàm dùng để đo sai số giữa output dự đoán và output mong muốn.

\[
    \mathcal{L}(\hat{y},y)
\]

Hàm mất mát phải trả về một số vô hướng:

\[
    \mathcal{L} \in \mathbb{R}
\]

Số này càng nhỏ thì dự đoán càng gần với nhãn đúng.

\begin{ghinho}
Hàm mất mát biến khái niệm ``mô hình dự đoán sai'' thành một con số cụ thể. Nhờ đó, thuật toán huấn luyện có thể biết cần điều chỉnh tham số theo hướng nào.
\end{ghinho}

\subsection{Huấn luyện là giảm sai số}

Mục tiêu của huấn luyện là tìm bộ tham số $\theta$ sao cho loss nhỏ nhất có thể:

\[
    \theta^* = \arg\min_{\theta} \mathcal{L}(\hat{y},y)
\]

Với toàn bộ tập dữ liệu:

\[
    J(\theta)=\frac{1}{N}\sum_{i=1}^{N}\mathcal{L}\left(f(x_i;\theta),y_i\right)
\]

Mục tiêu là:

\[
    \theta^* = \arg\min_{\theta} J(\theta)
\]

% ==========================================================
\section{Lan truyền ngược sai số}

\subsection{Backpropagation là gì?}

Lan truyền ngược sai số, tiếng Anh là \textit{backpropagation}, là kỹ thuật dùng sai số ở output để điều chỉnh các tham số bên trong mạng.

Quá trình tổng quát:

\[
    \text{Input} \rightarrow \text{Lan truyền tiến} \rightarrow \text{Output dự đoán}
\]

\[
    \text{So sánh với label} \rightarrow \text{Tính loss}
\]

\[
    \text{Lan truyền ngược sai số} \rightarrow \text{Cập nhật trọng số và bias}
\]

\begin{ghinho}
Backpropagation là một trong những kỹ thuật phổ biến nhất để huấn luyện mạng nơ-ron hiện nay, nhưng không phải là kỹ thuật duy nhất.
\end{ghinho}

\subsection{Ý tưởng trực quan}

Ban đầu, mạng giống như một đứa trẻ chưa biết làm gì. Các trọng số được khởi tạo ngẫu nhiên nên mạng trả lời sai. Mỗi khi mạng sai, ta dùng mức độ sai đó để chỉnh lại các trọng số một chút.

Sau rất nhiều lần chỉnh như vậy, mạng dần học được cách tạo output gần với label đúng.

\begin{vidu}
Quá trình học qua ba mẫu dữ liệu:

\begin{enumerate}
    \item Đưa ảnh con chó vào. Mạng đoán ``máy bay''. So với label ``chó'', mạng sai. Ta cập nhật trọng số.
    \item Đưa ảnh con mèo vào. Mạng đoán ``xe tải''. So với label ``mèo'', mạng sai. Ta tiếp tục cập nhật trọng số.
    \item Đưa ảnh con voi vào. Mạng đoán ``ngôi nhà''. So với label ``voi'', mạng sai. Ta lại cập nhật trọng số.
\end{enumerate}

Lặp lại quá trình này với rất nhiều mẫu dữ liệu, mạng có cơ hội học dần quan hệ giữa ảnh và nhãn.
\end{vidu}

\subsection{Cập nhật từng bước}

Giả sử tại bước huấn luyện thứ $t$, bộ tham số của mạng là:

\[
    \theta_t
\]

Sau khi tính loss, thuật toán cập nhật tham số:

\[
    \theta_{t+1} = \theta_t + \Delta \theta_t
\]

Trong thực tế, nếu dùng gradient descent, công thức thường có dạng:

\[
    \theta_{t+1} = \theta_t - \eta \nabla_{\theta} J(\theta_t)
\]

Trong đó:

\begin{itemize}
    \item $\eta$ là tốc độ học, tiếng Anh là \textit{learning rate};
    \item $\nabla_{\theta} J(\theta_t)$ là gradient của hàm mất mát theo tham số;
    \item dấu trừ thể hiện việc đi theo hướng làm giảm loss.
\end{itemize}

\begin{luuy}
Trong transcript, giảng viên chủ yếu trình bày ý tưởng tổng quát: dùng sai số giữa output hiện tại và output mong muốn để quay lại điều chỉnh các trọng số. Để lập trình được backpropagation, cần học thêm các chi tiết toán học như đạo hàm riêng, gradient, chain rule và thuật toán tối ưu.
\end{luuy}

\subsection{Liên hệ với điều khiển tự động}

Ý tưởng dùng sai lệch giữa kết quả mong muốn và kết quả hiện tại để điều chỉnh hệ thống có liên hệ với tư duy trong điều khiển tự động.

Trong điều khiển phản hồi:

\[
    e = r - y
\]

Trong đó:

\begin{itemize}
    \item $r$ là giá trị mong muốn;
    \item $y$ là giá trị thực tế;
    \item $e$ là sai lệch.
\end{itemize}

Hệ thống dùng sai lệch $e$ để điều chỉnh tín hiệu điều khiển.

Trong huấn luyện mạng nơ-ron:

\[
    \text{sai số} = \text{label đúng} - \text{output dự đoán}
\]

Mạng dùng sai số này để điều chỉnh trọng số và bias.

% ==========================================================
\section{Quá trình huấn luyện lặp}

\subsection{Huấn luyện qua từng mẫu}

Với mỗi cặp dữ liệu $(x_i,y_i)$, ta làm các bước:

\begin{enumerate}
    \item Đưa $x_i$ vào mạng.
    \item Lan truyền tiến để tính $\hat{y}_i$.
    \item So sánh $\hat{y}_i$ với $y_i$.
    \item Tính loss.
    \item Lan truyền ngược sai số.
    \item Cập nhật trọng số và bias.
\end{enumerate}

Sau đó chuyển sang mẫu tiếp theo.

\subsection{Epoch}

Một epoch là một lượt mô hình đi qua toàn bộ tập dữ liệu huấn luyện.

Nếu tập dữ liệu có $N$ mẫu, thì sau một epoch, mạng đã được huấn luyện trên $N$ mẫu đó một lần.

Thông thường, ta cần nhiều epoch:

\[
    \text{Training} = \text{nhiều epoch}
\]

\subsection{Vì sao cần lặp nhiều lần?}

Một lần cập nhật thường chỉ làm thay đổi tham số một chút. Vì vậy, mạng cần rất nhiều lần cập nhật để học được một quan hệ phức tạp.

\begin{vidu}
Nếu mạng có hàng triệu tham số và dữ liệu gồm hàng trăm nghìn ảnh, việc chỉ học một vài ảnh là không đủ. Mạng cần nhìn thấy nhiều ảnh chó, nhiều ảnh mèo, nhiều ảnh xe hơi trong nhiều điều kiện khác nhau để học được đặc trưng tổng quát.
\end{vidu}

\subsection{Thời gian huấn luyện}

Thời gian huấn luyện phụ thuộc vào:

\begin{itemize}
    \item kích thước tập dữ liệu;
    \item độ phức tạp của mạng;
    \item số lượng tham số;
    \item số epoch;
    \item cấu hình phần cứng;
    \item loại GPU hoặc card màn hình;
    \item cách cài đặt thuật toán.
\end{itemize}

\begin{luuy}
Huấn luyện mạng nơ-ron có thể mất vài giờ, vài ngày hoặc thậm chí vài tuần. Với các mô hình lớn, phần cứng mạnh, đặc biệt là GPU, đóng vai trò rất quan trọng.
\end{luuy}

% ==========================================================
\section{Bài toán phân loại nhiều lớp}

\subsection{Khái niệm class}

Trong bài toán phân loại, \textit{class} là lớp dữ liệu hoặc nhóm nhãn mà mô hình cần dự đoán.

Ví dụ bài toán có 4 class:

\[
    \{\text{chó}, \text{mèo}, \text{diều}, \text{xe hơi}\}
\]

Mỗi ảnh đầu vào được giả sử thuộc đúng một class.

\begin{ghinho}
Nếu bài toán yêu cầu mỗi mẫu chỉ thuộc một lớp duy nhất, ta gọi đó là bài toán phân loại một nhãn, tiếng Anh là \textit{single-label classification}.
\end{ghinho}

\subsection{Thiết kế output cho bài toán nhiều lớp}

Với 4 class, có nhiều cách thiết kế output.

\subsection{Cách 1: Một output nhận giá trị 1, 2, 3, 4}

Ta có thể quy ước:

\[
    1 = \text{chó}, \quad
    2 = \text{mèo}, \quad
    3 = \text{diều}, \quad
    4 = \text{xe hơi}
\]

Khi đó mạng chỉ có một output.

Tuy nhiên, cách này có vấn đề: các con số 1, 2, 3, 4 có thứ tự số học, nhưng các class không nhất thiết có quan hệ thứ tự. Ví dụ, không có nghĩa là ``xe hơi'' lớn hơn ``chó'' theo nghĩa toán học.

\subsection{Cách 2: Hai output theo mã nhị phân}

Với 4 class, ta có thể dùng 2 bit:

\[
    00,\quad 01,\quad 10,\quad 11
\]

Ví dụ:

\[
    00 = \text{chó}, \quad
    01 = \text{mèo}, \quad
    10 = \text{diều}, \quad
    11 = \text{xe hơi}
\]

Cách này tiết kiệm số output hơn nhưng không phải cách phổ biến nhất trong bài toán phân loại cơ bản.

\subsection{Cách 3: Bốn output theo one-hot vector}

Cách phổ biến là dùng số output bằng số class. Với 4 class, ta dùng 4 output.

Mỗi label được biểu diễn bằng một vector có đúng một phần tử bằng 1, các phần tử còn lại bằng 0.

\[
    \text{chó} =
    \begin{bmatrix}
    1\\0\\0\\0
    \end{bmatrix},
    \quad
    \text{mèo} =
    \begin{bmatrix}
    0\\1\\0\\0
    \end{bmatrix},
    \quad
    \text{diều} =
    \begin{bmatrix}
    0\\0\\1\\0
    \end{bmatrix},
    \quad
    \text{xe hơi} =
    \begin{bmatrix}
    0\\0\\0\\1
    \end{bmatrix}
\]

\begin{ghinho}
One-hot vector là cách biểu diễn nhãn rất phổ biến trong bài toán phân loại nhiều lớp.
\end{ghinho}

% ==========================================================
\section{One-hot encoding}

\subsection{One-hot vector là gì?}

One-hot vector là vector trong đó chỉ có một phần tử mang giá trị 1, các phần tử còn lại mang giá trị 0.

Nếu có $K$ class, one-hot vector sẽ có $K$ phần tử.

Ví dụ với $K=4$:

\[
    y =
    \begin{bmatrix}
    0\\
    1\\
    0\\
    0
    \end{bmatrix}
\]

Vector này biểu diễn class thứ 2.

\subsection{Bảng mã hóa one-hot}

\begin{center}
\begin{tabular}{lll}
\toprule
Class & Vector one-hot & Ý nghĩa \\
\midrule
Chó & $[1,0,0,0]^T$ & Phần tử thứ nhất bằng 1 \\
Mèo & $[0,1,0,0]^T$ & Phần tử thứ hai bằng 1 \\
Diều & $[0,0,1,0]^T$ & Phần tử thứ ba bằng 1 \\
Xe hơi & $[0,0,0,1]^T$ & Phần tử thứ tư bằng 1 \\
\bottomrule
\end{tabular}
\end{center}

\subsection{Output dự đoán của mạng}

Trong thực tế, output của mạng thường không ra đúng tuyệt đối các giá trị 0 và 1. Nó có thể ra một vector điểm số hoặc xác suất.

Ví dụ, với ảnh con chó:

\[
    y =
    \begin{bmatrix}
    1\\0\\0\\0
    \end{bmatrix}
\]

Nhưng mạng dự đoán:

\[
    \hat{y} =
    \begin{bmatrix}
    0.70\\
    0.20\\
    0.05\\
    0.05
    \end{bmatrix}
\]

Ta có thể hiểu mạng đang cho rằng:

\begin{itemize}
    \item xác suất là chó: $70\%$;
    \item xác suất là mèo: $20\%$;
    \item xác suất là diều: $5\%$;
    \item xác suất là xe hơi: $5\%$.
\end{itemize}

Dự đoán cuối cùng là class có giá trị lớn nhất:

\[
    \arg\max(\hat{y}) = \text{chó}
\]

\begin{vidu}
Nếu mạng cho output:

\[
    \hat{y} =
    \begin{bmatrix}
    0.1\\
    0.6\\
    0.2\\
    0.1
    \end{bmatrix}
\]

thì phần tử lớn nhất là $0.6$, nằm ở vị trí thứ 2. Mạng dự đoán ảnh thuộc class thứ 2, tức là ``mèo''.
\end{vidu}

% ==========================================================
\section{Đo sai số giữa hai vector}

\subsection{Vì sao cần đo sai số giữa hai vector?}

Trong bài toán phân loại nhiều lớp, label đúng thường là một vector one-hot, còn output dự đoán cũng là một vector.

Ví dụ:

\[
    y =
    \begin{bmatrix}
    1\\0\\0\\0
    \end{bmatrix},
    \qquad
    \hat{y} =
    \begin{bmatrix}
    0.7\\
    0.1\\
    0.1\\
    0.1
    \end{bmatrix}
\]

Ta cần một cách đo sự khác biệt giữa hai vector này.

\begin{ghinho}
Khi output và label đều là vector, sai số không còn đơn giản là phép trừ giữa hai số. Ta cần dùng các đại lượng như chuẩn vector, góc giữa hai vector hoặc cross entropy.
\end{ghinho}

% ==========================================================
\section{Hàm mất mát dựa trên chuẩn L2 - NORM 2}

\subsection{Chuẩn L2 của vector}

Với vector:

\[
    v =
    \begin{bmatrix}
    v_1\\
    v_2\\
    \vdots\\
    v_n
    \end{bmatrix}
\]

Chuẩn $L_2$ của $v$ là:

\[
    \|v\|_2 = \sqrt{v_1^2+v_2^2+\cdots+v_n^2}
\]

Đây chính là độ dài Euclid của vector.

\subsection{Sai số L2 giữa hai vector}

Với label đúng $y$ và output dự đoán $\hat{y}$, vector sai số là:

\[
    e = y-\hat{y}
\]

Sai số theo chuẩn $L_2$ là:

\[
    \|y-\hat{y}\|_2
\]

Hoặc thường dùng bình phương chuẩn $L_2$:

\[
    \|y-\hat{y}\|_2^2
\]

\subsection{Ví dụ tính sai số L2}

Giả sử:

\[
    y =
    \begin{bmatrix}
    1\\0\\0\\0
    \end{bmatrix},
    \qquad
    \hat{y} =
    \begin{bmatrix}
    0.7\\0.1\\0.1\\0.1
    \end{bmatrix}
\]

Khi đó:

\[
    e = y-\hat{y}
    =
    \begin{bmatrix}
    1-0.7\\
    0-0.1\\
    0-0.1\\
    0-0.1
    \end{bmatrix}
    =
    \begin{bmatrix}
    0.3\\
    -0.1\\
    -0.1\\
    -0.1
    \end{bmatrix}
\]

Chuẩn $L_2$:

\[
    \|e\|_2
    =
    \sqrt{0.3^2+(-0.1)^2+(-0.1)^2+(-0.1)^2}
\]

\[
    \|e\|_2
    =
    \sqrt{0.09+0.01+0.01+0.01}
    =
    \sqrt{0.12}
    \approx 0.346
\]

Bình phương chuẩn $L_2$:

\[
    \|e\|_2^2 = 0.12
\]

\begin{ghinho}
Sai số $L_2$ càng nhỏ thì hai vector càng gần nhau.
\end{ghinho}

% ==========================================================
\section{Đo khác biệt bằng góc giữa hai vector}

\subsection{Tích vô hướng}

Với hai vector:

\[
    p =
    \begin{bmatrix}
    p_1\\p_2\\ \vdots\\p_n
    \end{bmatrix},
    \qquad
    q =
    \begin{bmatrix}
    q_1\\q_2\\ \vdots\\q_n
    \end{bmatrix}
\]

Tích vô hướng là:

\[
    p \cdot q = p_1q_1+p_2q_2+\cdots+p_nq_n
\]

Tích vô hướng cho ra một số vô hướng, không phải một vector.

\begin{luuy}
Không nhầm tích vô hướng với tích có hướng. Tích vô hướng giữa hai vector cho ra một số thực. Tích có hướng mới cho ra một vector trong không gian ba chiều.
\end{luuy}

\subsection{Góc giữa hai vector}

Góc $\theta$ giữa hai vector $p$ và $q$ được tính bởi:

\[
    \cos\theta =
    \frac{p \cdot q}{\|p\|_2 \|q\|_2}
\]

Suy ra:

\[
    \theta = \arccos\left(\frac{p \cdot q}{\|p\|_2 \|q\|_2}\right)
\]

\subsection{Ý nghĩa}

\begin{itemize}
    \item Nếu $\theta$ gần $0^\circ$, hai vector gần cùng hướng, tức là tương đồng cao.
    \item Nếu $\theta$ gần $90^\circ$, hai vector gần vuông góc, tức là ít liên quan theo hướng.
    \item Nếu $\theta$ gần $180^\circ$, hai vector gần ngược hướng.
\end{itemize}

\begin{vidu}
Cho:

\[
    p =
    \begin{bmatrix}
    1\\0
    \end{bmatrix},
    \qquad
    q =
    \begin{bmatrix}
    0\\1
    \end{bmatrix}
\]

Ta có:

\[
    p \cdot q = 1\cdot 0 + 0\cdot 1 = 0
\]

\[
    \|p\|_2=1,\qquad \|q\|_2=1
\]

Do đó:

\[
    \cos\theta = \frac{0}{1\cdot 1}=0
\]

\[
    \theta = 90^\circ
\]

Hai vector vuông góc với nhau.
\end{vidu}

\subsection{Cosine similarity}

Trong nhiều ứng dụng, thay vì dùng trực tiếp góc, ta dùng cosine similarity:

\[
    \text{cosine similarity}(p,q)
    =
    \frac{p \cdot q}{\|p\|_2\|q\|_2}
\]

Giá trị này càng gần 1 thì hai vector càng giống nhau về hướng.

Một dạng loss có thể xây dựng từ cosine similarity là:

\[
    \mathcal{L}_{cos} = 1 -
    \frac{p \cdot q}{\|p\|_2\|q\|_2}
\]

% ==========================================================
\section{Cross entropy}

\subsection{Ý tưởng của cross entropy}

Cross entropy là một hàm mất mát rất phổ biến trong bài toán phân loại.

Giả sử:

\[
    p = \text{phân phối đúng}
\]

\[
    q = \text{phân phối dự đoán}
\]

Cross entropy giữa $p$ và $q$ được định nghĩa:

\[
    H(p,q) = -\sum_{i=1}^{K} p_i \log(q_i)
\]

Trong đó:

\begin{itemize}
    \item $K$ là số class;
    \item $p_i$ là xác suất đúng của class $i$;
    \item $q_i$ là xác suất mô hình dự đoán cho class $i$.
\end{itemize}

\subsection{Cross entropy với one-hot label}

Nếu label đúng là one-hot vector, chỉ có một phần tử $p_c=1$, các phần tử còn lại bằng 0.

Khi đó:

\[
    H(p,q) = -\log(q_c)
\]

Trong đó $q_c$ là xác suất mô hình gán cho class đúng.

\begin{ghinho}
Với one-hot label, cross entropy chỉ quan tâm trực tiếp đến xác suất mà mô hình gán cho class đúng. Nếu mô hình gán xác suất cao cho class đúng, loss nhỏ. Nếu mô hình gán xác suất thấp cho class đúng, loss lớn.
\end{ghinho}

\subsection{Ví dụ tính cross entropy}

Giả sử có 4 class:

\[
    \{\text{chó},\text{mèo},\text{diều},\text{xe hơi}\}
\]

Ảnh thật là chó, nên:

\[
    p =
    \begin{bmatrix}
    1\\0\\0\\0
    \end{bmatrix}
\]

Mạng dự đoán:

\[
    q =
    \begin{bmatrix}
    0.7\\
    0.1\\
    0.1\\
    0.1
    \end{bmatrix}
\]

Cross entropy:

\[
    H(p,q) =
    -[1\log(0.7)+0\log(0.1)+0\log(0.1)+0\log(0.1)]
\]

\[
    H(p,q) = -\log(0.7) \approx 0.357
\]

Nếu mạng dự đoán tệ hơn:

\[
    q =
    \begin{bmatrix}
    0.1\\
    0.6\\
    0.2\\
    0.1
    \end{bmatrix}
\]

thì:

\[
    H(p,q)=-\log(0.1)\approx 2.303
\]

Loss lớn hơn nhiều vì mô hình chỉ gán xác suất $10\%$ cho class đúng.

\subsection{So sánh các cách đo sai số}

\begin{center}
\begin{tabular}{p{0.25\textwidth}p{0.33\textwidth}p{0.32\textwidth}}
\toprule
Cách đo & Ý tưởng & Thường dùng khi \\
\midrule
Chuẩn $L_2$ & Đo độ dài vector sai số $y-\hat{y}$ & Bài toán hồi quy hoặc so sánh vector đơn giản \\
Góc giữa hai vector & Đo sự giống nhau về hướng & Bài toán so sánh độ tương đồng vector \\
Cross entropy & Đo sai khác giữa phân phối đúng và phân phối dự đoán & Bài toán phân loại nhiều lớp \\
\bottomrule
\end{tabular}
\end{center}

% ==========================================================
\section{Mô hình học được những gì?}

\subsection{Mô hình chỉ học trong phạm vi dữ liệu được dạy}

Một điểm quan trọng trong transcript là: mô hình học những gì nó được dạy.

Nếu tập dữ liệu chỉ gồm 10 loại con vật, mô hình chỉ được học để phân biệt 10 loại đó. Nếu đưa vào một con vật thứ 11 chưa từng xuất hiện trong dữ liệu huấn luyện, mô hình có thể trả lời sai hoặc trả lời một cách không đáng tin cậy.

\begin{vidu}
Giả sử mô hình được huấn luyện với 4 class:

\[
    \{\text{chó},\text{mèo},\text{diều},\text{xe hơi}\}
\]

Nếu đưa ảnh một con voi vào, nhưng mô hình không có class ``voi'', nó vẫn có thể buộc phải chọn một trong 4 class đã biết. Ví dụ, nó có thể dự đoán nhầm là ``chó'' hoặc ``xe hơi''.
\end{vidu}

\subsection{Số lượng class trong các bộ dữ liệu lớn}

Trong thực tế, nhiều mô hình phân loại ảnh được huấn luyện trên các tập dữ liệu có rất nhiều class. Ví dụ, một bộ dữ liệu có thể có 1000 class, bao gồm nhiều loại động vật, phương tiện, đồ vật và cảnh vật khác nhau.

Khi số class tăng, output layer cũng thường tăng số nơ-ron tương ứng nếu dùng one-hot encoding.

\[
    K \text{ class} \Rightarrow K \text{ output}
\]

\subsection{Không phải cứ huấn luyện là thành công}

Huấn luyện mạng nơ-ron không đảm bảo chắc chắn mô hình sẽ tốt. Kết quả phụ thuộc vào nhiều yếu tố:

\begin{itemize}
    \item dữ liệu có đủ lớn không;
    \item dữ liệu có đúng và sạch không;
    \item nhãn có chính xác không;
    \item kiến trúc mạng có phù hợp không;
    \item hàm mất mát có phù hợp không;
    \item tốc độ học có phù hợp không;
    \item số epoch có đủ không;
    \item có bị overfitting hay underfitting không.
\end{itemize}

\begin{canhbao}
Nếu chỉ hiểu ý tưởng tổng quát của backpropagation mà không hiểu các chi tiết toán học và kỹ thuật, việc tự lập trình hoặc huấn luyện mô hình phức tạp rất dễ thất bại.
\end{canhbao}

% ==========================================================
\section{Tổng hợp quy trình huấn luyện ANN}

\subsection{Quy trình tổng quát}

Một quy trình huấn luyện ANN có thể được mô tả như sau:

\begin{enumerate}
    \item Chuẩn bị tập dữ liệu gồm các cặp $(x_i,y_i)$.
    \item Thiết kế kiến trúc mạng.
    \item Chọn hàm kích hoạt.
    \item Chọn cách biểu diễn output, ví dụ one-hot vector.
    \item Chọn hàm mất mát.
    \item Khởi tạo ngẫu nhiên trọng số và bias.
    \item Thực hiện lan truyền tiến để tính output dự đoán.
    \item Tính loss giữa output dự đoán và label đúng.
    \item Dùng backpropagation để tính hướng điều chỉnh tham số.
    \item Cập nhật tham số.
    \item Lặp lại quá trình trên nhiều lần cho đến khi mô hình đạt chất lượng mong muốn.
\end{enumerate}

\subsection{Sơ đồ tổng quát}

\begin{center}
\begin{tikzpicture}[node distance=1.1cm,>=Stealth,
    box/.style={rectangle,draw,rounded corners,align=center,
    minimum width=4.2cm,minimum height=0.85cm,font=\small}
]

\node[box] (data) {Dữ liệu huấn luyện\\$(x_i,y_i)$};
\node[box,below=of data] (init) {Khởi tạo trọng số\\và bias ngẫu nhiên};
\node[box,below=of init] (forward) {Lan truyền tiến\\tính $\hat{y}_i$};
\node[box,below=of forward] (loss) {Tính hàm mất mát\\$\mathcal{L}(\hat{y}_i,y_i)$};
\node[box,below=of loss] (back) {Lan truyền ngược\\tính hướng điều chỉnh};
\node[box,below=of back] (update) {Cập nhật tham số};
\node[box,below=of update] (repeat) {Lặp lại với nhiều mẫu\\và nhiều epoch};

\draw[->] (data) -- (init);
\draw[->] (init) -- (forward);
\draw[->] (forward) -- (loss);
\draw[->] (loss) -- (back);
\draw[->] (back) -- (update);
\draw[->] (update) -- (repeat);
\draw[->] (repeat.east) .. controls +(2,0) and +(2,0) .. (forward.east);

\end{tikzpicture}
\end{center}

% ==========================================================
\section{Technical Glossary -- Bảng thuật ngữ chuyên ngành}

\small
\begin{longtable}{p{0.25\textwidth}p{0.48\textwidth}p{0.20\textwidth}}
\toprule
\textbf{Thuật ngữ} & \textbf{Giải thích} & \textbf{Ví dụ} \\
\midrule
\endfirsthead

\toprule
\textbf{Thuật ngữ} & \textbf{Giải thích} & \textbf{Ví dụ} \\
\midrule
\endhead

ANN & Artificial Neural Network, mạng nơ-ron nhân tạo. & Mạng phân loại ảnh. \\

Training & Quá trình huấn luyện mô hình bằng dữ liệu. & Cho mạng học từ ảnh và nhãn. \\

Supervised learning & Học có giám sát, mô hình học từ các cặp input--label. & Ảnh chó đi kèm nhãn ``chó''. \\

Input & Dữ liệu đầu vào của mạng. & Một ảnh, một vector giá cổ phiếu. \\

Label & Nhãn đúng của dữ liệu. & Chó, mèo, xe hơi. \\

Target & Đầu ra mong muốn, thường tương đương label. & Vector $[1,0,0,0]^T$. \\

Output & Kết quả mà mạng tạo ra. & $\hat{y}=[0.7,0.1,0.1,0.1]^T$. \\

Parameter & Tham số được học từ dữ liệu. & Trọng số $w$ và bias $b$. \\

Weight & Trọng số, biểu diễn mức ảnh hưởng của một kết nối. & $w_1=0.5$. \\

Bias & Hằng số cộng thêm vào tổng có trọng số. & $z=wx+b$. \\

Random initialization & Khởi tạo tham số ban đầu bằng giá trị ngẫu nhiên. & Trọng số ban đầu lấy ngẫu nhiên. \\

Forward propagation & Lan truyền tiến từ input đến output. & Tính lần lượt qua các lớp. \\

Backpropagation & Lan truyền ngược sai số để cập nhật tham số. & Dùng loss để chỉnh trọng số. \\

Loss function & Hàm mất mát, đo sai khác giữa dự đoán và label. & Cross entropy. \\

Error & Sai số giữa output dự đoán và output mong muốn. & $e=y-\hat{y}$. \\

Gradient & Vector đạo hàm cho biết hướng tăng nhanh nhất của loss. & $\nabla_\theta J(\theta)$. \\

Learning rate & Tốc độ học, quyết định bước cập nhật tham số lớn hay nhỏ. & $\eta=0.001$. \\

Epoch & Một lượt đi qua toàn bộ tập dữ liệu huấn luyện. & Huấn luyện 50 epoch. \\

Class & Lớp dữ liệu trong bài toán phân loại. & Chó, mèo, xe hơi. \\

Classification & Bài toán phân loại dữ liệu vào các class. & Phân loại ảnh động vật. \\

One-hot vector & Vector có đúng một phần tử bằng 1, còn lại bằng 0. & $[0,1,0,0]^T$. \\

One-hot encoding & Cách mã hóa nhãn bằng one-hot vector. & Mèo $\rightarrow [0,1,0,0]^T$. \\

$L_2$ norm & Chuẩn Euclid của vector. & $\|v\|_2=\sqrt{\sum v_i^2}$. \\

Dot product & Tích vô hướng giữa hai vector, cho ra một số. & $p\cdot q=\sum p_iq_i$. \\

Cosine similarity & Độ tương đồng theo góc giữa hai vector. & $\frac{p\cdot q}{\|p\|\|q\|}$. \\

Cross entropy & Hàm mất mát phổ biến cho phân loại. & $-\sum p_i\log q_i$. \\

GPU & Bộ xử lý đồ họa, thường dùng để tăng tốc huấn luyện mạng lớn. & Card màn hình NVIDIA. \\

\bottomrule
\end{longtable}
\normalsize

% ==========================================================
\section{Key Concepts Summary -- Tóm tắt kiến thức trọng tâm}

\subsection{Các ý chính cần nhớ}

\begin{enumerate}
    \item Ban đầu, mạng nơ-ron chưa biết suy luận vì các tham số được khởi tạo ngẫu nhiên.
    \item Khi đưa input vào mạng, mạng thực hiện lan truyền tiến để tạo output dự đoán.
    \item Output dự đoán được so sánh với label đúng để tính sai số.
    \item Hàm mất mát biến sai số thành một số vô hướng.
    \item Backpropagation dùng sai số này để quay ngược lại điều chỉnh các trọng số và bias.
    \item Quá trình cập nhật tham số được lặp lại rất nhiều lần.
    \item Trong phân loại nhiều lớp, label thường được biểu diễn bằng one-hot vector.
    \item Các cách đo sai số giữa vector gồm chuẩn $L_2$, góc giữa hai vector và cross entropy.
    \item Cross entropy đặc biệt phổ biến trong bài toán phân loại.
    \item Mô hình chỉ học tốt những gì được thể hiện đủ rõ và đủ đúng trong dữ liệu huấn luyện.
\end{enumerate}

\subsection{Công thức quan trọng}

\subsection*{Lan truyền tiến qua một lớp}

\[
    z^{(l)} = W^{(l)}a^{(l-1)} + b^{(l)}
\]

\[
    a^{(l)} = \varphi(z^{(l)})
\]

\subsection*{Sai số vector}

\[
    e = y-\hat{y}
\]

\subsection*{Chuẩn L2}

\[
    \|e\|_2 = \sqrt{e_1^2+e_2^2+\cdots+e_n^2}
\]

\subsection*{Góc giữa hai vector}

\[
    \cos\theta =
    \frac{p\cdot q}{\|p\|_2\|q\|_2}
\]

\subsection*{Cross entropy}

\[
    H(p,q)=-\sum_{i=1}^{K}p_i\log(q_i)
\]

\subsection*{Cập nhật tham số theo gradient descent}

\[
    \theta_{t+1}=\theta_t-\eta\nabla_{\theta}J(\theta_t)
\]

% ==========================================================
\section{Review Questions -- Câu hỏi ôn tập}

\begin{enumerate}
    \item Vì sao mạng nơ-ron ban đầu thường cho output sai?
    \item Khởi tạo ngẫu nhiên trọng số có ý nghĩa gì?
    \item Lan truyền tiến là gì?
    \item Backpropagation là gì?
    \item Vì sao cần hàm mất mát?
    \item Hàm mất mát phải trả về vector hay số vô hướng?
    \item One-hot vector là gì?
    \item Với 5 class, one-hot vector có bao nhiêu phần tử?
    \item Viết công thức chuẩn $L_2$ của một vector.
    \item Tích vô hướng giữa hai vector cho ra số hay vector?
    \item Góc giữa hai vector cho biết điều gì?
    \item Cross entropy thường dùng cho loại bài toán nào?
    \item Vì sao mô hình có thể trả lời sai khi gặp class chưa từng được học?
    \item Vì sao huấn luyện mạng nơ-ron có thể mất nhiều giờ hoặc nhiều ngày?
    \item Vì sao chỉ hiểu ý tưởng backpropagation là chưa đủ để lập trình hoàn chỉnh?
\end{enumerate}

% ==========================================================
\section{Practice Exercises -- Bài tập vận dụng}

\subsection*{Bài tập 1: One-hot encoding}
\addcontentsline{toc}{section}{Bài tập 1: One-hot encoding}

Cho 5 class:

\[
    \{\text{chó}, \text{mèo}, \text{chim}, \text{xe hơi}, \text{máy bay}\}
\]

Hãy viết one-hot vector cho từng class theo đúng thứ tự trên.

\subsection*{Lời giải gợi ý}

\[
    \text{chó} =
    \begin{bmatrix}
    1\\0\\0\\0\\0
    \end{bmatrix},
    \quad
    \text{mèo} =
    \begin{bmatrix}
    0\\1\\0\\0\\0
    \end{bmatrix}
\]

\[
    \text{chim} =
    \begin{bmatrix}
    0\\0\\1\\0\\0
    \end{bmatrix},
    \quad
    \text{xe hơi} =
    \begin{bmatrix}
    0\\0\\0\\1\\0
    \end{bmatrix},
    \quad
    \text{máy bay} =
    \begin{bmatrix}
    0\\0\\0\\0\\1
    \end{bmatrix}
\]

\subsection*{Bài tập 2: Tính chuẩn L2}
\addcontentsline{toc}{section}{Bài tập 2: Tính chuẩn L2}

Cho:

\[
    y =
    \begin{bmatrix}
    1\\0\\0
    \end{bmatrix},
    \qquad
    \hat{y} =
    \begin{bmatrix}
    0.8\\0.1\\0.1
    \end{bmatrix}
\]

Tính:

\[
    \|y-\hat{y}\|_2
\]

\subsection*{Lời giải gợi ý}

\[
    e=y-\hat{y}
    =
    \begin{bmatrix}
    0.2\\-0.1\\-0.1
    \end{bmatrix}
\]

\[
    \|e\|_2
    =
    \sqrt{0.2^2+(-0.1)^2+(-0.1)^2}
\]

\[
    \|e\|_2
    =
    \sqrt{0.04+0.01+0.01}
    =
    \sqrt{0.06}
    \approx 0.245
\]

\subsection*{Bài tập 3: Tính cross entropy}
\addcontentsline{toc}{section}{Bài tập 3: Tính cross entropy}

Cho label đúng:

\[
    p =
    \begin{bmatrix}
    0\\1\\0
    \end{bmatrix}
\]

Mạng dự đoán:

\[
    q =
    \begin{bmatrix}
    0.2\\0.7\\0.1
    \end{bmatrix}
\]

Tính cross entropy:

\[
    H(p,q)=-\sum_{i=1}^{3}p_i\log(q_i)
\]

\subsection*{Lời giải gợi ý}

Vì class đúng là class thứ 2:

\[
    H(p,q)=-\log(0.7)\approx 0.357
\]

\subsection*{Bài tập 4: Xác định output layer}
\addcontentsline{toc}{section}{Bài tập 4: Xác định output layer}

Một bài toán phân loại ảnh có 12 class. Nếu dùng one-hot encoding, lớp output nên có bao nhiêu nơ-ron?

\subsection*{Lời giải gợi ý}

Vì có 12 class, one-hot vector có 12 phần tử. Do đó, lớp output nên có 12 nơ-ron.

\[
    12\ \text{class} \Rightarrow 12\ \text{output neurons}
\]

% ==========================================================
\section*{Kết luận}
\addcontentsline{toc}{section}{Kết luận}

ANN3 giải thích phần cốt lõi của quá trình huấn luyện mạng nơ-ron. Ban đầu, mạng có các trọng số và bias ngẫu nhiên nên output thường sai. Thông qua lan truyền tiến, mạng tạo dự đoán. Dự đoán này được so sánh với label đúng bằng một hàm mất mát. Sau đó, backpropagation dùng sai số để điều chỉnh các tham số của mạng.

Trong bài toán phân loại nhiều lớp, nhãn thường được mã hóa bằng one-hot vector. Khi đó, output và label đều là vector, nên cần các cách đo sai số giữa vector như chuẩn $L_2$, góc giữa hai vector hoặc cross entropy. Trong thực tế, cross entropy là lựa chọn rất phổ biến cho bài toán phân loại.

Điều quan trọng cần nhớ là huấn luyện mạng nơ-ron không chỉ là chạy một thuật toán có sẵn. Muốn mô hình hoạt động tốt, người học cần hiểu dữ liệu, cách mã hóa nhãn, hàm mất mát, thuật toán tối ưu, kiến trúc mạng và nhiều chi tiết kỹ thuật khác.



\chapter{ANN4: Gradient Descent -- Hạ gradient}


% ==========================================================
\section*{Learning Objectives -- Mục tiêu bài học}
\addcontentsline{toc}{section}{Mục tiêu bài học}

Bài ANN4 tiếp tục phần huấn luyện mạng nơ-ron, tập trung vào bản chất tối ưu hóa của quá trình học. Sau khi học xong, người học cần nắm được:

\begin{enumerate}[label=\arabic*.]
    \item Nhắc lại vai trò của hàm mất mát trong huấn luyện mạng nơ-ron.
    \item Hiểu vì sao huấn luyện ANN là một bài toán tối ưu.
    \item Nắm được ý tưởng của phương pháp Gradient Descent.
    \item Biết cách dùng đạo hàm để quyết định hướng cập nhật tham số.
    \item Hiểu vai trò của learning rate trong quá trình hội tụ.
    \item Nhận biết sự khác nhau giữa bài toán một biến và bài toán nhiều biến.
    \item Hiểu khái niệm cực trị địa phương và cực trị toàn cục.
    \item Nắm được trực giác về exploration, exploitation, momentum và thuật toán di truyền.
    \item Biết khi nào nên dùng các hàm mất mát phổ biến như cross entropy hoặc mean squared error.
\end{enumerate}

% ==========================================================
\section{Neural Network Training Recap -- Nhắc lại quá trình huấn luyện mạng nơ-ron}

\subsection{Từ dữ liệu đến sai số}

Trong học có giám sát, ta cung cấp cho mạng nơ-ron các cặp dữ liệu:

\[
    (x_i, y_i)
\]

Trong đó:

\begin{itemize}
    \item $x_i$ là dữ liệu đầu vào, còn gọi là \textit{data} hoặc \textit{input};
    \item $y_i$ là nhãn đúng, còn gọi là \textit{label} hoặc \textit{target}.
\end{itemize}

Mạng nơ-ron nhận $x_i$ và tạo ra dự đoán:

\[
    \hat{y}_i = f(x_i;\theta)
\]

Trong đó $\theta$ là toàn bộ tham số của mạng, gồm các trọng số và bias.

Sai số xuất hiện khi dự đoán $\hat{y}_i$ khác với nhãn đúng $y_i$.

\[
    \text{Sai số} = \text{khác biệt giữa } y_i \text{ và } \hat{y}_i
\]

\begin{ghinho}
Huấn luyện mạng nơ-ron là quá trình dùng sai số giữa output mong muốn và output hiện tại để điều chỉnh các tham số của mạng.
\end{ghinho}

\subsection{Hàm mất mát}

Hàm mất mát, tiếng Anh là \textit{loss function}, là hàm biến sự khác biệt giữa $y$ và $\hat{y}$ thành một số vô hướng.

\[
    \mathcal{L}(y,\hat{y}) \in \mathbb{R}
\]

Giá trị loss càng nhỏ thì mô hình dự đoán càng tốt.

Một số dạng hàm mất mát phổ biến:

\begin{itemize}
    \item chuẩn $L_2$ hoặc bình phương sai số;
    \item góc giữa hai vector;
    \item cross entropy.
\end{itemize}

\subsection{Mục tiêu lý tưởng}

Trong trường hợp lý tưởng, nếu mô hình dự đoán hoàn hảo, sai số sẽ tiến về 0.

\[
    \mathcal{L}(y,\hat{y}) = 0
\]

Tuy nhiên, trong thực tế, dữ liệu có thể nhiễu, mô hình có thể chưa đủ tốt, quá trình tối ưu có thể chưa đạt cực trị tốt nhất. Do đó, ta thường chỉ kỳ vọng loss giảm dần đến một mức đủ nhỏ.

\begin{vidu}
Giả sử bài toán phân loại ảnh chó và mèo.

Nếu ảnh thật là chó:

\[
    y = \text{chó}
\]

Mạng dự đoán đúng:

\[
    \hat{y} = \text{chó}
\]

thì sai số rất nhỏ.

Nếu mạng dự đoán:

\[
    \hat{y} = \text{mèo}
\]

thì sai số lớn hơn. Hàm mất mát giúp lượng hóa mức sai này thành một con số cụ thể.
\end{vidu}

% ==========================================================
\section{ANN Training as Optimization -- Huấn luyện ANN là một bài toán tối ưu}

\subsection{Bản chất của quá trình huấn luyện}

Trong mạng nơ-ron, ta cần tìm bộ tham số:

\[
    \theta = \{W,b\}
\]

sao cho hàm mất mát trên tập dữ liệu là nhỏ nhất.

Với một mẫu dữ liệu:

\[
    \mathcal{L}(y_i,\hat{y}_i)
\]

Với toàn bộ tập huấn luyện gồm $N$ mẫu, ta thường xét hàm mục tiêu:

\[
    J(\theta)
    =
    \frac{1}{N}
    \sum_{i=1}^{N}
    \mathcal{L}\left(y_i, f(x_i;\theta)\right)
\]

Mục tiêu huấn luyện là:

\[
    \theta^*
    =
    \operatorname*{arg\,min}_{\theta} J(\theta)
\]

Trong đó:

\begin{itemize}
    \item $J(\theta)$ là hàm mục tiêu cần tối thiểu hóa;
    \item $\theta$ là bộ tham số cần điều chỉnh;
    \item $\theta^*$ là bộ tham số tối ưu hoặc gần tối ưu.
\end{itemize}

\begin{ghinho}
Bản chất toán học của huấn luyện mạng nơ-ron là bài toán tối ưu hóa: tìm các trọng số và bias để hàm mất mát đạt giá trị nhỏ nhất có thể.
\end{ghinho}

\subsection{Số lượng biến rất lớn}

Trong bài toán tối ưu thông thường, ta có thể chỉ cần tìm một biến $x$. Nhưng trong mạng nơ-ron, số biến chính là số tham số của mạng.

Một mạng nơ-ron có thể có:

\begin{itemize}
    \item vài nghìn tham số;
    \item vài triệu tham số;
    \item vài chục triệu tham số;
    \item thậm chí vài trăm triệu hoặc nhiều hơn.
\end{itemize}

Vì vậy, bài toán tối ưu trong ANN thường là bài toán đa biến quy mô lớn.

\subsection{Không phải lúc nào cũng có lời giải giải tích}

Với một số hàm đơn giản, ví dụ hàm bậc hai, ta có thể tìm cực trị bằng cách giải phương trình đạo hàm bằng 0.

Ví dụ:

\[
    f(x) = x^2 - 4x + 5
\]

Đạo hàm:

\[
    f'(x) = 2x - 4
\]

Cho:

\[
    f'(x)=0
\]

Ta được:

\[
    2x-4=0 \Rightarrow x=2
\]

Đây là lời giải giải tích.

Tuy nhiên, trong mạng nơ-ron, hàm mất mát thường rất phức tạp vì phụ thuộc vào nhiều lớp, nhiều hàm kích hoạt, nhiều tham số và dữ liệu lớn. Do đó, ta thường không có lời giải giải tích trực tiếp, mà phải dùng phương pháp tính lặp.

\begin{ghinho}
Trong ANN, ta thường không giải trực tiếp được $\theta^*$ bằng công thức đóng. Thay vào đó, ta bắt đầu từ một bộ tham số ban đầu và cập nhật lặp để giảm loss.
\end{ghinho}

% ==========================================================
\section{Gradient Descent: Single Variable -- Gradient Descent trong trường hợp một biến}

\subsection{Bài toán minh họa}

Giả sử ta cần tìm giá trị $x$ sao cho hàm:

\[
    f(x)
\]

đạt cực tiểu.

Ta có thể hình dung đồ thị của $f(x)$ như một thung lũng. Mục tiêu là tìm điểm thấp nhất của thung lũng.

\begin{center}
\begin{tikzpicture}[scale=1.0,>=Stealth]
    % Axes
    \draw[->] (-4,0) -- (4,0) node[right] {$x$};
    \draw[->] (0,-0.2) -- (0,4) node[above] {$f(x)$};

    % Function curve
    \draw[domain=-3.2:3.2,smooth,variable=\x,thick,blue]
        plot ({\x},{0.25*(\x)^2 + 0.4});

    % Minimum
    \filldraw[red] (0,0.4) circle (2pt);
    \node[below right] at (0,0.4) {cực tiểu};

    % x1 left
    \filldraw[black] (-2,1.4) circle (2pt);
    \node[above left] at (-2,1.4) {$x_1$};

    % x2 right
    \filldraw[black] (2,1.4) circle (2pt);
    \node[above right] at (2,1.4) {$x_2$};

    % Arrows
    \draw[->,thick,red] (-2,1.2) -- (-1.2,0.8);
    \draw[->,thick,red] (2,1.2) -- (1.2,0.8);

\end{tikzpicture}
\end{center}

Nếu ta bắt đầu ở bên trái điểm cực tiểu, ta cần tăng $x$ để đi về phía cực tiểu.

Nếu ta bắt đầu ở bên phải điểm cực tiểu, ta cần giảm $x$ để đi về phía cực tiểu.

Câu hỏi quan trọng là: làm sao biết nên tăng hay giảm $x$?

Câu trả lời là: dùng đạo hàm.

\subsection{Đạo hàm cho biết hướng đi}

Đạo hàm $f'(x)$ cho biết độ dốc của hàm tại vị trí hiện tại.

\begin{itemize}
    \item Nếu $f'(x) < 0$, hàm đang giảm khi $x$ tăng. Muốn giảm $f(x)$, ta nên tăng $x$.
    \item Nếu $f'(x) > 0$, hàm đang tăng khi $x$ tăng. Muốn giảm $f(x)$, ta nên giảm $x$.
\end{itemize}

Công thức cập nhật của Gradient Descent trong trường hợp một biến là:

\[
    x_{\text{new}}
    =
    x_{\text{old}}
    -
    \eta f'(x_{\text{old}})
\]

Trong đó:

\begin{itemize}
    \item $x_{\text{old}}$ là giá trị hiện tại;
    \item $x_{\text{new}}$ là giá trị sau khi cập nhật;
    \item $f'(x_{\text{old}})$ là đạo hàm tại điểm hiện tại;
    \item $\eta$ là learning rate.
\end{itemize}

\begin{ghinho}
Dấu trừ trong công thức Gradient Descent thể hiện việc đi ngược hướng tăng của hàm để làm hàm giảm xuống.
\end{ghinho}

\subsection{Ví dụ trực quan}

Giả sử:

\[
    f(x) = (x-3)^2
\]

Hàm này đạt cực tiểu tại:

\[
    x=3
\]

Đạo hàm:

\[
    f'(x)=2(x-3)
\]

Chọn learning rate:

\[
    \eta = 0.1
\]

Giả sử bắt đầu từ:

\[
    x_0 = 0
\]

Ta cập nhật:

\[
    x_{1}=x_0-\eta f'(x_0)
\]

\[
    f'(0)=2(0-3)=-6
\]

\[
    x_1=0-0.1(-6)=0.6
\]

Tiếp tục:

\[
    f'(0.6)=2(0.6-3)=-4.8
\]

\[
    x_2=0.6-0.1(-4.8)=1.08
\]

Tiếp tục:

\[
    f'(1.08)=2(1.08-3)=-3.84
\]

\[
    x_3=1.08-0.1(-3.84)=1.464
\]

Ta thấy $x$ đang tăng dần và tiến về 3.

\begin{vidu}
Bảng vài bước cập nhật:

\begin{center}
\begin{tabular}{cccc}
\toprule
Bước & $x$ & $f'(x)$ & $f(x)$ \\
\midrule
0 & 0.000 & -6.000 & 9.000 \\
1 & 0.600 & -4.800 & 5.760 \\
2 & 1.080 & -3.840 & 3.686 \\
3 & 1.464 & -3.072 & 2.360 \\
4 & 1.771 & -2.458 & 1.511 \\
\bottomrule
\end{tabular}
\end{center}

Giá trị $f(x)$ giảm dần, tức thuật toán đang tiến về cực tiểu.
\end{vidu}

% ==========================================================
\section{Learning Rate -- Learning rate}

\subsection{Learning rate là gì?}

Learning rate, thường ký hiệu là $\eta$, là hệ số quyết định mỗi lần cập nhật tham số sẽ đi một bước lớn hay nhỏ.

Trong công thức:

\[
    x_{\text{new}}
    =
    x_{\text{old}}
    -
    \eta f'(x_{\text{old}})
\]

đại lượng:

\[
    \eta f'(x_{\text{old}})
\]

quyết định độ lớn của bước cập nhật.

\subsection{Learning rate quá nhỏ}

Nếu $\eta$ quá nhỏ, mỗi lần cập nhật chỉ thay đổi rất ít. Thuật toán vẫn có thể đi đúng hướng nhưng sẽ hội tụ rất chậm.

\begin{vidu}
Nếu điểm hiện tại còn rất xa cực tiểu nhưng learning rate quá nhỏ, thuật toán giống như đi từng bước rất ngắn. Kết quả là cần rất nhiều vòng lặp mới đến gần nghiệm.
\end{vidu}

\subsection{Learning rate quá lớn}

Nếu $\eta$ quá lớn, bước cập nhật quá dài. Thuật toán có thể nhảy qua điểm cực tiểu, dao động qua lại hoặc thậm chí không hội tụ.

\begin{center}
\begin{tikzpicture}[scale=1.0,>=Stealth]
    % Axes
    \draw[->] (-4,0) -- (4,0) node[right] {$x$};
    \draw[->] (0,-0.2) -- (0,4) node[above] {$f(x)$};

    % Function curve
    \draw[domain=-3.2:3.2,smooth,variable=\x,thick,blue]
        plot ({\x},{0.25*(\x)^2 + 0.4});

    % Minimum
    \filldraw[red] (0,0.4) circle (2pt);
    \node[below right] at (0,0.4) {cực tiểu};

    % Big jumps
    \filldraw[black] (-2.5,1.9625) circle (2pt);
    \filldraw[black] (2.0,1.4) circle (2pt);
    \filldraw[black] (-1.6,1.04) circle (2pt);
    \draw[->,thick,orange] (-2.5,1.9) -- (2.0,1.35);
    \draw[->,thick,orange] (2.0,1.35) -- (-1.6,1.0);

    \node[above] at (-2.5,1.9625) {bước lớn};
\end{tikzpicture}
\end{center}

\begin{canhbao}
Learning rate quá lớn có thể làm quá trình học không ổn định. Loss có thể không giảm mà dao động hoặc tăng lên.
\end{canhbao}

\subsection{Learning rate thay đổi theo thời gian}

Trong nhiều thuật toán, learning rate không cố định. Một chiến lược thường gặp là:

\begin{itemize}
    \item ban đầu dùng learning rate lớn để đi nhanh;
    \item về sau giảm learning rate để tinh chỉnh gần điểm cực tiểu.
\end{itemize}

Ta có thể ký hiệu learning rate tại bước $t$ là:

\[
    \eta_t
\]

và công thức cập nhật:

\[
    x_{t+1}=x_t-\eta_t f'(x_t)
\]

Trong đó $\eta_t$ có thể giảm dần theo thời gian.

\begin{ghinho}
Chọn learning rate là một trong những quyết định quan trọng nhất khi huấn luyện mạng nơ-ron.
\end{ghinho}

% ==========================================================
\section{Gradient Descent trong bài toán nhiều biến}

\subsection{Từ đạo hàm sang gradient}

Trong mạng nơ-ron, ta không chỉ có một biến $x$, mà có rất nhiều tham số:

\[
    \theta =
    \begin{bmatrix}
    \theta_1\\
    \theta_2\\
    \vdots\\
    \theta_m
    \end{bmatrix}
\]

Hàm mục tiêu:

\[
    J(\theta)
\]

là hàm nhiều biến. Khi đó, ta không dùng đạo hàm thông thường mà dùng gradient:

\[
    \nabla_{\theta}J(\theta)
    =
    \begin{bmatrix}
    \frac{\partial J}{\partial \theta_1}\\
    \frac{\partial J}{\partial \theta_2}\\
    \vdots\\
    \frac{\partial J}{\partial \theta_m}
    \end{bmatrix}
\]

Gradient là vector gồm các đạo hàm riêng của hàm mất mát theo từng tham số.

\subsection{Công thức cập nhật tổng quát}

Công thức Gradient Descent trong bài toán nhiều biến:

\[
    \theta_{t+1}
    =
    \theta_t
    -
    \eta \nabla_{\theta}J(\theta_t)
\]

Nếu xét riêng một trọng số $w$, ta có:

\[
    w_{\text{new}}
    =
    w_{\text{old}}
    -
    \eta
    \frac{\partial J}{\partial w}
\]

Nếu xét bias $b$:

\[
    b_{\text{new}}
    =
    b_{\text{old}}
    -
    \eta
    \frac{\partial J}{\partial b}
\]

\begin{ghinho}
Trong ANN, mỗi trọng số và mỗi bias đều được cập nhật dựa trên đạo hàm riêng của hàm mất mát theo chính tham số đó.
\end{ghinho}

\subsection{Liên hệ với backpropagation}

Backpropagation giúp tính các đạo hàm riêng:

\[
    \frac{\partial J}{\partial w}
    \quad \text{và} \quad
    \frac{\partial J}{\partial b}
\]

cho toàn bộ mạng một cách hiệu quả. Sau khi có các đạo hàm này, ta dùng thuật toán tối ưu như Gradient Descent để cập nhật tham số.

Tóm lại:

\[
    \text{Backpropagation} \Rightarrow \text{tính gradient}
\]

\[
    \text{Gradient Descent} \Rightarrow \text{cập nhật tham số}
\]

\subsection{Ví dụ cập nhật một trọng số}

Giả sử tại một bước huấn luyện:

\[
    w = 0.8
\]

Gradient theo trọng số này là:

\[
    \frac{\partial J}{\partial w}=0.5
\]

Chọn learning rate:

\[
    \eta=0.1
\]

Cập nhật:

\[
    w_{\text{new}}
    =
    0.8 - 0.1 \cdot 0.5
    =
    0.8 - 0.05
    =
    0.75
\]

Nếu gradient âm:

\[
    \frac{\partial J}{\partial w}=-0.5
\]

thì:

\[
    w_{\text{new}}
    =
    0.8 - 0.1(-0.5)
    =
    0.85
\]

Như vậy, dấu của gradient quyết định tham số tăng hay giảm.

% ==========================================================
\section{Cực trị địa phương và cực trị toàn cục}

\subsection{Cực trị toàn cục}

Cực trị toàn cục, tiếng Anh là \textit{global minimum}, là điểm có giá trị hàm mất mát nhỏ nhất trên toàn bộ không gian tìm kiếm.

Nếu $\theta^*$ là cực tiểu toàn cục thì:

\[
    J(\theta^*) \leq J(\theta)
\]

với mọi $\theta$ trong miền xét.

\subsection{Cực trị địa phương}

Cực trị địa phương, tiếng Anh là \textit{local minimum}, là điểm thấp nhất trong một vùng lân cận, nhưng chưa chắc là điểm thấp nhất toàn cục.

\begin{center}
\begin{tikzpicture}[scale=1.0,>=Stealth]
    \draw[->] (-4,0) -- (4.5,0) node[right] {$x$};
    \draw[->] (0,-0.5) -- (0,4) node[above] {$J(x)$};

    \draw[domain=-3.6:4,smooth,variable=\x,thick,blue]
        plot ({\x},{0.08*(\x+3)*(\x+1)*(\x-1.2)*(\x-3)+1.8});

    \filldraw[red] (-2.15,0.95) circle (2pt);
    \node[below] at (-2.15,0.95) {cực trị địa phương};

    \filldraw[red] (2.55,0.35) circle (2pt);
    \node[below] at (2.55,0.35) {cực trị toàn cục};

\end{tikzpicture}
\end{center}

\subsection{Vấn đề mắc kẹt ở cực trị địa phương}

Gradient Descent dựa vào thông tin cục bộ tại điểm hiện tại. Nếu thuật toán đi xuống một ``hố'' địa phương, nó có thể bị mắc kẹt ở đó và không tìm được cực trị toàn cục.

\begin{luuy}
Một biến cũng có thể có nhiều cực trị. Không phải chỉ bài toán nhiều biến mới có nhiều cực trị. Tuy nhiên, trong bài toán nhiều biến, bề mặt mất mát phức tạp hơn rất nhiều.
\end{luuy}

\subsection{Ảnh hưởng trong huấn luyện ANN}

Trong mạng nơ-ron, hàm mất mát theo hàng triệu tham số tạo thành một bề mặt rất phức tạp. Bề mặt này có thể có:

\begin{itemize}
    \item nhiều cực trị địa phương;
    \item vùng yên ngựa;
    \item vùng phẳng;
    \item dốc rất lớn hoặc rất nhỏ.
\end{itemize}

Do đó, việc tối ưu không hề đơn giản. Không phải cứ chạy huấn luyện là chắc chắn tìm được nghiệm tốt nhất.

% ==========================================================
\section{Exploration và Exploitation}

\subsection{Exploitation là gì?}

Exploitation có thể hiểu là khai thác thông tin hiện có ở khu vực hiện tại để cải thiện lời giải.

Trong Gradient Descent, tại một điểm hiện tại, ta dùng gradient ở điểm đó để quyết định hướng đi tiếp theo. Đây là dạng khai thác thông tin cục bộ.

\[
    \text{Thông tin cục bộ} \Rightarrow \text{đi theo hướng giảm loss}
\]

\subsection{Exploration là gì?}

Exploration là thăm dò các khu vực khác của không gian tìm kiếm. Mục tiêu là tránh việc chỉ tập trung vào một vùng hiện tại và bỏ lỡ các nghiệm tốt hơn ở vùng khác.

\begin{vidu}
Có thể hình dung như việc tìm vàng:

\begin{itemize}
    \item Exploitation: tập trung khai thác mỏ vàng hiện tại.
    \item Exploration: cử một số người đi tìm các mỏ vàng khác.
\end{itemize}

Nếu chỉ khai thác mỏ hiện tại, ta có thể bỏ lỡ một mỏ vàng lớn hơn ở nơi khác.
\end{vidu}

\subsection{Cân bằng exploration và exploitation}

Một thuật toán tối ưu tốt thường cần cân bằng:

\[
    \text{khai thác lời giải hiện tại}
    \quad \text{và} \quad
    \text{tìm kiếm vùng mới}
\]

Nếu quá thiên về exploitation, thuật toán dễ mắc kẹt tại cực trị địa phương.

Nếu quá thiên về exploration, thuật toán có thể tìm kiếm lan man và hội tụ chậm.

\begin{ghinho}
Trong tối ưu hóa, một trong những khó khăn lớn là cân bằng giữa việc khai thác nghiệm hiện tại và khám phá các vùng nghiệm mới.
\end{ghinho}

% ==========================================================
\section{Thuật toán di truyền và metaheuristic}

\subsection{Metaheuristic là gì?}

Metaheuristic là nhóm thuật toán tối ưu cấp cao, thường được thiết kế để tìm nghiệm tốt trong những bài toán khó, đặc biệt khi:

\begin{itemize}
    \item không có lời giải giải tích;
    \item không có đạo hàm;
    \item không gian tìm kiếm quá lớn;
    \item hàm mục tiêu có nhiều cực trị địa phương.
\end{itemize}

Các thuật toán metaheuristic thường không đảm bảo tìm được nghiệm tối ưu toàn cục, nhưng có thể tìm được nghiệm tốt trong thời gian chấp nhận được.

\subsection{Thuật toán di truyền}

Thuật toán di truyền, tiếng Anh là \textit{genetic algorithm}, là một ví dụ của metaheuristic. Nó lấy cảm hứng từ tiến hóa tự nhiên.

Một quần thể gồm nhiều cá thể được tạo ra. Mỗi cá thể đại diện cho một lời giải. Qua nhiều thế hệ, các lời giải tốt được chọn, lai ghép và đột biến để tạo lời giải mới.

Các bước khái quát:

\begin{enumerate}
    \item Khởi tạo một quần thể lời giải.
    \item Đánh giá độ tốt của từng lời giải bằng hàm mục tiêu.
    \item Chọn các lời giải tốt hơn.
    \item Lai ghép và đột biến để tạo thế hệ mới.
    \item Lặp lại cho đến khi đạt điều kiện dừng.
\end{enumerate}

\subsection{Liên hệ với ANN}

Về nguyên tắc, thuật toán di truyền có thể được dùng để tối ưu trọng số của mạng nơ-ron, nhất là ở quy mô nhỏ.

Tuy nhiên, với mạng lớn có hàng triệu hoặc hàng trăm triệu tham số, thuật toán di truyền thường không hiệu quả bằng các phương pháp dựa trên gradient.

\begin{luuy}
Trong huấn luyện mạng nơ-ron hiện đại, các phương pháp dựa trên gradient vẫn là nền tảng chính. Metaheuristic có thể hữu ích trong một số bài toán đặc biệt hoặc quy mô nhỏ.
\end{luuy}

% ==========================================================
\section{Momentum và ý tưởng vượt cực trị địa  - local minimum}

\subsection{Hạn chế của Gradient Descent cơ bản}

Gradient Descent cơ bản chỉ dựa vào gradient hiện tại. Vì vậy, khi gặp một cực trị địa phương hoặc vùng trũng nhỏ, thuật toán có thể bị mắc kẹt.

Ta có thể hình dung một hòn bi lăn xuống dốc. Nếu hòn bi không có quán tính, nó có thể dừng lại ở một hố nhỏ. Nhưng nếu có đủ động lượng, nó có thể vượt qua hố nhỏ và tiếp tục lăn xuống vùng thấp hơn.

\subsection{Ý tưởng momentum}

Momentum bổ sung yếu tố quán tính vào quá trình cập nhật. Thay vì chỉ dùng gradient hiện tại, thuật toán còn xét hướng di chuyển trước đó.

Một dạng công thức đơn giản:

\[
    v_t = \beta v_{t-1} + \eta \nabla_{\theta}J(\theta_t)
\]

\[
    \theta_{t+1} = \theta_t - v_t
\]

Trong đó:

\begin{itemize}
    \item $v_t$ là vận tốc cập nhật tại bước $t$;
    \item $\beta$ là hệ số momentum;
    \item $\eta$ là learning rate;
    \item $\nabla_{\theta}J(\theta_t)$ là gradient hiện tại.
\end{itemize}

\begin{ghinho}
Momentum giúp thuật toán có xu hướng duy trì hướng đi, nhờ đó có thể giảm dao động và đôi khi vượt qua các vùng cực trị địa phương nông.
\end{ghinho}

\subsection{Không đảm bảo tìm được cực trị toàn cục}

Momentum và các kỹ thuật tối ưu hiện đại có thể giúp quá trình huấn luyện tốt hơn, nhưng không đảm bảo luôn tìm được cực trị toàn cục.

Kết quả còn phụ thuộc vào:

\begin{itemize}
    \item hình dạng hàm mất mát;
    \item cách khởi tạo tham số;
    \item learning rate;
    \item kiến trúc mạng;
    \item dữ liệu;
    \item thuật toán tối ưu.
\end{itemize}

\begin{canhbao}
Không nên hiểu rằng dùng thuật toán tối ưu hiện đại là chắc chắn tìm được nghiệm tốt nhất. Trong thực tế, ta thường tìm nghiệm đủ tốt thay vì nghiệm tối ưu tuyệt đối.
\end{canhbao}

% ==========================================================
\section{Lựa chọn hàm mất mát}

\subsection{Các hàm mất mát phổ biến}

Trong các bài ANN trước, ta đã gặp một số cách đo sai số:

\begin{enumerate}
    \item Sai số theo chuẩn $L_2$ hoặc bình phương sai số.
    \item Sai số dựa trên góc giữa hai vector.
    \item Cross entropy.
\end{enumerate}

Mỗi hàm mất mát có đặc điểm riêng và phù hợp với từng loại bài toán.

\subsection{Mean Squared Error}

Mean Squared Error, viết tắt là MSE, là sai số bình phương trung bình:

\[
    \text{MSE}
    =
    \frac{1}{N}
    \sum_{i=1}^{N}
    (y_i-\hat{y}_i)^2
\]

Với output dạng vector:

\[
    \text{MSE}
    =
    \frac{1}{N}
    \sum_{i=1}^{N}
    \|y_i-\hat{y}_i\|_2^2
\]

MSE thường được dùng trong:

\begin{itemize}
    \item bài toán hồi quy;
    \item dự đoán giá trị liên tục;
    \item một số bài toán thị giác máy tính đặc thù.
\end{itemize}

\begin{vidu}
Dự đoán giá cổ phiếu:

\[
    y = 105,\qquad \hat{y}=102
\]

Sai số bình phương:

\[
    (105-102)^2=9
\]
\end{vidu}

\subsection{Cross entropy}

Cross entropy thường được dùng trong bài toán phân loại.

Với $K$ class:

\[
    H(p,q)
    =
    -\sum_{i=1}^{K}p_i\log(q_i)
\]

Trong đó:

\begin{itemize}
    \item $p_i$ là phân phối đúng;
    \item $q_i$ là phân phối dự đoán;
    \item $K$ là số class.
\end{itemize}

Nếu $p$ là one-hot vector và class đúng là $c$, công thức rút gọn:

\[
    H(p,q)=-\log(q_c)
\]

\begin{vidu}
Ảnh thật là chó. Mạng dự đoán xác suất cho class chó là:

\[
    q_c=0.9
\]

Cross entropy:

\[
    -\log(0.9)\approx 0.105
\]

Nếu mạng chỉ dự đoán xác suất chó là:

\[
    q_c=0.1
\]

Cross entropy:

\[
    -\log(0.1)\approx 2.303
\]

Dự đoán càng tự tin đúng thì loss càng nhỏ. Dự đoán class đúng với xác suất thấp thì loss lớn.
\end{vidu}

\subsection{Vì sao cross entropy thường tốt cho phân loại?}

Trong bài toán phân loại, ta thường muốn mô hình gán xác suất cao cho class đúng. Cross entropy phạt mạnh khi xác suất của class đúng thấp.

Ví dụ, nếu label đúng là:

\[
    p =
    \begin{bmatrix}
    1\\0
    \end{bmatrix}
\]

và mô hình dự đoán:

\[
    q =
    \begin{bmatrix}
    0.5\\0.5
    \end{bmatrix}
\]

thì mô hình còn phân vân. Nếu mô hình dự đoán:

\[
    q =
    \begin{bmatrix}
    0.99\\0.01
    \end{bmatrix}
\]

thì mô hình rất chắc vào class đúng và loss sẽ nhỏ hơn nhiều.

\begin{ghinho}
Cross entropy có xu hướng khuyến khích mô hình tiến gần hơn đến phân phối nhãn đúng, đặc biệt trong bài toán phân loại.
\end{ghinho}

\subsection{Không có hàm mất mát tốt nhất cho mọi bài toán}

Dù cross entropy rất phổ biến trong phân loại, điều đó không có nghĩa là mọi nghiên cứu hoặc mọi mô hình đều dùng cross entropy.

Một số bài toán vẫn dùng MSE hoặc các biến thể của MSE. Ví dụ, trong các hệ thống phát hiện vật thể như YOLO, hàm mất mát thường là tổ hợp của nhiều thành phần, trong đó có các thành phần liên quan đến sai số bình phương cho tọa độ hộp bao.

\begin{luuy}
Chọn hàm mất mát phải dựa vào bản chất bài toán. Bài toán phân loại, hồi quy, phát hiện vật thể, phân đoạn ảnh hoặc sinh dữ liệu có thể cần các hàm loss khác nhau.
\end{luuy}

% ==========================================================
\section{So sánh các khái niệm quan trọng}

\subsection{Backpropagation và Gradient Descent}

\begin{center}
\begin{tabular}{p{0.30\textwidth}p{0.30\textwidth}p{0.30\textwidth}}
\toprule
Khái niệm & Vai trò & Trả lời câu hỏi \\
\midrule
Backpropagation & Tính gradient của loss theo các tham số & Tham số nào ảnh hưởng đến loss như thế nào? \\
Gradient Descent & Dùng gradient để cập nhật tham số & Cần tăng hay giảm tham số bao nhiêu? \\
\bottomrule
\end{tabular}
\end{center}

\subsection{Local minimum và global minimum}

\begin{center}
\begin{tabular}{p{0.35\textwidth}p{0.55\textwidth}}
\toprule
Khái niệm & Ý nghĩa \\
\midrule
Cực trị địa phương & Tốt nhất trong một vùng lân cận, nhưng có thể không phải tốt nhất toàn cục. \\
Cực trị toàn cục & Tốt nhất trên toàn bộ không gian tìm kiếm. \\
\bottomrule
\end{tabular}
\end{center}

\subsection{Exploration và exploitation}

\begin{center}
\begin{tabular}{p{0.35\textwidth}p{0.55\textwidth}}
\toprule
Khái niệm & Ý nghĩa \\
\midrule
Exploitation & Khai thác thông tin hiện có để cải thiện nghiệm hiện tại. \\
Exploration & Tìm kiếm ở các vùng mới để tránh bỏ lỡ nghiệm tốt hơn. \\
\bottomrule
\end{tabular}
\end{center}

% ==========================================================
\section{Technical Glossary -- Bảng thuật ngữ chuyên ngành}

\small
\begin{longtable}{p{0.26\textwidth}p{0.48\textwidth}p{0.19\textwidth}}
\toprule
\textbf{Thuật ngữ} & \textbf{Giải thích} & \textbf{Ví dụ} \\
\midrule
\endfirsthead

\toprule
\textbf{Thuật ngữ} & \textbf{Giải thích} & \textbf{Ví dụ} \\
\midrule
\endhead

Loss function & Hàm mất mát, đo sai khác giữa output dự đoán và label đúng. & Cross entropy, MSE. \\

Error & Sai số giữa kết quả mong muốn và kết quả hiện tại. & $e=y-\hat{y}$. \\

Optimization & Tối ưu hóa, quá trình tìm tham số để hàm mục tiêu đạt giá trị tốt nhất. & Tối thiểu hóa loss. \\

Objective function & Hàm mục tiêu cần tối ưu. & $J(\theta)$. \\

Parameter & Tham số mô hình được điều chỉnh trong huấn luyện. & Trọng số, bias. \\

Gradient Descent & Phương pháp suy giảm theo hướng ngược gradient để giảm hàm mục tiêu. & $x_{t+1}=x_t-\eta f'(x_t)$. \\

Derivative & Đạo hàm, cho biết độ dốc của hàm một biến. & $f'(x)$. \\

Gradient & Vector đạo hàm riêng của hàm nhiều biến. & $\nabla_\theta J$. \\

Partial derivative & Đạo hàm riêng theo một biến khi các biến khác xem như cố định. & $\frac{\partial J}{\partial w}$. \\

Learning rate & Hệ số quyết định độ lớn bước cập nhật. & $\eta=0.001$. \\

Convergence & Hội tụ, trạng thái thuật toán tiến dần đến nghiệm ổn định. & Loss giảm dần. \\

Divergence & Phân kỳ, trạng thái thuật toán không tiến đến nghiệm ổn định. & Loss tăng hoặc dao động. \\

Local minimum & Cực tiểu địa phương, điểm thấp nhất trong một vùng lân cận. & Hố nhỏ trên đồ thị loss. \\

Global minimum & Cực tiểu toàn cục, điểm thấp nhất trên toàn miền. & Loss nhỏ nhất có thể. \\

Exploitation & Khai thác thông tin cục bộ hiện có. & Đi theo gradient hiện tại. \\

Exploration & Khám phá vùng tìm kiếm mới. & Tìm khu vực có nghiệm tốt hơn. \\

Metaheuristic & Nhóm thuật toán tối ưu tổng quát, thường dùng cho bài toán khó. & Thuật toán di truyền. \\

Genetic Algorithm & Thuật toán di truyền, mô phỏng chọn lọc tự nhiên để tìm nghiệm tốt. & Quần thể lời giải. \\

Momentum & Kỹ thuật thêm quán tính vào cập nhật gradient. & Giúp giảm dao động. \\

MSE & Mean Squared Error, sai số bình phương trung bình. & Dự đoán giá trị liên tục. \\

Cross entropy & Hàm mất mát phổ biến cho phân loại xác suất. & $-\sum p_i\log q_i$. \\

YOLO & Hệ thống phát hiện vật thể, thường dùng hàm loss gồm nhiều thành phần. & Nhận diện vật thể trong ảnh. \\

\bottomrule
\end{longtable}
\normalsize

% ==========================================================
\section{Key Concepts Summary -- Tóm tắt kiến thức trọng tâm}

\subsection{Các ý chính cần nhớ}

\begin{enumerate}
    \item Huấn luyện mạng nơ-ron là bài toán tối ưu hóa hàm mất mát.
    \item Các tham số cần tối ưu là toàn bộ trọng số và bias trong mạng.
    \item Gradient Descent cập nhật tham số theo hướng làm giảm loss.
    \item Trong bài toán một biến, đạo hàm cho biết nên tăng hay giảm biến.
    \item Trong bài toán nhiều biến, gradient thay thế vai trò của đạo hàm.
    \item Learning rate quyết định độ lớn của mỗi bước cập nhật.
    \item Learning rate quá nhỏ làm học chậm; quá lớn có thể gây dao động hoặc không hội tụ.
    \item Hàm mất mát của ANN thường rất phức tạp và không có lời giải giải tích đơn giản.
    \item Gradient Descent có thể mắc kẹt ở cực trị địa phương.
    \item Momentum và các kỹ thuật tối ưu hiện đại giúp cải thiện khả năng hội tụ nhưng không đảm bảo tìm cực trị toàn cục.
    \item Cross entropy phổ biến trong bài toán phân loại, nhưng không phải là lựa chọn duy nhất.
    \item Chọn hàm loss phải dựa vào bản chất bài toán cụ thể.
\end{enumerate}

\subsection{Công thức quan trọng}

\subsection*{Hàm mục tiêu trên tập dữ liệu}

\[
    J(\theta)
    =
    \frac{1}{N}
    \sum_{i=1}^{N}
    \mathcal{L}
    \left(
    y_i,
    f(x_i;\theta)
    \right)
\]

\subsection*{Bài toán tối ưu}

\[
    \theta^*
    =
    \operatorname*{arg\,min}_{\theta}J(\theta)
\]

\subsection*{Gradient Descent một biến}

\[
    x_{t+1}
    =
    x_t
    -
    \eta f'(x_t)
\]

\subsection*{Gradient Descent nhiều biến}

\[
    \theta_{t+1}
    =
    \theta_t
    -
    \eta \nabla_{\theta}J(\theta_t)
\]

\subsection*{Cập nhật một trọng số}

\[
    w_{\text{new}}
    =
    w_{\text{old}}
    -
    \eta
    \frac{\partial J}{\partial w}
\]

\subsection*{MSE}

\[
    \text{MSE}
    =
    \frac{1}{N}
    \sum_{i=1}^{N}
    (y_i-\hat{y}_i)^2
\]

\subsection*{Cross entropy}

\[
    H(p,q)
    =
    -\sum_{i=1}^{K}
    p_i\log(q_i)
\]

% ==========================================================
\section{Review Questions -- Câu hỏi ôn tập}

\begin{enumerate}
    \item Vì sao huấn luyện mạng nơ-ron được xem là một bài toán tối ưu?
    \item Trong ANN, các biến cần tối ưu là gì?
    \item Hàm mất mát có vai trò gì trong quá trình huấn luyện?
    \item Trong trường hợp lý tưởng, loss nên tiến về giá trị nào?
    \item Gradient Descent dùng thông tin gì để cập nhật tham số?
    \item Vì sao công thức Gradient Descent có dấu trừ?
    \item Learning rate ảnh hưởng như thế nào đến tốc độ hội tụ?
    \item Điều gì xảy ra nếu learning rate quá nhỏ?
    \item Điều gì xảy ra nếu learning rate quá lớn?
    \item Gradient khác đạo hàm thông thường như thế nào?
    \item Cực trị địa phương khác gì cực trị toàn cục?
    \item Vì sao Gradient Descent có thể mắc kẹt ở cực trị địa phương?
    \item Exploration và exploitation khác nhau như thế nào?
    \item Thuật toán di truyền thuộc nhóm thuật toán nào?
    \item Momentum giúp ích gì trong tối ưu hóa?
    \item Cross entropy thường phù hợp với bài toán nào?
    \item MSE thường phù hợp với bài toán nào?
    \item Vì sao không có một hàm mất mát tốt nhất cho mọi bài toán?
\end{enumerate}

% ==========================================================
\section{Practice Exercises -- Bài tập vận dụng}

\subsection*{Bài tập 1: Gradient Descent một biến}
\addcontentsline{toc}{section}{Bài tập 1: Gradient Descent một biến}

Cho:

\[
    f(x) = (x-4)^2
\]

\[
    x_0 = 0,\qquad \eta=0.1
\]

\begin{enumerate}
    \item Tính đạo hàm $f'(x)$.
    \item Tính $x_1$ theo công thức Gradient Descent.
    \item Tính $x_2$.
\end{enumerate}

\subsection*{Lời giải gợi ý}

Đạo hàm:

\[
    f'(x)=2(x-4)
\]

Bước 1:

\[
    f'(0)=2(0-4)=-8
\]

\[
    x_1=0-0.1(-8)=0.8
\]

Bước 2:

\[
    f'(0.8)=2(0.8-4)=-6.4
\]

\[
    x_2=0.8-0.1(-6.4)=1.44
\]

\subsection*{Bài tập 2: Cập nhật trọng số}
\addcontentsline{toc}{section}{Bài tập 2: Cập nhật trọng số}

Cho một trọng số:

\[
    w=1.2
\]

Gradient:

\[
    \frac{\partial J}{\partial w}=0.3
\]

Learning rate:

\[
    \eta=0.05
\]

Tính trọng số mới.

\subsection*{Lời giải gợi ý}

\[
    w_{\text{new}}
    =
    w_{\text{old}}
    -
    \eta
    \frac{\partial J}{\partial w}
\]

\[
    w_{\text{new}}
    =
    1.2 - 0.05\cdot 0.3
    =
    1.2 - 0.015
    =
    1.185
\]

\subsection*{Bài tập 3: So sánh learning rate}
\addcontentsline{toc}{section}{Bài tập 3: So sánh learning rate}

Giải thích ngắn gọn hiện tượng xảy ra trong hai trường hợp:

\begin{enumerate}
    \item Learning rate rất nhỏ.
    \item Learning rate rất lớn.
\end{enumerate}

\subsection*{Lời giải gợi ý}

\begin{enumerate}
    \item Learning rate rất nhỏ làm mỗi bước cập nhật rất ngắn, thuật toán học chậm và cần nhiều vòng lặp.
    \item Learning rate rất lớn làm bước cập nhật quá dài, thuật toán có thể vượt qua cực tiểu, dao động hoặc không hội tụ.
\end{enumerate}

\subsection*{Bài tập 4: Tính cross entropy}
\addcontentsline{toc}{section}{Bài tập 4: Tính cross entropy}

Một bài toán phân loại có 3 lớp. Label đúng là lớp thứ nhất:

\[
    p =
    \begin{bmatrix}
    1\\0\\0
    \end{bmatrix}
\]

Mạng dự đoán:

\[
    q =
    \begin{bmatrix}
    0.8\\0.1\\0.1
    \end{bmatrix}
\]

Tính cross entropy.

\subsection*{Lời giải gợi ý}

Vì class đúng là class thứ nhất:

\[
    H(p,q)=-\log(0.8)
\]

\[
    H(p,q)\approx 0.223
\]

\subsection*{Bài tập 5: Nhận diện loại bài toán}
\addcontentsline{toc}{section}{Bài tập 5: Nhận diện loại bài toán}

Với mỗi bài toán sau, hãy đề xuất hàm mất mát phù hợp hơn: MSE hoặc cross entropy.

\begin{enumerate}
    \item Dự đoán giá nhà.
    \item Phân loại ảnh chó, mèo, chim.
    \item Dự đoán nhiệt độ ngày mai.
    \item Phân loại email spam hoặc không spam.
\end{enumerate}

\subsection*{Lời giải gợi ý}

\begin{center}
\begin{tabular}{lll}
\toprule
Bài toán & Loại bài toán & Hàm mất mát gợi ý \\
\midrule
Dự đoán giá nhà & Hồi quy & MSE \\
Phân loại ảnh & Phân loại nhiều lớp & Cross entropy \\
Dự đoán nhiệt độ & Hồi quy & MSE \\
Phân loại email & Phân loại nhị phân & Binary cross entropy \\
\bottomrule
\end{tabular}
\end{center}

% ==========================================================
\section*{Kết luận}
\addcontentsline{toc}{section}{Kết luận}

ANN4 làm rõ bản chất tối ưu hóa của quá trình huấn luyện mạng nơ-ron. Khi mạng tạo ra output khác với label đúng, hàm mất mát được dùng để đo sai số. Mục tiêu của huấn luyện là điều chỉnh toàn bộ trọng số và bias sao cho hàm mất mát giảm dần.

Gradient Descent là phương pháp cơ bản để cập nhật tham số theo hướng làm giảm loss. Trong trường hợp một biến, ta dùng đạo hàm; trong trường hợp nhiều biến, ta dùng gradient. Learning rate quyết định độ lớn bước cập nhật và ảnh hưởng mạnh đến quá trình hội tụ.

Bên cạnh đó, bài học cũng giới thiệu các khó khăn trong tối ưu hóa như cực trị địa phương, sự cần thiết của exploration và exploitation, cũng như ý tưởng của momentum và thuật toán di truyền. Cuối cùng, việc lựa chọn hàm mất mát phải gắn với bản chất bài toán: cross entropy thường dùng cho phân loại, còn MSE thường dùng cho hồi quy hoặc một số bài toán đặc thù.



\chapter{ANN6: Backpropagation Example -- Ví dụ lan truyền ngược}


% ==========================================================
\section*{Learning Objectives -- Mục tiêu của tài liệu}
\addcontentsline{toc}{section}{Mục tiêu của tài liệu}

Tài liệu này chỉ tập trung giải thích ví dụ tính toán trong video ANN6. Nội dung không trình bày lại lý thuyết tổng quát, mà đi theo đúng tinh thần của ví dụ:

\begin{enumerate}[label=\arabic*.]
    \item Tính xuôi để lấy output và sai số tổng.
    \item Tính ngược để cập nhật các trọng số từ lớp output: $w_5,w_6,w_7,w_8$.
    \item Tính ngược tiếp để cập nhật các trọng số phía trước: $w_1,w_2,w_3,w_4$.
    \item Làm rõ điểm dễ sai nhất: khi tính đạo hàm theo $w_1$, output của hidden neuron $h_1$ ảnh hưởng đồng thời đến cả $o_1$ và $o_2$, nên phải cộng cả hai nhánh.
\end{enumerate}

% ==========================================================
\section{Example Description -- Mô tả ví dụ}

\subsection{Cấu trúc mạng trong ví dụ}

Ví dụ sử dụng một mạng nơ-ron rất nhỏ:

\begin{itemize}
    \item 2 input: $i_1,i_2$;
    \item 2 neuron ẩn: $h_1,h_2$;
    \item 2 output: $o_1,o_2$.
\end{itemize}

\begin{center}
\begin{tikzpicture}[x=2.4cm,y=1.2cm,>=Stealth,
    neuron/.style={circle,draw,minimum size=9mm},
    every node/.style={font=\small}
]

% Input
\node[neuron] (i1) at (0,0.8) {$i_1$};
\node[neuron] (i2) at (0,-0.8) {$i_2$};

% Hidden
\node[neuron] (h1) at (1.8,0.8) {$h_1$};
\node[neuron] (h2) at (1.8,-0.8) {$h_2$};

% Output
\node[neuron] (o1) at (3.6,0.8) {$o_1$};
\node[neuron] (o2) at (3.6,-0.8) {$o_2$};

% Connections input-hidden
\draw[->] (i1) -- node[above] {$w_1$} (h1);
\draw[->] (i2) -- node[below] {$w_2$} (h1);
\draw[->] (i1) -- node[above] {$w_3$} (h2);
\draw[->] (i2) -- node[below] {$w_4$} (h2);

% Connections hidden-output
\draw[->] (h1) -- node[above] {$w_5$} (o1);
\draw[->] (h2) -- node[below] {$w_6$} (o1);
\draw[->] (h1) -- node[above] {$w_7$} (o2);
\draw[->] (h2) -- node[below] {$w_8$} (o2);

\node at (0,1.8) {\textbf{Input}};
\node at (1.8,1.8) {\textbf{Hidden}};
\node at (3.6,1.8) {\textbf{Output}};

\end{tikzpicture}
\end{center}

\subsection{Dữ liệu ban đầu}

Trong ví dụ, ta dùng các giá trị sau:

\begin{center}
\begin{tabular}{lll}
\toprule
Đại lượng & Giá trị & Ý nghĩa \\
\midrule
$i_1$ & $0.05$ & Input thứ nhất \\
$i_2$ & $0.10$ & Input thứ hai \\
$t_1$ & $0.01$ & Target mong muốn của output $o_1$ \\
$t_2$ & $0.99$ & Target mong muốn của output $o_2$ \\
$b_h$ & $0.35$ & Bias cho lớp hidden \\
$b_o$ & $0.60$ & Bias cho lớp output \\
$\eta$ & $0.5$ & Learning rate \\
\bottomrule
\end{tabular}
\end{center}

Các trọng số ban đầu:

\begin{center}
\begin{tabular}{cccccccc}
\toprule
$w_1$ & $w_2$ & $w_3$ & $w_4$ & $w_5$ & $w_6$ & $w_7$ & $w_8$ \\
\midrule
0.15 & 0.20 & 0.25 & 0.30 & 0.40 & 0.45 & 0.50 & 0.55 \\
\bottomrule
\end{tabular}
\end{center}

Hàm kích hoạt được dùng trong ví dụ là sigmoid:

\[
    \sigma(x)=\frac{1}{1+e^{-x}}
\]

Đạo hàm sigmoid tại output của nó:

\[
    \sigma'(x)=\sigma(x)\left(1-\sigma(x)\right)
\]

Trong khi tính, nếu đã biết output $a=\sigma(x)$, ta dùng nhanh:

\[
    \frac{da}{dx}=a(1-a)
\]

% ==========================================================
\section{Forward Pass Example -- Tính xuôi trong ví dụ}

\subsection{Tính output của hidden neuron $h_1$}

Net input vào $h_1$:

\[
    net_{h_1}=i_1w_1+i_2w_2+b_h
\]

Thay số:

\[
    net_{h_1}
    =
    0.05\cdot 0.15
    +
    0.10\cdot 0.20
    +
    0.35
\]

\[
    net_{h_1}=0.3775
\]

Output của $h_1$:

\[
    out_{h_1}=\sigma(0.3775)
\]

\[
    out_{h_1}\approx 0.593269992
\]

\subsection{Tính output của hidden neuron $h_2$}

Net input vào $h_2$:

\[
    net_{h_2}=i_1w_3+i_2w_4+b_h
\]

Thay số:

\[
    net_{h_2}
    =
    0.05\cdot 0.25
    +
    0.10\cdot 0.30
    +
    0.35
\]

\[
    net_{h_2}=0.3925
\]

Output của $h_2$:

\[
    out_{h_2}=\sigma(0.3925)
\]

\[
    out_{h_2}\approx 0.596884378
\]

\subsection{Tính output $o_1$}

Net input vào $o_1$:

\[
    net_{o_1}=out_{h_1}w_5+out_{h_2}w_6+b_o
\]

Thay số:

\[
    net_{o_1}
    =
    0.593269992\cdot 0.40
    +
    0.596884378\cdot 0.45
    +
    0.60
\]

\[
    net_{o_1}\approx 1.105905967
\]

Output:

\[
    out_{o_1}=\sigma(1.105905967)
\]

\[
    out_{o_1}\approx 0.751365070
\]

\subsection{Tính output $o_2$}

Net input vào $o_2$:

\[
    net_{o_2}=out_{h_1}w_7+out_{h_2}w_8+b_o
\]

Thay số:

\[
    net_{o_2}
    =
    0.593269992\cdot 0.50
    +
    0.596884378\cdot 0.55
    +
    0.60
\]

\[
    net_{o_2}\approx 1.224921404
\]

Output:

\[
    out_{o_2}=\sigma(1.224921404)
\]

\[
    out_{o_2}\approx 0.772928465
\]

% ==========================================================
\section{Error Calculation Example -- Tính sai số của ví dụ}

\subsection{Sai số tại output $o_1$}

Target của $o_1$ là:

\[
    t_1=0.01
\]

Output hiện tại:

\[
    out_{o_1}=0.751365070
\]

Sai số tại $o_1$:

\[
    E_{o_1}
    =
    \frac{1}{2}
    (t_1-out_{o_1})^2
\]

Thay số:

\[
    E_{o_1}
    =
    \frac{1}{2}
    (0.01-0.751365070)^2
\]

\[
    E_{o_1}\approx 0.274811083
\]

\subsection{Sai số tại output $o_2$}

Target của $o_2$ là:

\[
    t_2=0.99
\]

Output hiện tại:

\[
    out_{o_2}=0.772928465
\]

Sai số tại $o_2$:

\[
    E_{o_2}
    =
    \frac{1}{2}
    (t_2-out_{o_2})^2
\]

Thay số:

\[
    E_{o_2}
    =
    \frac{1}{2}
    (0.99-0.772928465)^2
\]

\[
    E_{o_2}\approx 0.023560026
\]

\subsection{Sai số tổng}

\[
    E_{total}=E_{o_1}+E_{o_2}
\]

\[
    E_{total}
    =
    0.274811083
    +
    0.023560026
\]

\[
    E_{total}\approx 0.298371109
\]

Làm tròn:

\[
    E_{total}\approx 0.2984
\]

\begin{ghinho}
Trong video, giá trị sai số được nhắc đến là khoảng $0.2984$. Đây chính là sai số tổng sau bước tính xuôi đầu tiên.
\end{ghinho}

% ==========================================================
\section{Updating Output Layer Weights -- Cập nhật các trọng số lớp output}

\subsection{Cập nhật $w_5$}

Trọng số $w_5$ nối từ $h_1$ đến $o_1$.

Ta cần tính:

\[
    \frac{\partial E_{total}}{\partial w_5}
\]

Vì $w_5$ chỉ đi trực tiếp vào $o_1$, nên:

\[
    \frac{\partial E_{total}}{\partial w_5}
    =
    \frac{\partial E_{o_1}}{\partial out_{o_1}}
    \cdot
    \frac{\partial out_{o_1}}{\partial net_{o_1}}
    \cdot
    \frac{\partial net_{o_1}}{\partial w_5}
\]

Từng thành phần:

\[
    \frac{\partial E_{o_1}}{\partial out_{o_1}}
    =
    out_{o_1}-t_1
\]

\[
    =
    0.751365070-0.01
    =
    0.741365070
\]

\[
    \frac{\partial out_{o_1}}{\partial net_{o_1}}
    =
    out_{o_1}(1-out_{o_1})
\]

\[
    =
    0.751365070(1-0.751365070)
    \approx 0.186815602
\]

\[
    \frac{\partial net_{o_1}}{\partial w_5}
    =
    out_{h_1}
    =
    0.593269992
\]

Vậy:

\[
    \frac{\partial E_{total}}{\partial w_5}
    =
    0.741365070
    \cdot
    0.186815602
    \cdot
    0.593269992
\]

\[
    \frac{\partial E_{total}}{\partial w_5}
    \approx 0.082167041
\]

Cập nhật:

\[
    w_5^{new}
    =
    w_5
    -
    \eta
    \frac{\partial E_{total}}{\partial w_5}
\]

\[
    w_5^{new}
    =
    0.40
    -
    0.5\cdot 0.082167041
\]

\[
    w_5^{new}
    \approx 0.358916480
\]

\subsection{Cập nhật $w_6,w_7,w_8$}

Cách làm tương tự như $w_5$.

\subsection*{Với $w_6$}

$w_6$ nối từ $h_2$ đến $o_1$:

\[
    \frac{\partial E_{total}}{\partial w_6}
    =
    \frac{\partial E_{o_1}}{\partial out_{o_1}}
    \cdot
    \frac{\partial out_{o_1}}{\partial net_{o_1}}
    \cdot
    \frac{\partial net_{o_1}}{\partial w_6}
\]

Trong đó:

\[
    \frac{\partial net_{o_1}}{\partial w_6}=out_{h_2}
\]

\[
    \frac{\partial E_{total}}{\partial w_6}
    =
    0.741365070
    \cdot
    0.186815602
    \cdot
    0.596884378
\]

\[
    \frac{\partial E_{total}}{\partial w_6}
    \approx 0.082667628
\]

\[
    w_6^{new}
    =
    0.45
    -
    0.5\cdot 0.082667628
\]

\[
    w_6^{new}\approx 0.408666186
\]

\subsection*{Với $w_7$}

$w_7$ nối từ $h_1$ đến $o_2$:

\[
    \frac{\partial E_{total}}{\partial w_7}
    =
    \frac{\partial E_{o_2}}{\partial out_{o_2}}
    \cdot
    \frac{\partial out_{o_2}}{\partial net_{o_2}}
    \cdot
    \frac{\partial net_{o_2}}{\partial w_7}
\]

\[
    \frac{\partial E_{o_2}}{\partial out_{o_2}}
    =
    out_{o_2}-t_2
    =
    0.772928465-0.99
\]

\[
    =
    -0.217071535
\]

\[
    \frac{\partial out_{o_2}}{\partial net_{o_2}}
    =
    out_{o_2}(1-out_{o_2})
\]

\[
    =
    0.772928465(1-0.772928465)
    \approx 0.175510052
\]

\[
    \frac{\partial net_{o_2}}{\partial w_7}
    =
    out_{h_1}
    =
    0.593269992
\]

\[
    \frac{\partial E_{total}}{\partial w_7}
    =
    -0.217071535
    \cdot
    0.175510052
    \cdot
    0.593269992
\]

\[
    \frac{\partial E_{total}}{\partial w_7}
    \approx -0.022602540
\]

\[
    w_7^{new}
    =
    0.50
    -
    0.5(-0.022602540)
\]

\[
    w_7^{new}
    \approx 0.511301270
\]

\subsection*{Với $w_8$}

$w_8$ nối từ $h_2$ đến $o_2$:

\[
    \frac{\partial net_{o_2}}{\partial w_8}
    =
    out_{h_2}
    =
    0.596884378
\]

\[
    \frac{\partial E_{total}}{\partial w_8}
    =
    -0.217071535
    \cdot
    0.175510052
    \cdot
    0.596884378
\]

\[
    \frac{\partial E_{total}}{\partial w_8}
    \approx -0.022740242
\]

\[
    w_8^{new}
    =
    0.55
    -
    0.5(-0.022740242)
\]

\[
    w_8^{new}
    \approx 0.561370121
\]

\subsection{Bảng kết quả cập nhật lớp output}

\begin{center}
\begin{tabular}{cccc}
\toprule
Trọng số & Giá trị cũ & Gradient & Giá trị mới \\
\midrule
$w_5$ & 0.400000000 & 0.082167041 & 0.358916480 \\
$w_6$ & 0.450000000 & 0.082667628 & 0.408666186 \\
$w_7$ & 0.500000000 & -0.022602540 & 0.511301270 \\
$w_8$ & 0.550000000 & -0.022740242 & 0.561370121 \\
\bottomrule
\end{tabular}
\end{center}

% ==========================================================
\section{Cập nhật bias ở lớp output}

\subsection{Bias ứng với output $o_1$}

Với bias của $o_1$:

\[
    \frac{\partial net_{o_1}}{\partial b_{o_1}}=1
\]

Do đó:

\[
    \frac{\partial E_{total}}{\partial b_{o_1}}
    =
    \frac{\partial E_{o_1}}{\partial out_{o_1}}
    \cdot
    \frac{\partial out_{o_1}}{\partial net_{o_1}}
\]

\[
    =
    0.741365070
    \cdot
    0.186815602
\]

\[
    \frac{\partial E_{total}}{\partial b_{o_1}}
    \approx 0.138498562
\]

Nếu bias ban đầu là $0.60$:

\[
    b_{o_1}^{new}
    =
    0.60
    -
    0.5\cdot 0.138498562
\]

\[
    b_{o_1}^{new}
    \approx 0.530750719
\]

\subsection{Bias ứng với output $o_2$}

\[
    \frac{\partial E_{total}}{\partial b_{o_2}}
    =
    \frac{\partial E_{o_2}}{\partial out_{o_2}}
    \cdot
    \frac{\partial out_{o_2}}{\partial net_{o_2}}
\]

\[
    =
    -0.217071535
    \cdot
    0.175510052
\]

\[
    \frac{\partial E_{total}}{\partial b_{o_2}}
    \approx -0.038098237
\]

\[
    b_{o_2}^{new}
    =
    0.60
    -
    0.5(-0.038098237)
\]

\[
    b_{o_2}^{new}
    \approx 0.619049118
\]

\begin{luuy}
Nếu trong tài liệu gốc dùng chung một bias $b_o$ cho cả hai output, cách cập nhật có thể được trình bày khác. Trong ví dụ học tập này, ta tách thành $b_{o_1}$ và $b_{o_2}$ để thấy rõ mỗi output có một nhánh sai số riêng.
\end{luuy}

% ==========================================================
\section{Cập nhật các trọng số phía trước}

\subsection{Điểm quan trọng khi tính $w_1$}

Trọng số $w_1$ nối từ input $i_1$ đến hidden neuron $h_1$.

Khi $w_1$ thay đổi, chuỗi ảnh hưởng không chỉ là:

\[
    w_1
    \rightarrow
    net_{h_1}
    \rightarrow
    out_{h_1}
    \rightarrow
    net_{o_1}
    \rightarrow
    out_{o_1}
    \rightarrow
    E
\]

Chuỗi này chưa đủ.

Vì $out_{h_1}$ đi đến cả $o_1$ và $o_2$, nên $w_1$ ảnh hưởng đến sai số qua cả hai nhánh:

\[
    out_{h_1}
    \rightarrow
    o_1
    \rightarrow
    E_{o_1}
\]

và:

\[
    out_{h_1}
    \rightarrow
    o_2
    \rightarrow
    E_{o_2}
\]

\begin{ghinho}
Đây là điểm quan trọng nhất trong ví dụ ANN6: khi tính đạo hàm theo $w_1$, phải cộng ảnh hưởng qua cả $o_1$ và $o_2$. Nếu chỉ lấy nhánh qua $o_1$, kết quả sẽ sai.
\end{ghinho}

\subsection{Sơ đồ đúng khi tính đạo hàm theo $w_1$}

\begin{center}
\begin{tikzpicture}[x=1.8cm,y=1.1cm,>=Stealth,
    box/.style={rectangle,draw,rounded corners,align=center,minimum width=0.5cm,minimum height=0.7cm},
    every node/.style={font=\small}
]

\node[box] (w1) at (0,0) {$w_1$};
\node[box] (neth1) at (1.4,0) {$net_{h_1}$};
\node[box] (outh1) at (2.9,0) {$out_{h_1}$};

\node[box] (neto1) at (4.4,0.8) {$net_{o_1}$};
\node[box] (outo1) at (5.9,0.8) {$out_{o_1}$};
\node[box] (eo1) at (7.4,0.8) {$E_{o_1}$};

\node[box] (neto2) at (4.4,-0.8) {$net_{o_2}$};
\node[box] (outo2) at (5.9,-0.8) {$out_{o_2}$};
\node[box] (eo2) at (7.4,-0.8) {$E_{o_2}$};

\node[box] (et) at (8.9,0) {$E_{total}$};

\draw[->] (w1) -- (neth1);
\draw[->] (neth1) -- (outh1);

\draw[->] (outh1) -- (neto1);
\draw[->] (neto1) -- (outo1);
\draw[->] (outo1) -- (eo1);

\draw[->] (outh1) -- (neto2);
\draw[->] (neto2) -- (outo2);
\draw[->] (outo2) -- (eo2);

\draw[->] (eo1) -- (et);
\draw[->] (eo2) -- (et);

\end{tikzpicture}
\end{center}

\subsection{Tính $\frac{\partial E_{total}}{\partial out_{h_1}}$}

Vì $out_{h_1}$ ảnh hưởng đến cả $o_1$ và $o_2$:

\[
    \frac{\partial E_{total}}{\partial out_{h_1}}
    =
    \frac{\partial E_{o_1}}{\partial out_{h_1}}
    +
    \frac{\partial E_{o_2}}{\partial out_{h_1}}
\]

Nhánh qua $o_1$:

\[
    \frac{\partial E_{o_1}}{\partial out_{h_1}}
    =
    \frac{\partial E_{o_1}}{\partial net_{o_1}}
    \cdot
    \frac{\partial net_{o_1}}{\partial out_{h_1}}
\]

Ta đã có:

\[
    \frac{\partial E_{o_1}}{\partial net_{o_1}}
    =
    0.138498562
\]

và:

\[
    \frac{\partial net_{o_1}}{\partial out_{h_1}}
    =
    w_5
    =
    0.40
\]

Do đó:

\[
    \frac{\partial E_{o_1}}{\partial out_{h_1}}
    =
    0.138498562\cdot 0.40
\]

\[
    \frac{\partial E_{o_1}}{\partial out_{h_1}}
    \approx 0.055399425
\]

Nhánh qua $o_2$:

\[
    \frac{\partial E_{o_2}}{\partial out_{h_1}}
    =
    \frac{\partial E_{o_2}}{\partial net_{o_2}}
    \cdot
    \frac{\partial net_{o_2}}{\partial out_{h_1}}
\]

Ta đã có:

\[
    \frac{\partial E_{o_2}}{\partial net_{o_2}}
    =
    -0.038098237
\]

và:

\[
    \frac{\partial net_{o_2}}{\partial out_{h_1}}
    =
    w_7
    =
    0.50
\]

Do đó:

\[
    \frac{\partial E_{o_2}}{\partial out_{h_1}}
    =
    -0.038098237\cdot 0.50
\]

\[
    \frac{\partial E_{o_2}}{\partial out_{h_1}}
    \approx -0.019049119
\]

Cộng hai nhánh:

\[
    \frac{\partial E_{total}}{\partial out_{h_1}}
    =
    0.055399425
    +
    (-0.019049119)
\]

\[
    \frac{\partial E_{total}}{\partial out_{h_1}}
    \approx 0.036350306
\]

\subsection{Tính gradient của $w_1$}

Ta cần:

\[
    \frac{\partial E_{total}}{\partial w_1}
\]

Theo chuỗi:

\[
    \frac{\partial E_{total}}{\partial w_1}
    =
    \frac{\partial E_{total}}{\partial out_{h_1}}
    \cdot
    \frac{\partial out_{h_1}}{\partial net_{h_1}}
    \cdot
    \frac{\partial net_{h_1}}{\partial w_1}
\]

Trong đó:

\[
    \frac{\partial E_{total}}{\partial out_{h_1}}
    \approx 0.036350306
\]

\[
    \frac{\partial out_{h_1}}{\partial net_{h_1}}
    =
    out_{h_1}(1-out_{h_1})
\]

\[
    =
    0.593269992(1-0.593269992)
\]

\[
    \frac{\partial out_{h_1}}{\partial net_{h_1}}
    \approx 0.241300709
\]

Và:

\[
    \frac{\partial net_{h_1}}{\partial w_1}=i_1=0.05
\]

Vậy:

\[
    \frac{\partial E_{total}}{\partial w_1}
    =
    0.036350306
    \cdot
    0.241300709
    \cdot
    0.05
\]

\[
    \frac{\partial E_{total}}{\partial w_1}
    \approx 0.000438568
\]

Cập nhật:

\[
    w_1^{new}
    =
    0.15
    -
    0.5\cdot 0.000438568
\]

\[
    w_1^{new}
    \approx 0.149780716
\]

\subsection{Cập nhật $w_2$}

$w_2$ cũng đi vào $h_1$, nên phần đầu giống $w_1$. Điểm khác là:

\[
    \frac{\partial net_{h_1}}{\partial w_2}=i_2=0.10
\]

Do đó:

\[
    \frac{\partial E_{total}}{\partial w_2}
    =
    0.036350306
    \cdot
    0.241300709
    \cdot
    0.10
\]

\[
    \frac{\partial E_{total}}{\partial w_2}
    \approx 0.000877135
\]

Cập nhật:

\[
    w_2^{new}
    =
    0.20
    -
    0.5\cdot 0.000877135
\]

\[
    w_2^{new}
    \approx 0.199561432
\]

% ==========================================================
\section{Cập nhật $w_3$ và $w_4$}

\subsection{Tính ảnh hưởng qua hidden neuron $h_2$}

Tương tự $h_1$, output $out_{h_2}$ cũng đi đến cả $o_1$ và $o_2$.

Do đó:

\[
    \frac{\partial E_{total}}{\partial out_{h_2}}
    =
    \frac{\partial E_{o_1}}{\partial out_{h_2}}
    +
    \frac{\partial E_{o_2}}{\partial out_{h_2}}
\]

Nhánh qua $o_1$:

\[
    \frac{\partial E_{o_1}}{\partial out_{h_2}}
    =
    \frac{\partial E_{o_1}}{\partial net_{o_1}}
    \cdot
    \frac{\partial net_{o_1}}{\partial out_{h_2}}
\]

\[
    =
    0.138498562\cdot w_6
\]

\[
    =
    0.138498562\cdot 0.45
\]

\[
    \approx 0.062324353
\]

Nhánh qua $o_2$:

\[
    \frac{\partial E_{o_2}}{\partial out_{h_2}}
    =
    \frac{\partial E_{o_2}}{\partial net_{o_2}}
    \cdot
    \frac{\partial net_{o_2}}{\partial out_{h_2}}
\]

\[
    =
    -0.038098237\cdot w_8
\]

\[
    =
    -0.038098237\cdot 0.55
\]

\[
    \approx -0.020954030
\]

Cộng lại:

\[
    \frac{\partial E_{total}}{\partial out_{h_2}}
    =
    0.062324353
    -
    0.020954030
\]

\[
    \frac{\partial E_{total}}{\partial out_{h_2}}
    \approx 0.041370323
\]

\subsection{Cập nhật $w_3$}

\[
    \frac{\partial E_{total}}{\partial w_3}
    =
    \frac{\partial E_{total}}{\partial out_{h_2}}
    \cdot
    \frac{\partial out_{h_2}}{\partial net_{h_2}}
    \cdot
    \frac{\partial net_{h_2}}{\partial w_3}
\]

Ta có:

\[
    \frac{\partial out_{h_2}}{\partial net_{h_2}}
    =
    out_{h_2}(1-out_{h_2})
\]

\[
    =
    0.596884378(1-0.596884378)
\]

\[
    \approx 0.240613417
\]

Và:

\[
    \frac{\partial net_{h_2}}{\partial w_3}=i_1=0.05
\]

Vậy:

\[
    \frac{\partial E_{total}}{\partial w_3}
    =
    0.041370323
    \cdot
    0.240613417
    \cdot
    0.05
\]

\[
    \frac{\partial E_{total}}{\partial w_3}
    \approx 0.000497713
\]

Cập nhật:

\[
    w_3^{new}
    =
    0.25
    -
    0.5\cdot 0.000497713
\]

\[
    w_3^{new}
    \approx 0.249751144
\]

\subsection{Cập nhật $w_4$}

Với $w_4$:

\[
    \frac{\partial net_{h_2}}{\partial w_4}=i_2=0.10
\]

\[
    \frac{\partial E_{total}}{\partial w_4}
    =
    0.041370323
    \cdot
    0.240613417
    \cdot
    0.10
\]

\[
    \frac{\partial E_{total}}{\partial w_4}
    \approx 0.000995425
\]

Cập nhật:

\[
    w_4^{new}
    =
    0.30
    -
    0.5\cdot 0.000995425
\]

\[
    w_4^{new}
    \approx 0.299502287
\]

\subsection{Bảng kết quả cập nhật lớp hidden}

\begin{center}
\begin{tabular}{cccc}
\toprule
Trọng số & Giá trị cũ & Gradient & Giá trị mới \\
\midrule
$w_1$ & 0.150000000 & 0.000438568 & 0.149780716 \\
$w_2$ & 0.200000000 & 0.000877135 & 0.199561432 \\
$w_3$ & 0.250000000 & 0.000497713 & 0.249751144 \\
$w_4$ & 0.300000000 & 0.000995425 & 0.299502287 \\
\bottomrule
\end{tabular}
\end{center}

% ==========================================================
\section{Cập nhật bias ở lớp hidden}

\subsection{Bias ứng với $h_1$}

Vì:

\[
    \frac{\partial net_{h_1}}{\partial b_{h_1}}=1
\]

nên:

\[
    \frac{\partial E_{total}}{\partial b_{h_1}}
    =
    \frac{\partial E_{total}}{\partial out_{h_1}}
    \cdot
    \frac{\partial out_{h_1}}{\partial net_{h_1}}
\]

\[
    =
    0.036350306
    \cdot
    0.241300709
\]

\[
    \frac{\partial E_{total}}{\partial b_{h_1}}
    \approx 0.008771355
\]

Cập nhật:

\[
    b_{h_1}^{new}
    =
    0.35
    -
    0.5\cdot 0.008771355
\]

\[
    b_{h_1}^{new}
    \approx 0.345614323
\]

\subsection{Bias ứng với $h_2$}

\[
    \frac{\partial E_{total}}{\partial b_{h_2}}
    =
    \frac{\partial E_{total}}{\partial out_{h_2}}
    \cdot
    \frac{\partial out_{h_2}}{\partial net_{h_2}}
\]

\[
    =
    0.041370323
    \cdot
    0.240613417
\]

\[
    \frac{\partial E_{total}}{\partial b_{h_2}}
    \approx 0.009954255
\]

Cập nhật:

\[
    b_{h_2}^{new}
    =
    0.35
    -
    0.5\cdot 0.009954255
\]

\[
    b_{h_2}^{new}
    \approx 0.345022873
\]

\begin{luuy}
Tương tự bias ở lớp output, nếu tài liệu gốc dùng chung một bias cho cả lớp hidden thì cách ghi có thể khác. Việc tách $b_{h_1}$ và $b_{h_2}$ giúp dễ theo dõi từng nhánh tính toán.
\end{luuy}

% ==========================================================
\section{Tổng kết một lượt cập nhật}

\subsection{Bảng trọng số trước và sau cập nhật}

\begin{center}
\begin{tabular}{ccc}
\toprule
Trọng số & Trước cập nhật & Sau cập nhật \\
\midrule
$w_1$ & 0.150000000 & 0.149780716 \\
$w_2$ & 0.200000000 & 0.199561432 \\
$w_3$ & 0.250000000 & 0.249751144 \\
$w_4$ & 0.300000000 & 0.299502287 \\
$w_5$ & 0.400000000 & 0.358916480 \\
$w_6$ & 0.450000000 & 0.408666186 \\
$w_7$ & 0.500000000 & 0.511301270 \\
$w_8$ & 0.550000000 & 0.561370121 \\
\bottomrule
\end{tabular}
\end{center}

\subsection{Bảng bias trước và sau cập nhật}

\begin{center}
\begin{tabular}{ccc}
\toprule
Bias & Trước cập nhật & Sau cập nhật \\
\midrule
$b_{o_1}$ & 0.600000000 & 0.530750719 \\
$b_{o_2}$ & 0.600000000 & 0.619049118 \\
$b_{h_1}$ & 0.350000000 & 0.345614323 \\
$b_{h_2}$ & 0.350000000 & 0.345022873 \\
\bottomrule
\end{tabular}
\end{center}

\subsection{Ý nghĩa của một lượt cập nhật}

Sau khi cập nhật xong tất cả trọng số và bias, mạng đã thay đổi một chút so với trạng thái ban đầu. Đây mới chỉ là một lượt cập nhật với một cặp dữ liệu.

Với cặp dữ liệu tiếp theo, ta lại làm:

\[
    \text{tính xuôi}
    \rightarrow
    \text{tính sai số}
    \rightarrow
    \text{tính ngược}
    \rightarrow
    \text{cập nhật trọng số}
\]

Quá trình này lặp đi lặp lại rất nhiều lần cho đến khi sai số giảm đến mức chấp nhận được.

% ==========================================================
\section{Những lỗi dễ mắc trong ví dụ}

\subsection{Lỗi 1: Chỉ tính một nhánh khi cập nhật $w_1$}

Sai lầm thường gặp là viết:

\[
    \frac{\partial E_{total}}{\partial w_1}
    =
    \frac{\partial E_{o_1}}{\partial out_{o_1}}
    \cdot
    \frac{\partial out_{o_1}}{\partial net_{o_1}}
    \cdot
    \frac{\partial net_{o_1}}{\partial out_{h_1}}
    \cdot
    \frac{\partial out_{h_1}}{\partial net_{h_1}}
    \cdot
    \frac{\partial net_{h_1}}{\partial w_1}
\]

Biểu thức này chỉ đi qua nhánh $o_1$, nên thiếu nhánh $o_2$.

Biểu thức đúng phải là:

\[
    \frac{\partial E_{total}}{\partial w_1}
    =
    \left(
    \frac{\partial E_{o_1}}{\partial out_{h_1}}
    +
    \frac{\partial E_{o_2}}{\partial out_{h_1}}
    \right)
    \cdot
    \frac{\partial out_{h_1}}{\partial net_{h_1}}
    \cdot
    \frac{\partial net_{h_1}}{\partial w_1}
\]

\subsection{Lỗi 2: Nhầm dấu khi cập nhật}

Công thức cập nhật là:

\[
    w^{new}=w-\eta\frac{\partial E}{\partial w}
\]

Nếu gradient dương, trọng số giảm.

Nếu gradient âm, trọng số tăng.

Ví dụ:

\[
    \frac{\partial E}{\partial w_7}\approx -0.022602540
\]

Do đó:

\[
    w_7^{new}
    =
    0.50
    -
    0.5(-0.022602540)
\]

\[
    w_7^{new}>0.50
\]

Vì gradient âm nên $w_7$ tăng.

\subsection{Lỗi 3: Dùng nhầm output của neuron}

Khi tính gradient của $w_5$:

\[
    \frac{\partial net_{o_1}}{\partial w_5}=out_{h_1}
\]

không phải $out_{h_2}$.

Khi tính gradient của $w_6$:

\[
    \frac{\partial net_{o_1}}{\partial w_6}=out_{h_2}
\]

không phải $out_{h_1}$.

\begin{luuy}
Mỗi trọng số nằm trên một cạnh cụ thể. Đạo hàm của $net$ theo trọng số đó bằng đúng output của neuron ở đầu vào của cạnh đó.
\end{luuy}

% ==========================================================
\section{Bài tập tự kiểm tra theo ví dụ}

\subsection*{Bài tập 1: Kiểm tra sai số tổng}
\addcontentsline{toc}{section}{Bài tập 1: Kiểm tra sai số tổng}

Dùng các giá trị:

\[
    out_{o_1}=0.751365070,\quad t_1=0.01
\]

\[
    out_{o_2}=0.772928465,\quad t_2=0.99
\]

Hãy tính:

\[
    E_{total}
    =
    \frac{1}{2}(t_1-out_{o_1})^2
    +
    \frac{1}{2}(t_2-out_{o_2})^2
\]

\subsection*{Đáp án}

\[
    E_{total}\approx 0.298371109
\]

\[
    E_{total}\approx 0.2984
\]

\subsection*{Bài tập 2: Tính lại $w_5^{new}$}
\addcontentsline{toc}{section}{Bài tập 2: Tính lại $w_5^{new}$}

Cho:

\[
    \frac{\partial E}{\partial w_5}=0.082167041
\]

\[
    w_5=0.40,\qquad \eta=0.5
\]

Tính $w_5^{new}$.

\subsection*{Đáp án}

\[
    w_5^{new}
    =
    0.40
    -
    0.5\cdot 0.082167041
\]

\[
    w_5^{new}
    \approx 0.358916480
\]

\subsection*{Bài tập 3: Vì sao $w_7$ tăng?}
\addcontentsline{toc}{section}{Bài tập 3: Vì sao $w_7$ tăng?}

Cho:

\[
    \frac{\partial E}{\partial w_7}\approx -0.022602540
\]

\[
    w_7=0.50,\qquad \eta=0.5
\]

Tính $w_7^{new}$ và giải thích vì sao nó tăng.

\subsection*{Đáp án}

\[
    w_7^{new}
    =
    0.50
    -
    0.5(-0.022602540)
\]

\[
    w_7^{new}
    =
    0.50
    +
    0.011301270
\]

\[
    w_7^{new}
    \approx 0.511301270
\]

Vì gradient âm nên khi trừ gradient, giá trị trọng số tăng.

\subsection*{Bài tập 4: Viết đúng biểu thức cho $w_1$}
\addcontentsline{toc}{section}{Bài tập 4: Viết đúng biểu thức cho $w_1$}

Điền vào chỗ trống:

\[
    \frac{\partial E_{total}}{\partial w_1}
    =
    \left(
    \underline{\hspace{3cm}}
    +
    \underline{\hspace{3cm}}
    \right)
    \cdot
    \frac{\partial out_{h_1}}{\partial net_{h_1}}
    \cdot
    \frac{\partial net_{h_1}}{\partial w_1}
\]

\subsection*{Đáp án}

\[
    \frac{\partial E_{total}}{\partial w_1}
    =
    \left(
    \frac{\partial E_{o_1}}{\partial out_{h_1}}
    +
    \frac{\partial E_{o_2}}{\partial out_{h_1}}
    \right)
    \cdot
    \frac{\partial out_{h_1}}{\partial net_{h_1}}
    \cdot
    \frac{\partial net_{h_1}}{\partial w_1}
\]

% ==========================================================
\section*{Kết luận}
\addcontentsline{toc}{section}{Kết luận}

Ví dụ ANN6 minh họa cách tính lan truyền ngược bằng tay trên một mạng nhỏ. Sau khi tính xuôi, ta thu được sai số tổng khoảng:

\[
    E_{total}\approx 0.2984
\]

Sau đó, ta tính ngược từ lớp output để cập nhật $w_5,w_6,w_7,w_8$, rồi tiếp tục đi về lớp trước để cập nhật $w_1,w_2,w_3,w_4$.

Điểm quan trọng nhất của ví dụ là khi một neuron ở lớp hidden nối đến nhiều output, sai số truyền về neuron đó là tổng ảnh hưởng từ tất cả các nhánh phía sau. Cụ thể, khi tính $w_1$, không được chỉ xét nhánh từ $h_1$ đến $o_1$, mà phải cộng cả nhánh từ $h_1$ đến $o_2$.

\[
    out_{h_1}
    \rightarrow
    o_1
    \rightarrow
    E_{o_1}
\]

\[
    out_{h_1}
    \rightarrow
    o_2
    \rightarrow
    E_{o_2}
\]

Vì vậy:

\[
    \frac{\partial E_{total}}{\partial out_{h_1}}
    =
    \frac{\partial E_{o_1}}{\partial out_{h_1}}
    +
    \frac{\partial E_{o_2}}{\partial out_{h_1}}
\]

Đây là phần ``ruột'' quan trọng của ví dụ, giúp hiểu thật sự điều gì xảy ra bên trong quá trình backpropagation thay vì chỉ dùng thư viện có sẵn.



\chapter{ANN7: Training Issues -- Vấn đề huấn luyện}


% ==========================================================
\section*{Learning Objectives -- Mục tiêu bài học}
\addcontentsline{toc}{section}{Mục tiêu bài học}

Bài ANN7 tập trung vào một số vấn đề thực tế khi huấn luyện mạng nơ-ron. Sau khi học xong, người học cần nắm được:

\begin{enumerate}[label=\arabic*.]
    \item Vì sao về mặt lý thuyết ta có thể dùng backpropagation và gradient descent để huấn luyện mạng, nhưng thực tế việc huấn luyện không phải lúc nào cũng dễ.
    \item Hiểu hiện tượng \textit{vanishing gradient}: gradient quá nhỏ làm trọng số cập nhật rất chậm hoặc gần như không cập nhật.
    \item Biết vì sao các mạng sâu dùng sigmoid dễ gặp vanishing gradient.
    \item Hiểu vì sao ReLU và các biến thể của nó được dùng phổ biến hơn sigmoid trong nhiều mạng sâu hiện đại.
    \item Phân biệt được \textit{underfitting}, \textit{good fitting} và \textit{overfitting}.
    \item Hiểu overfitting qua hai dạng bài toán: phân loại và hồi quy.
    \item Hiểu vai trò của nhiễu dữ liệu, sai số đo lường và khả năng tổng quát hóa trên dữ liệu mới.
    \item Phân biệt bài toán phân loại và bài toán hồi quy.
\end{enumerate}

% ==========================================================
\section{From Theory to Practice -- Từ lý thuyết huấn luyện đến vấn đề thực tế}

\subsection{Nhắc lại quá trình huấn luyện}

Trong các bài trước, ta đã học quy trình huấn luyện mạng nơ-ron:

\[
    \text{Input}
    \rightarrow
    \text{Tính xuôi}
    \rightarrow
    \text{Output dự đoán}
    \rightarrow
    \text{Tính loss}
    \rightarrow
    \text{Lan truyền ngược}
    \rightarrow
    \text{Cập nhật trọng số}
\]

Về mặt logic, nếu có dữ liệu, có hàm mất mát và có thuật toán tối ưu, ta có thể kỳ vọng rằng mô hình sẽ dần hội tụ.

Tuy nhiên, thực tế phức tạp hơn. Có những trường hợp mạng học rất chậm, không học được, hoặc học quá sát dữ liệu huấn luyện đến mức làm việc kém trên dữ liệu mới.

\begin{ghinho}
Một mô hình học tốt không chỉ là mô hình có loss nhỏ trên dữ liệu huấn luyện. Điều quan trọng hơn là mô hình phải làm việc tốt trên dữ liệu mới chưa từng thấy.
\end{ghinho}

\subsection{Ba vấn đề chính trong bài ANN7}

Bài học tập trung vào ba vấn đề quan trọng:

\begin{enumerate}
    \item \textbf{Vanishing gradient}: gradient biến mất hoặc gần bằng 0.
    \item \textbf{Underfitting}: mô hình quá đơn giản, chưa học được xu hướng dữ liệu.
    \item \textbf{Overfitting}: mô hình quá khớp với dữ liệu huấn luyện, làm mất khả năng tổng quát hóa.
\end{enumerate}

% ==========================================================
\section{Vanishing Gradient -- Vanishing Gradient}

\subsection{Vanishing gradient là gì?}

Trong backpropagation, gradient được dùng để cập nhật trọng số:

\[
    w_{\text{new}}
    =
    w_{\text{old}}
    -
    \eta
    \frac{\partial E}{\partial w}
\]

Trong đó:

\begin{itemize}
    \item $w$ là trọng số;
    \item $\eta$ là learning rate;
    \item $\frac{\partial E}{\partial w}$ là gradient của hàm mất mát theo trọng số $w$.
\end{itemize}

Nếu gradient rất nhỏ:

\[
    \frac{\partial E}{\partial w} \approx 0
\]

thì bước cập nhật:

\[
    \eta \frac{\partial E}{\partial w}
\]

cũng rất nhỏ. Khi đó trọng số gần như không thay đổi.

\begin{ghinho}
Vanishing gradient là hiện tượng gradient trở nên quá nhỏ trong quá trình lan truyền ngược, khiến các trọng số ở những lớp phía trước cập nhật rất chậm hoặc gần như không cập nhật.
\end{ghinho}

\subsection{Ví dụ số đơn giản}

Giả sử:

\[
    w = 0.5,\qquad \eta = 0.1
\]

Nếu gradient bình thường:

\[
    \frac{\partial E}{\partial w}=0.8
\]

thì:

\[
    w_{\text{new}}
    =
    0.5 - 0.1 \cdot 0.8
    =
    0.42
\]

Trọng số thay đổi đáng kể.

Nhưng nếu gradient rất nhỏ:

\[
    \frac{\partial E}{\partial w}=0.000001
\]

thì:

\[
    w_{\text{new}}
    =
    0.5 - 0.1 \cdot 0.000001
\]

\[
    w_{\text{new}}
    =
    0.4999999
\]

Trọng số gần như không thay đổi.

\begin{vidu}
Nếu một mạng có hàng triệu trọng số, và nhiều trọng số ở các lớp đầu chỉ cập nhật với lượng cực nhỏ như vậy, quá trình huấn luyện có thể trở nên rất chậm. Mạng có thể cần rất nhiều thời gian mà vẫn không học được tốt.
\end{vidu}

\subsection{Vì sao gradient có thể biến mất?}

Trong backpropagation, gradient thường là tích của nhiều đạo hàm thành phần.

Ví dụ đơn giản:

\[
    \frac{\partial E}{\partial w_1}
    =
    a_1 \cdot a_2 \cdot a_3 \cdot \ldots \cdot a_n
\]

Nếu nhiều thành phần trong đó nhỏ hơn 1, tích của chúng sẽ càng ngày càng nhỏ.

\[
    0.5^5 = 0.03125
\]

\[
    0.5^{10} = 0.0009765625
\]

\[
    0.5^{20} \approx 0.0000009537
\]

Khi mạng càng sâu, chuỗi nhân trong backpropagation càng dài. Vì vậy, gradient truyền về các lớp đầu có thể trở nên cực kỳ nhỏ.

\begin{ghinho}
Mạng càng sâu thì gradient càng phải đi qua nhiều lớp. Nếu các đạo hàm trung gian nhỏ, gradient ở các lớp đầu có thể bị ``triệt tiêu''.
\end{ghinho}

% ==========================================================
\section{Vanishing Gradient \& Sigmoid -- Vanishing Gradient và hàm sigmoid}

\subsection{Hàm sigmoid}

Hàm sigmoid được định nghĩa:

\[
    \sigma(x)=\frac{1}{1+e^{-x}}
\]

Đạo hàm của sigmoid có dạng:

\[
    \sigma'(x)=\sigma(x)(1-\sigma(x))
\]

Giá trị lớn nhất của đạo hàm sigmoid là:

\[
    \sigma'(0)=0.25
\]

Điều này có nghĩa là đạo hàm sigmoid không bao giờ lớn hơn $0.25$.

\begin{center}
\begin{tikzpicture}[scale=1.0,>=Stealth]
    % Axes
    \draw[->] (-4.5,0) -- (4.5,0) node[right] {$x$};
    \draw[->] (0,-0.1) -- (0,1.2) node[above] {$y$};

    % Sigmoid rough curve
    \draw[domain=-4:4,smooth,variable=\x,thick,blue]
        plot ({\x},{1/(1+exp(-\x))});

    \node[blue] at (2.8,0.85) {$\sigma(x)$};

    % derivative max indicator
    \draw[dashed] (0,0) -- (0,0.5);
    \filldraw[red] (0,0.5) circle (2pt);
    \node[right] at (0.1,0.5) {$\sigma(0)=0.5$};

\end{tikzpicture}
\end{center}

\subsection{Vấn đề của sigmoid trong mạng sâu}

Vì đạo hàm sigmoid thường nhỏ, khi nhân nhiều đạo hàm sigmoid qua nhiều lớp, gradient có thể giảm rất nhanh.

Ví dụ, nếu mỗi lớp tạo ra một hệ số đạo hàm khoảng $0.25$, sau 10 lớp:

\[
    0.25^{10}
    =
    0.00000095367431640625
\]

Đây là một số rất nhỏ.

Sau 20 lớp:

\[
    0.25^{20}
    \approx 9.09 \times 10^{-13}
\]

Gradient lúc này gần như biến mất.

\begin{canhbao}
Đây là một lý do quan trọng khiến sigmoid không còn được dùng phổ biến ở các lớp ẩn trong nhiều mạng sâu hiện đại.
\end{canhbao}

\subsection{Sigmoid bị bão hòa}

Khi $x$ rất lớn hoặc rất nhỏ, sigmoid gần như nằm ngang.

\[
    x \gg 0 \Rightarrow \sigma(x)\approx 1
\]

\[
    x \ll 0 \Rightarrow \sigma(x)\approx 0
\]

Ở hai vùng này, đạo hàm rất nhỏ:

\[
    \sigma'(x)\approx 0
\]

Khi đạo hàm gần 0, gradient truyền qua neuron đó cũng rất nhỏ.

\begin{vidu}
Giả sử một neuron sigmoid nhận $x=10$.

\[
    \sigma(10)\approx 0.99995
\]

Đạo hàm:

\[
    \sigma'(10)
    =
    \sigma(10)(1-\sigma(10))
\]

\[
    \approx
    0.99995(1-0.99995)
    \approx 0.00005
\]

Đạo hàm này rất nhỏ. Nếu nhiều neuron đều rơi vào vùng như vậy, gradient sẽ yếu đi rất nhanh.
\end{vidu}

% ==========================================================
\section{ReLU: Mitigating Vanishing Gradient -- ReLU và ý tưởng giảm vanishing gradient}

\subsection{Hàm ReLU}

ReLU, viết đầy đủ là \textit{Rectified Linear Unit}, được định nghĩa:

\[
    \operatorname{ReLU}(x)=\max(0,x)
\]

Tức là:

\[
    \operatorname{ReLU}(x)=
    \begin{cases}
    0, & x\leq 0\\
    x, & x>0
    \end{cases}
\]

Đạo hàm của ReLU:

\[
    \operatorname{ReLU}'(x)=
    \begin{cases}
    0, & x<0\\
    1, & x>0
    \end{cases}
\]

\begin{center}
\begin{tikzpicture}[scale=1.0,>=Stealth]
    % Axes
    \draw[->] (-4,0) -- (4,0) node[right] {$x$};
    \draw[->] (0,-0.5) -- (0,3.5) node[above] {$y$};

    % ReLU
    \draw[thick,blue] (-3.5,0) -- (0,0);
    \draw[thick,blue] (0,0) -- (3,3);
    \node[blue] at (2.5,2.4) {ReLU};

    \node[below] at (0,0) {$0$};
\end{tikzpicture}
\end{center}

\subsection{Vì sao ReLU giúp giảm vanishing gradient?}

Với $x>0$, đạo hàm của ReLU bằng 1.

\[
    \operatorname{ReLU}'(x)=1
\]

Khi lan truyền ngược qua nhiều lớp, nhân với 1 không làm gradient nhỏ đi như khi nhân với các số nhỏ hơn 1.

\begin{vidu}
So sánh hai chuỗi nhân:

\[
    0.25^{10}\approx 0.0000009537
\]

nhưng:

\[
    1^{10}=1
\]

Vì vậy, ở vùng ReLU hoạt động, gradient có thể truyền ngược tốt hơn nhiều so với sigmoid.
\end{vidu}

\subsection{Hạn chế của ReLU}

ReLU không hoàn hảo. Với $x<0$, đạo hàm của ReLU bằng 0. Khi một neuron rơi vào vùng âm và luôn cho output 0, nó có thể khó được cập nhật.

Hiện tượng này thường được gọi là \textit{dying ReLU}.

Để khắc phục, nhiều biến thể của ReLU được đề xuất, ví dụ:

\begin{itemize}
    \item Leaky ReLU;
    \item Parametric ReLU;
    \item ELU;
    \item SELU.
\end{itemize}

\begin{luuy}
Ý chính trong bài học là: các hàm kích hoạt hiện đại được phát triển một phần để giảm vấn đề vanishing gradient và giúp mạng sâu huấn luyện hiệu quả hơn.
\end{luuy}

% ==========================================================
\section{Underfitting và Overfitting}

\subsection{Bài toán fitting là gì?}

Khi huấn luyện mô hình, ta muốn mô hình học được quan hệ giữa input và output từ dữ liệu huấn luyện.

Nếu mô hình học quá ít, nó không nắm được xu hướng dữ liệu. Đây là underfitting.

Nếu mô hình học quá sát từng điểm dữ liệu huấn luyện, kể cả nhiễu và điểm bất thường, nó có thể làm việc kém trên dữ liệu mới. Đây là overfitting.

Trường hợp mong muốn là mô hình học được xu hướng chung của dữ liệu.

\begin{ghinho}
Mục tiêu không phải là nhớ chính xác mọi điểm dữ liệu huấn luyện, mà là học được quy luật tổng quát để áp dụng cho dữ liệu mới.
\end{ghinho}

\subsection{Ba trạng thái thường gặp}

\begin{center}
\begin{tabular}{p{0.25\textwidth}p{0.35\textwidth}p{0.30\textwidth}}
\toprule
Trạng thái & Mô tả & Hậu quả \\
\midrule
Underfitting & Mô hình quá đơn giản, chưa khớp dữ liệu & Sai nhiều trên cả train và test \\
Good fitting & Mô hình học được xu hướng chung & Làm việc tốt trên dữ liệu mới \\
Overfitting & Mô hình quá khớp dữ liệu huấn luyện & Train tốt nhưng test kém \\
\bottomrule
\end{tabular}
\end{center}

% ==========================================================
\section{Ví dụ phân loại 2D}

\subsection{Mô tả ví dụ}

Giả sử ta có dữ liệu 2D, ví dụ:

\[
    x_1 = \text{chiều cao}, \qquad x_2 = \text{cân nặng}
\]

Mỗi điểm dữ liệu thuộc một trong hai lớp:

\[
    \text{Lớp A} \quad \text{hoặc} \quad \text{Lớp B}
\]

Ví dụ:

\begin{itemize}
    \item người thuộc nhóm A và nhóm B;
    \item người thừa cân và không thừa cân;
    \item người mắc bệnh và không mắc bệnh.
\end{itemize}

Mô hình cần tìm đường ranh giới để phân chia hai lớp.

\subsection{Underfitting trong phân loại}

Underfitting xảy ra khi đường phân lớp quá đơn giản so với dữ liệu.

Ví dụ, dữ liệu có xu hướng cong, nhưng mô hình chỉ dùng một đường thẳng đơn giản. Khi đó, nhiều điểm bị phân loại sai.

\begin{center}
\begin{tikzpicture}[scale=0.9,>=Stealth]
    \draw[->] (0,0) -- (5,0) node[right] {$x_1$};
    \draw[->] (0,0) -- (0,4) node[above] {$x_2$};

    % Class A circles
    \foreach \x/\y in {0.8/0.8,1.1/1.2,1.4/1.5,1.8/1.8,2.2/2.0,2.6/2.2,3.0/2.4}
        \draw[blue,thick] (\x,\y) circle (2pt);

    % Class B crosses
    \foreach \x/\y in {1.2/3.0,1.7/3.2,2.2/3.0,2.7/2.9,3.2/2.7,3.7/2.5,4.0/2.3}
        \draw[red,thick] (\x-0.06,\y-0.06) -- (\x+0.06,\y+0.06)
                          (\x-0.06,\y+0.06) -- (\x+0.06,\y-0.06);

    % Bad linear boundary
    \draw[thick,orange] (0.6,2.2) -- (4.5,1.6);
    \node[orange] at (3.3,0.8) {underfitting};

\end{tikzpicture}
\end{center}

\begin{vidu}
Nếu dữ liệu thật có ranh giới cong, nhưng ta dùng mô hình tuyến tính quá đơn giản, mô hình sẽ không học được cấu trúc dữ liệu. Kết quả là nhiều điểm ở cả hai lớp bị phân loại sai.
\end{vidu}

\subsection{ fitting trong phân loại}

Good fitting là khi mô hình học được xu hướng chung của dữ liệu. Nó có thể chấp nhận một vài điểm sai do nhiễu hoặc điểm bất thường, nhưng đường phân lớp vẫn hợp lý.

\begin{center}
\begin{tikzpicture}[scale=0.9,>=Stealth]
    \draw[->] (0,0) -- (5,0) node[right] {$x_1$};
    \draw[->] (0,0) -- (0,4) node[above] {$x_2$};

    % Class A circles
    \foreach \x/\y in {0.8/0.8,1.1/1.2,1.4/1.5,1.8/1.8,2.2/2.0,2.6/2.2,3.0/2.4,3.4/2.55}
        \draw[blue,thick] (\x,\y) circle (2pt);

    % Class B crosses
    \foreach \x/\y in {1.2/3.0,1.7/3.2,2.2/3.0,2.7/2.9,3.2/2.7,3.7/2.5,4.0/2.3}
        \draw[red,thick] (\x-0.06,\y-0.06) -- (\x+0.06,\y+0.06)
                          (\x-0.06,\y+0.06) -- (\x+0.06,\y-0.06);

    % A noisy point
    \draw[blue,thick] (3.6,2.85) circle (2pt);

    % Good curved boundary
    \draw[thick,orange,smooth] plot coordinates {
        (0.7,2.2) (1.3,2.35) (2.0,2.45) (2.8,2.55) (3.6,2.70) (4.3,2.95)
    };
    \node[orange] at (3.5,1.0) {good fitting};

\end{tikzpicture}
\end{center}

\begin{ghinho}
Một mô hình tốt không nhất thiết phải phân loại đúng 100\% dữ liệu huấn luyện. Nó cần học được xu hướng chung để dự đoán tốt trên dữ liệu mới.
\end{ghinho}

\subsection{Overfitting trong phân loại}

Overfitting xảy ra khi mô hình cố gắng phân loại đúng hoàn hảo mọi điểm trong tập huấn luyện, kể cả điểm nhiễu hoặc điểm bất thường.

Đường phân lớp khi đó có thể trở nên rất ngoằn ngoèo, phức tạp và phi tự nhiên.

\begin{center}
\begin{tikzpicture}[scale=0.9,>=Stealth]
    \draw[->] (0,0) -- (5,0) node[right] {$x_1$};
    \draw[->] (0,0) -- (0,4) node[above] {$x_2$};

    % Class A circles
    \foreach \x/\y in {0.8/0.8,1.1/1.2,1.4/1.5,1.8/1.8,2.2/2.0,2.6/2.2,3.0/2.4,3.6/2.85}
        \draw[blue,thick] (\x,\y) circle (2pt);

    % Class B crosses
    \foreach \x/\y in {1.2/3.0,1.7/3.2,2.2/3.0,2.7/2.9,3.2/2.7,3.7/2.5,4.0/2.3}
        \draw[red,thick] (\x-0.06,\y-0.06) -- (\x+0.06,\y+0.06)
                          (\x-0.06,\y+0.06) -- (\x+0.06,\y-0.06);

    % Overfit boundary
    \draw[thick,orange,smooth] plot coordinates {
        (0.6,2.4) (1.1,2.45) (1.6,2.55) (2.1,2.50)
        (2.6,2.65) (3.1,2.55) (3.45,3.1) (3.8,2.75) (4.4,2.9)
    };
    \node[orange] at (3.3,1.0) {overfitting};

\end{tikzpicture}
\end{center}

\begin{canhbao}
Nếu mô hình quá khớp với dữ liệu huấn luyện, nó có thể đạt sai số gần 0 trên tập huấn luyện nhưng lại dự đoán sai trên dữ liệu thực tế.
\end{canhbao}

\subsection{Vì sao phân loại hoàn hảo chưa chắc là tốt?}

Giả sử mô hình phân loại hoàn hảo mọi điểm huấn luyện. Điều này nghe có vẻ tốt, nhưng chưa chắc đúng.

Lý do là bài toán thực tế không chỉ yêu cầu xử lý các điểm đã có. Ta muốn dùng mô hình cho dữ liệu mới trong tương lai.

Nếu đường phân lớp quá phức tạp chỉ để ôm sát dữ liệu huấn luyện, nó có thể đánh mất xu hướng chung.

\begin{vidu}
Giả sử mô hình dùng để phân loại người mắc bệnh và không mắc bệnh. Nếu trong dữ liệu huấn luyện có một điểm nhiễu, mô hình overfit có thể bẻ cong đường phân lớp chỉ để phân loại đúng điểm đó.

Khi gặp một người mới có đặc trưng gần khu vực đó, mô hình có thể dự đoán sai nghiêm trọng. Ví dụ, người không thuộc nhóm nguy cơ lại bị dự đoán là có bệnh.
\end{vidu}

% ==========================================================
\section{Nhiễu dữ liệu và sai số đo lường}

\subsection{Dữ liệu đo được không hoàn hảo}

Trong thực tế, dữ liệu không bao giờ hoàn hảo tuyệt đối. Dữ liệu có thể bị ảnh hưởng bởi:

\begin{itemize}
    \item sai số cảm biến;
    \item điều kiện môi trường;
    \item cách đo;
    \item sai số do con người nhập liệu;
    \item nhiễu ngẫu nhiên;
    \item điểm bất thường.
\end{itemize}

\subsection{Ví dụ đo chiều dài}

Giả sử ta cần đo chiều dài một viên gạch bằng thước hoặc cảm biến laser.

Con số đo được không chắc là chiều dài thật tuyệt đối của vật thể. Nó chỉ là giá trị đo lường.

Giá trị đo có thể bị ảnh hưởng bởi:

\begin{itemize}
    \item độ chính xác của thiết bị;
    \item góc đặt cảm biến;
    \item bề mặt phản xạ;
    \item nhiệt độ và độ ẩm;
    \item lỗi ghi chép hoặc nhập dữ liệu.
\end{itemize}

\begin{ghinho}
Dữ liệu huấn luyện thường chứa nhiễu. Vì vậy, mô hình tốt cần học xu hướng chung thay vì cố gắng khớp hoàn hảo từng điểm dữ liệu.
\end{ghinho}

\subsection{Liên hệ với overfitting}

Nếu mô hình quá phức tạp, nó có thể học luôn cả nhiễu trong dữ liệu.

Khi đó, mô hình không chỉ học quy luật thật mà còn học cả sai số đo, điểm bất thường và lỗi nhập liệu.

Đây là nguyên nhân quan trọng dẫn đến overfitting.

% ==========================================================
\section{Ví dụ hồi quy}

\subsection{Bài toán hồi quy là gì?}

Hồi quy, tiếng Anh là \textit{regression}, là bài toán dự đoán một giá trị liên tục.

Ví dụ:

\begin{itemize}
    \item dự đoán giá cổ phiếu ngày mai là bao nhiêu;
    \item dự đoán nhiệt độ ngày mai;
    \item dự đoán sản lượng điện;
    \item dự đoán giá nhà;
    \item dự đoán chi phí sản xuất.
\end{itemize}

Khác với phân loại, output của hồi quy không phải là một lớp rời rạc, mà là một con số liên tục.

\subsection{Dữ liệu hồi quy có nhiễu}

Giả sử dữ liệu thật đi theo một đường cong nào đó. Nhưng do nhiễu, các điểm dữ liệu đo được không nằm chính xác trên đường cong thật.

Mô hình cần tìm một đường xấp xỉ phù hợp với xu hướng chung của dữ liệu.

\subsection{Underfitting trong hồi quy}

Underfitting xảy ra khi mô hình quá đơn giản.

Ví dụ, dữ liệu có dạng cong hoặc dạng sóng, nhưng ta dùng một đường thẳng để xấp xỉ.

\begin{center}
\begin{tikzpicture}[scale=0.95,>=Stealth]
    \draw[->] (0,0) -- (6,0) node[right] {$x$};
    \draw[->] (0,0) -- (0,4) node[above] {$y$};

    % Data points
    \foreach \x/\y in {0.5/1.0,1/1.5,1.5/2.2,2/2.6,2.5/2.4,3/2.0,3.5/1.7,4/2.0,4.5/2.7,5/3.2,5.5/3.1}
        \filldraw[blue] (\x,\y) circle (2pt);

    % Underfit line
    \draw[thick,orange] (0.4,1.5) -- (5.6,2.4);
    \node[orange] at (4.2,0.8) {underfitting};

\end{tikzpicture}
\end{center}

\begin{vidu}
Nếu dữ liệu nhiệt độ theo thời gian có xu hướng tăng giảm theo mùa, nhưng mô hình chỉ dùng một đường thẳng, mô hình có thể không bắt được dao động thật của dữ liệu.
\end{vidu}

\subsection{Good fitting trong hồi quy}

Good fitting là khi mô hình xấp xỉ được xu hướng chính của dữ liệu, dù không đi qua mọi điểm.

\begin{center}
\begin{tikzpicture}[scale=0.95,>=Stealth]
    \draw[->] (0,0) -- (6,0) node[right] {$x$};
    \draw[->] (0,0) -- (0,4) node[above] {$y$};

    % Data points
    \foreach \x/\y in {0.5/1.0,1/1.5,1.5/2.2,2/2.6,2.5/2.4,3/2.0,3.5/1.7,4/2.0,4.5/2.7,5/3.2,5.5/3.1}
        \filldraw[blue] (\x,\y) circle (2pt);

    % Good curve
    \draw[thick,orange,smooth] plot coordinates {
        (0.4,0.9) (1.0,1.5) (1.7,2.25) (2.4,2.45)
        (3.0,2.1) (3.6,1.8) (4.2,2.2) (4.8,2.85) (5.6,3.2)
    };
    \node[orange] at (4.2,0.8) {good fitting};

\end{tikzpicture}
\end{center}

\begin{ghinho}
Trong hồi quy, mô hình tốt không nhất thiết đi qua tất cả điểm dữ liệu. Nó cần mô tả được xu hướng tổng quát phía sau dữ liệu nhiễu.
\end{ghinho}

\subsection{Overfitting trong hồi quy}

Overfitting xảy ra khi mô hình quá phức tạp và cố gắng đi qua gần như tất cả các điểm dữ liệu.

Ví dụ, dùng đa thức bậc rất cao để fit một tập điểm nhiễu.

\begin{center}
\begin{tikzpicture}[scale=0.95,>=Stealth]
    \draw[->] (0,0) -- (6,0) node[right] {$x$};
    \draw[->] (0,0) -- (0,4) node[above] {$y$};

    % Data points
    \foreach \x/\y in {0.5/1.0,1/1.5,1.5/2.2,2/2.6,2.5/2.4,3/2.0,3.5/1.7,4/2.0,4.5/2.7,5/3.2,5.5/3.1}
        \filldraw[blue] (\x,\y) circle (2pt);

    % Overfit curve
    \draw[thick,orange,smooth] plot coordinates {
        (0.5,1.0) (0.8,1.9) (1.0,1.5) (1.25,1.1)
        (1.5,2.2) (1.8,3.0) (2.0,2.6) (2.3,1.6)
        (2.5,2.4) (2.8,2.9) (3.0,2.0) (3.3,1.2)
        (3.5,1.7) (3.8,2.7) (4.0,2.0) (4.3,1.5)
        (4.5,2.7) (4.8,3.6) (5.0,3.2) (5.3,2.5) (5.5,3.1)
    };
    \node[orange] at (4.2,0.8) {overfitting};

\end{tikzpicture}
\end{center}

\begin{canhbao}
Một đường cong đi qua tất cả điểm huấn luyện có thể có sai số huấn luyện bằng 0, nhưng vẫn là mô hình xấu nếu nó không phản ánh xu hướng thật của dữ liệu.
\end{canhbao}

\subsection{Ví dụ đa thức bậc cao}

Giả sử dữ liệu thật có thể được mô tả tốt bằng một hàm bậc 4. Nhưng nếu ta dùng đa thức bậc 15, mô hình có thể uốn lượn rất mạnh để đi qua từng điểm dữ liệu.

Khi đó:

\[
    \text{Training error} \approx 0
\]

nhưng:

\[
    \text{Test error} \text{ có thể rất lớn}
\]

Vì mô hình đã học cả nhiễu, không chỉ học quy luật.

% ==========================================================
\section{Phân loại và hồi quy}

\subsection{Bài toán phân loại}

Phân loại, tiếng Anh là \textit{classification}, là bài toán dự đoán một nhãn rời rạc.

Ví dụ:

\begin{itemize}
    \item ảnh là chó hay mèo;
    \item email là spam hay không spam;
    \item người có mắc COVID hay không;
    \item giá cổ phiếu ngày mai tăng, giảm hay không đổi;
    \item vật thể là xe hơi, máy bay, người hay con vật.
\end{itemize}

Output thường là một class:

\[
    y \in \{C_1,C_2,\ldots,C_K\}
\]

\subsection{Bài toán hồi quy}

Hồi quy là bài toán dự đoán giá trị liên tục.

Ví dụ:

\begin{itemize}
    \item giá cổ phiếu ngày mai bằng bao nhiêu;
    \item giá nhà là bao nhiêu;
    \item sản lượng điện tháng tới là bao nhiêu;
    \item nhiệt độ ngày mai là bao nhiêu.
\end{itemize}

Output thường là một số thực:

\[
    y \in \mathbb{R}
\]

\subsection{So sánh classification và regression}

\begin{center}
\begin{tabular}{p{0.26\textwidth}p{0.32\textwidth}p{0.32\textwidth}}
\toprule
Tiêu chí & Classification & Regression \\
\midrule
Loại output & Nhãn rời rạc & Giá trị liên tục \\
Ví dụ output & Chó, mèo, xe hơi & 25.3 độ C, 1.2 tỷ đồng \\
Ví dụ bài toán & Email spam hay không & Giá nhà bao nhiêu \\
Sai số thường dùng & Cross entropy & MSE, MAE \\
\bottomrule
\end{tabular}
\end{center}

\begin{vidu}
Cùng là giá cổ phiếu, có thể có hai bài toán khác nhau:

\begin{itemize}
    \item Dự đoán ngày mai cổ phiếu tăng hay giảm: classification.
    \item Dự đoán ngày mai cổ phiếu có giá chính xác bao nhiêu: regression.
\end{itemize}
\end{vidu}

% ==========================================================
\section{Tổng hợp: Khi nào mô hình học tốt?}

\subsection{Mô hình học tốt cần cân bằng}

Một mô hình học tốt cần cân bằng giữa hai cực:

\[
    \text{quá đơn giản}
    \quad \longleftrightarrow \quad
    \text{quá phức tạp}
\]

Nếu quá đơn giản, mô hình underfit.

Nếu quá phức tạp, mô hình dễ overfit.

Trường hợp tốt là mô hình đủ linh hoạt để học xu hướng dữ liệu, nhưng không quá phức tạp đến mức học cả nhiễu.

\subsection{Dấu hiệu nhận biết}

\begin{center}
\begin{tabular}{p{0.25\textwidth}p{0.30\textwidth}p{0.35\textwidth}}
\toprule
Trường hợp & Dữ liệu huấn luyện & Dữ liệu kiểm tra \\
\midrule
Underfitting & Sai số cao & Sai số cao \\
Good fitting & Sai số thấp vừa phải & Sai số thấp \\
Overfitting & Sai số rất thấp & Sai số cao \\
\bottomrule
\end{tabular}
\end{center}

\subsection{Liên hệ với loss bằng 0}

Trong lý thuyết, ta thường muốn loss tiến về 0. Nhưng trong thực tế, loss huấn luyện bằng 0 không phải lúc nào cũng là mục tiêu tốt.

Nếu dữ liệu có nhiễu, việc cố gắng đưa loss huấn luyện về 0 có thể làm mô hình học luôn cả nhiễu.

\begin{ghinho}
Loss thấp trên training set là cần thiết, nhưng chưa đủ. Mô hình phải có khả năng tổng quát hóa tốt trên dữ liệu mới.
\end{ghinho}

% ==========================================================
\section{Technical Glossary -- Bảng thuật ngữ chuyên ngành}

\small
\begin{longtable}{p{0.26\textwidth}p{0.48\textwidth}p{0.19\textwidth}}
\toprule
\textbf{Thuật ngữ} & \textbf{Giải thích} & \textbf{Ví dụ} \\
\midrule
\endfirsthead

\toprule
\textbf{Thuật ngữ} & \textbf{Giải thích} & \textbf{Ví dụ} \\
\midrule
\endhead

Vanishing gradient & Gradient trở nên rất nhỏ khi lan truyền ngược, khiến trọng số cập nhật rất chậm. & Gradient $\approx 10^{-9}$. \\

Gradient & Đạo hàm của loss theo tham số, dùng để cập nhật trọng số. & $\frac{\partial E}{\partial w}$. \\

Backpropagation & Quá trình lan truyền ngược sai số để tính gradient. & Tính gradient từ output về input. \\

Sigmoid & Hàm kích hoạt có output từ 0 đến 1. & $\sigma(x)=\frac{1}{1+e^{-x}}$. \\

ReLU & Hàm kích hoạt bằng 0 nếu $x\leq 0$, bằng $x$ nếu $x>0$. & $\max(0,x)$. \\

Dying ReLU & Hiện tượng neuron ReLU rơi vào vùng âm và gần như không hoạt động. & Output luôn bằng 0. \\

Underfitting & Mô hình quá đơn giản, chưa học được xu hướng dữ liệu. & Dùng đường thẳng cho dữ liệu cong. \\

Overfitting & Mô hình quá khớp dữ liệu huấn luyện, học cả nhiễu. & Đường cong đi qua mọi điểm train. \\

Good fitting & Mô hình học được xu hướng chung và làm việc tốt trên dữ liệu mới. & Đường cong xấp xỉ hợp lý. \\

Noise & Nhiễu trong dữ liệu do đo lường, môi trường hoặc lỗi nhập liệu. & Cảm biến laser sai số. \\

Generalization & Khả năng tổng quát hóa, tức làm việc tốt trên dữ liệu mới. & Test error thấp. \\

Training data & Dữ liệu dùng để huấn luyện mô hình. & 500 điểm dữ liệu đã có. \\

Test data & Dữ liệu dùng để kiểm tra khả năng tổng quát hóa. & Dữ liệu mới chưa học. \\

Classification & Bài toán phân loại vào các class rời rạc. & Chó hoặc mèo. \\

Regression & Bài toán dự đoán giá trị liên tục. & Dự đoán giá nhà. \\

Decision boundary & Đường hoặc mặt phân chia giữa các class. & Đường đỏ phân lớp. \\

MSE & Mean Squared Error, sai số bình phương trung bình. & Dùng trong hồi quy. \\

Cross entropy & Hàm mất mát phổ biến cho phân loại. & Phân loại nhiều lớp. \\

\bottomrule
\end{longtable}
\normalsize

% ==========================================================
\section{Key Concepts Summary -- Tóm tắt kiến thức trọng tâm}

\subsection{Các ý chính cần nhớ}

\begin{enumerate}
    \item Backpropagation về lý thuyết giúp tính gradient và cập nhật trọng số, nhưng thực tế huấn luyện có thể gặp nhiều vấn đề.
    \item Vanishing gradient xảy ra khi gradient quá nhỏ, làm trọng số cập nhật rất chậm.
    \item Mạng càng sâu thì gradient càng dễ bị nhỏ đi nếu phải nhân nhiều đạo hàm nhỏ.
    \item Sigmoid dễ gây vanishing gradient vì đạo hàm của nó nhỏ, đặc biệt ở vùng bão hòa.
    \item ReLU giúp giảm vanishing gradient ở vùng $x>0$ vì đạo hàm bằng 1.
    \item Underfitting là mô hình chưa học được xu hướng dữ liệu.
    \item Overfitting là mô hình học quá sát dữ liệu huấn luyện, kể cả nhiễu.
    \item Mô hình tốt là mô hình học được xu hướng chung và làm việc tốt trên dữ liệu mới.
    \item Dữ liệu thực tế luôn có nhiễu, sai số đo lường và điểm bất thường.
    \item Classification dự đoán nhãn rời rạc; regression dự đoán giá trị liên tục.
\end{enumerate}

\subsection{Công thức quan trọng}

\subsection*{Cập nhật trọng số}

\[
    w_{\text{new}}
    =
    w_{\text{old}}
    -
    \eta
    \frac{\partial E}{\partial w}
\]

\subsection*{Sigmoid}

\[
    \sigma(x)=\frac{1}{1+e^{-x}}
\]

\subsection*{Đạo hàm sigmoid}

\[
    \sigma'(x)=\sigma(x)(1-\sigma(x))
\]

\subsection*{ReLU}

\[
    \operatorname{ReLU}(x)=\max(0,x)
\]

\subsection*{Đạo hàm ReLU}

\[
    \operatorname{ReLU}'(x)=
    \begin{cases}
    0, & x<0\\
    1, & x>0
    \end{cases}
\]

% ==========================================================
\section{Review Questions -- Câu hỏi ôn tập}

\begin{enumerate}
    \item Vanishing gradient là gì?
    \item Khi gradient gần bằng 0, trọng số được cập nhật như thế nào?
    \item Vì sao mạng sâu dễ gặp vanishing gradient hơn mạng nông?
    \item Vì sao sigmoid có thể gây vanishing gradient?
    \item Đạo hàm lớn nhất của sigmoid là bao nhiêu?
    \item ReLU có công thức như thế nào?
    \item Vì sao ReLU giúp gradient truyền tốt hơn ở vùng $x>0$?
    \item Underfitting là gì?
    \item Overfitting là gì?
    \item Vì sao mô hình có training loss bằng 0 chưa chắc là mô hình tốt?
    \item Vì sao dữ liệu thực tế thường có nhiễu?
    \item Trong bài toán phân loại 2D, đường phân lớp tốt có nhất thiết phân loại đúng mọi điểm huấn luyện không?
    \item Trong hồi quy, vì sao đường cong đi qua mọi điểm dữ liệu có thể là mô hình xấu?
    \item Classification khác regression như thế nào?
    \item Cho ví dụ một bài toán classification và một bài toán regression.
\end{enumerate}

% ==========================================================
\section{Practice Exercises -- Bài tập vận dụng}

\subsection*{Bài tập 1: Kiểm tra vanishing gradient}
\addcontentsline{toc}{section}{Bài tập 1: Kiểm tra vanishing gradient}

Giả sử gradient qua mỗi lớp bị nhân với hệ số $0.25$. Tính giá trị gradient còn lại sau 5 lớp và sau 10 lớp nếu gradient ban đầu là 1.

\subsection*{Lời giải gợi ý}

Sau 5 lớp:

\[
    1\cdot 0.25^5
    =
    0.0009765625
\]

Sau 10 lớp:

\[
    1\cdot 0.25^{10}
    \approx
    0.0000009537
\]

Gradient giảm rất nhanh khi số lớp tăng.

\subsection*{Bài tập 2: Tính đạo hàm sigmoid}
\addcontentsline{toc}{section}{Bài tập 2: Tính đạo hàm sigmoid}

Cho:

\[
    \sigma(x)=0.9
\]

Tính:

\[
    \sigma'(x)=\sigma(x)(1-\sigma(x))
\]

\subsection*{Lời giải gợi ý}

\[
    \sigma'(x)=0.9(1-0.9)
\]

\[
    \sigma'(x)=0.9\cdot 0.1=0.09
\]

\subsection*{Bài tập 3: Nhận diện underfitting và overfitting}
\addcontentsline{toc}{section}{Bài tập 3: Nhận diện underfitting và overfitting}

Hãy xác định trường hợp nào là underfitting, good fitting hoặc overfitting:

\begin{enumerate}
    \item Mô hình sai nhiều trên cả dữ liệu huấn luyện và dữ liệu kiểm tra.
    \item Mô hình đúng gần như tuyệt đối trên dữ liệu huấn luyện nhưng sai nhiều trên dữ liệu kiểm tra.
    \item Mô hình có sai số huấn luyện vừa phải và sai số kiểm tra thấp.
\end{enumerate}

\subsection*{Lời giải gợi ý}

\begin{enumerate}
    \item Underfitting.
    \item Overfitting.
    \item Good fitting.
\end{enumerate}

\subsection*{Bài tập 4: Phân biệt classification và regression}
\addcontentsline{toc}{section}{Bài tập 4: Phân biệt classification và regression}

Xác định mỗi bài toán sau là classification hay regression:

\begin{enumerate}
    \item Dự đoán giá nhà.
    \item Phân loại email spam hoặc không spam.
    \item Dự đoán sản lượng điện tháng tới.
    \item Nhận dạng ảnh là chó, mèo hay xe hơi.
    \item Dự đoán cổ phiếu ngày mai tăng, giảm hay không đổi.
    \item Dự đoán giá cổ phiếu ngày mai bằng bao nhiêu.
\end{enumerate}

\subsection*{Lời giải gợi ý}

\begin{center}
\begin{tabular}{ll}
\toprule
Bài toán & Loại bài toán \\
\midrule
Dự đoán giá nhà & Regression \\
Email spam hoặc không spam & Classification \\
Dự đoán sản lượng điện & Regression \\
Nhận dạng ảnh chó, mèo, xe hơi & Classification \\
Cổ phiếu tăng, giảm hay không đổi & Classification \\
Giá cổ phiếu bằng bao nhiêu & Regression \\
\bottomrule
\end{tabular}
\end{center}

% ==========================================================
\section*{Kết luận}
\addcontentsline{toc}{section}{Kết luận}

ANN7 nhấn mạnh rằng huấn luyện mạng nơ-ron không chỉ là áp dụng backpropagation và gradient descent một cách máy móc. Trong thực tế, mạng có thể gặp các vấn đề như vanishing gradient, underfitting và overfitting.

Vanishing gradient làm cho các trọng số ở những lớp đầu cập nhật rất chậm, đặc biệt trong các mạng sâu dùng sigmoid. Đây là một trong những lý do các hàm kích hoạt như ReLU và các biến thể của ReLU trở nên phổ biến.

Underfitting xảy ra khi mô hình quá đơn giản, chưa học được xu hướng dữ liệu. Overfitting xảy ra khi mô hình quá phức tạp, học quá sát dữ liệu huấn luyện và học cả nhiễu. Mô hình tốt là mô hình không nhất thiết đúng hoàn hảo trên dữ liệu huấn luyện, nhưng phải học được xu hướng chung và tổng quát hóa tốt trên dữ liệu mới.



\chapter{ANN8: Overfitting \& Underfitting -- Quá khớp và dưới khớp}


% ==========================================================
\section*{Learning Objectives -- Mục tiêu bài học}
\addcontentsline{toc}{section}{Mục tiêu bài học}

Bài ANN8 tiếp tục nội dung về underfitting và overfitting, nhưng tập trung hơn vào nguyên nhân và các kỹ thuật xử lý. Sau khi học xong, người học cần nắm được:

\begin{enumerate}[label=\arabic*.]
    \item Nhắc lại ý nghĩa của underfitting và overfitting trong bài toán phân loại và hồi quy.
    \item Hiểu nguyên nhân gây underfitting: mô hình quá đơn giản hoặc huấn luyện chưa đủ.
    \item Hiểu nguyên nhân gây overfitting: mô hình quá phức tạp, dữ liệu ít hoặc dữ liệu thiếu đa dạng.
    \item Biết một số hướng xử lý underfitting.
    \item Biết một số hướng xử lý overfitting: đơn giản hóa mô hình, bỏ bớt đặc trưng, tăng dữ liệu, data augmentation, early stopping và dropout.
    \item Hiểu ý nghĩa của train/validation/test set và nguyên tắc sử dụng từng tập.
    \item Hiểu vì sao dùng sai tập test có thể làm sai lệch kết quả nghiên cứu.
    \item Nắm được trực giác của dropout: không cho mạng phụ thuộc quá mạnh vào một số neuron cụ thể.
\end{enumerate}

% ==========================================================
\section{Underfitting \& Overfitting Recap -- Ôn lại Underfitting và Overfitting}

\subsection{Underfitting}

Underfitting là tình huống mô hình chưa khớp được với dữ liệu. Nói cách khác, mô hình chưa học được xu hướng tổng quát của dữ liệu.

Thường gặp khi:

\begin{itemize}
    \item mô hình quá đơn giản;
    \item số tham số quá ít;
    \item kiến trúc mạng quá yếu;
    \item thời gian huấn luyện chưa đủ;
    \item đặc trưng đầu vào chưa đủ thông tin.
\end{itemize}

\begin{ghinho}
Underfitting nghĩa là mô hình chưa đủ khả năng học bài toán. Dù có tiếp tục dùng mô hình quá đơn giản, sai số vẫn có thể không giảm đến mức mong muốn.
\end{ghinho}

\subsection{Overfitting}

Overfitting là tình huống mô hình quá khớp với dữ liệu huấn luyện. Khi đó, sai số trên tập huấn luyện có thể rất thấp, thậm chí gần bằng 0, nhưng mô hình lại làm việc kém trên dữ liệu mới.

Thường gặp khi:

\begin{itemize}
    \item mô hình quá phức tạp so với bài toán;
    \item dữ liệu huấn luyện quá ít;
    \item dữ liệu thiếu đa dạng;
    \item mô hình học cả nhiễu và chi tiết bất thường;
    \item quá trình huấn luyện kéo dài quá lâu.
\end{itemize}

\begin{canhbao}
Overfitting nguy hiểm hơn underfitting ở chỗ: trên tập huấn luyện, mô hình có vẻ rất tốt vì loss thấp. Nếu không kiểm tra đúng cách, ta có thể nhầm rằng mô hình đã làm việc tốt.
\end{canhbao}

\subsection{So sánh ngắn gọn}

\begin{center}
\begin{tabular}{p{0.24\textwidth}p{0.34\textwidth}p{0.32\textwidth}}
\toprule
Trường hợp & Biểu hiện & Nguyên nhân thường gặp \\
\midrule
Underfitting & Sai nhiều trên cả train và test & Mô hình quá đơn giản hoặc train chưa đủ \\
Good fitting & Học được xu hướng chung, test tốt & Mô hình phù hợp với độ phức tạp bài toán \\
Overfitting & Train rất tốt, test kém & Mô hình quá phức tạp hoặc dữ liệu ít/thiếu đa dạng \\
\bottomrule
\end{tabular}
\end{center}

% ==========================================================
\section{Underfitting \& Overfitting: Classification -- Ví dụ Underfitting và Overfitting trong phân loại}

\subsection{Mô tả bài toán phân loại}

Giả sử ta có dữ liệu 2D, ví dụ:

\[
    x_1 = \text{chiều cao}, \qquad x_2 = \text{cân nặng}
\]

Mỗi điểm dữ liệu thuộc một trong hai class:

\[
    C_1 \quad \text{hoặc} \quad C_2
\]

Mục tiêu là tìm một đường ranh giới để phân biệt hai class.

\subsection{Underfitting trong phân loại}

Nếu ranh giới thật cần là một đường cong, nhưng ta chỉ dùng một đường thẳng, mô hình có thể không phân chia tốt hai class.

\begin{center}
\begin{tikzpicture}[scale=0.9,>=Stealth]
    \draw[->] (0,0) -- (5,0) node[right] {$x_1$};
    \draw[->] (0,0) -- (0,4) node[above] {$x_2$};

    % Class 1
    \foreach \x/\y in {0.8/0.7,1.2/1.0,1.5/1.4,2.0/1.8,2.5/2.1,3.0/2.3}
        \draw[blue,thick] (\x,\y) circle (2pt);

    % Class 2
    \foreach \x/\y in {1.0/3.1,1.5/3.2,2.1/3.0,2.7/2.85,3.3/2.65,3.9/2.45}
        \draw[red,thick] (\x-0.06,\y-0.06) -- (\x+0.06,\y+0.06)
                          (\x-0.06,\y+0.06) -- (\x+0.06,\y-0.06);

    % Underfit boundary
    \draw[thick,orange] (0.5,2.0) -- (4.5,1.75);
    \node[orange] at (3.6,0.8) {Đường thẳng quá đơn giản};

\end{tikzpicture}
\end{center}

\begin{vidu}
Nếu dữ liệu có xu hướng cong nhưng mô hình chỉ có khả năng tạo đường thẳng, ta xoay hoặc dịch đường thẳng thế nào cũng khó phân loại tốt. Đây là ví dụ điển hình của underfitting.
\end{vidu}

\subsection{Overfitting trong phân loại}

Ngược lại, nếu mô hình quá mạnh, nó có thể tạo ra một đường ranh giới rất ngoằn ngoèo để phân loại đúng mọi điểm huấn luyện, kể cả các điểm nhiễu.

\begin{center}
\begin{tikzpicture}[scale=0.9,>=Stealth]
    \draw[->] (0,0) -- (5,0) node[right] {$x_1$};
    \draw[->] (0,0) -- (0,4) node[above] {$x_2$};

    % Class 1
    \foreach \x/\y in {0.8/0.7,1.2/1.0,1.5/1.4,2.0/1.8,2.5/2.1,3.0/2.3,3.5/2.85}
        \draw[blue,thick] (\x,\y) circle (2pt);

    % Class 2
    \foreach \x/\y in {1.0/3.1,1.5/3.2,2.1/3.0,2.7/2.85,3.3/2.65,3.9/2.45}
        \draw[red,thick] (\x-0.06,\y-0.06) -- (\x+0.06,\y+0.06)
                          (\x-0.06,\y+0.06) -- (\x+0.06,\y-0.06);

    % Overfit boundary
    \draw[thick,orange,smooth] plot coordinates {
        (0.6,2.25) (1.1,2.35) (1.6,2.50) (2.1,2.43)
        (2.6,2.62) (3.05,2.48) (3.35,3.12) (3.75,2.68) (4.3,2.85)
    };
    \node[orange] at (3.5,0.8) {Ranh giới quá phức tạp};

\end{tikzpicture}
\end{center}

\begin{ghinho}
Overfitting trong phân loại thường biểu hiện bằng ranh giới quyết định quá phức tạp, cố gắng ``né'' từng điểm dữ liệu thay vì học xu hướng chung.
\end{ghinho}

% ==========================================================
\section{Underfitting \& Overfitting: Regression -- Ví dụ Underfitting và Overfitting trong hồi quy}

\subsection{Mô tả bài toán hồi quy}

Hồi quy là bài toán dự đoán giá trị liên tục.

Ví dụ:

\begin{itemize}
    \item sản lượng điện tiêu thụ trong ngày;
    \item giá nhà;
    \item nhiệt độ ngày mai;
    \item giá cổ phiếu ngày mai.
\end{itemize}

Ta có một tập điểm dữ liệu và cần tìm một mô hình khớp với xu hướng chung của dữ liệu.

\subsection{Underfitting trong hồi quy}

Nếu dữ liệu có dạng cong nhưng ta dùng đường thẳng, mô hình không đủ khả năng mô tả dữ liệu.

\begin{center}
\begin{tikzpicture}[scale=0.95,>=Stealth]
    \draw[->] (0,0) -- (6,0) node[right] {$x$};
    \draw[->] (0,0) -- (0,4) node[above] {$y$};

    \foreach \x/\y in {0.5/1.0,1.0/1.4,1.5/2.1,2.0/2.6,2.5/2.4,3.0/2.0,3.5/1.75,4.0/2.1,4.5/2.7,5.0/3.2,5.5/3.15}
        \filldraw[blue] (\x,\y) circle (2pt);

    \draw[thick,orange] (0.5,1.5) -- (5.5,2.35);
    \node[orange] at (4.1,0.7) {underfitting};

\end{tikzpicture}
\end{center}

\subsection{Good fitting trong hồi quy}

Mô hình tốt không cần đi qua toàn bộ điểm dữ liệu. Nó cần xấp xỉ được xu hướng tổng quát.

\begin{center}
\begin{tikzpicture}[scale=0.95,>=Stealth]
    \draw[->] (0,0) -- (6,0) node[right] {$x$};
    \draw[->] (0,0) -- (0,4) node[above] {$y$};

    \foreach \x/\y in {0.5/1.0,1.0/1.4,1.5/2.1,2.0/2.6,2.5/2.4,3.0/2.0,3.5/1.75,4.0/2.1,4.5/2.7,5.0/3.2,5.5/3.15}
        \filldraw[blue] (\x,\y) circle (2pt);

    \draw[thick,orange,smooth] plot coordinates {
        (0.45,0.9) (1.0,1.45) (1.6,2.25) (2.2,2.5)
        (2.8,2.2) (3.4,1.85) (4.0,2.10) (4.7,2.85) (5.5,3.2)
    };
    \node[orange] at (4.1,0.7) {good fitting};

\end{tikzpicture}
\end{center}

\subsection{Overfitting trong hồi quy}

Nếu dùng đa thức bậc rất cao, mô hình có thể uốn lượn để đi qua gần như toàn bộ điểm dữ liệu. Sai số train có thể bằng 0, nhưng mô hình mất xu hướng tổng quát.

\begin{center}
\begin{tikzpicture}[scale=0.95,>=Stealth]
    \draw[->] (0,0) -- (6,0) node[right] {$x$};
    \draw[->] (0,0) -- (0,4) node[above] {$y$};

    \foreach \x/\y in {0.5/1.0,1.0/1.4,1.5/2.1,2.0/2.6,2.5/2.4,3.0/2.0,3.5/1.75,4.0/2.1,4.5/2.7,5.0/3.2,5.5/3.15}
        \filldraw[blue] (\x,\y) circle (2pt);

    \draw[thick,orange,smooth] plot coordinates {
        (0.5,1.0) (0.8,1.9) (1.0,1.4) (1.25,1.0)
        (1.5,2.1) (1.8,3.0) (2.0,2.6) (2.3,1.55)
        (2.5,2.4) (2.8,2.9) (3.0,2.0) (3.3,1.2)
        (3.5,1.75) (3.8,2.75) (4.0,2.1) (4.3,1.5)
        (4.5,2.7) (4.8,3.6) (5.0,3.2) (5.3,2.5) (5.5,3.15)
    };
    \node[orange] at (4.1,0.7) {overfitting};

\end{tikzpicture}
\end{center}

\begin{canhbao}
Trong hồi quy, việc tăng bậc đa thức quá cao có thể làm mô hình đi qua tất cả điểm training, nhưng mô hình đó chưa chắc dự đoán tốt cho điểm dữ liệu mới.
\end{canhbao}

% ==========================================================
\section{Causes of Underfitting -- Nguyên nhân gây Underfitting}

\subsection{Mô hình quá đơn giản}

Nguyên nhân chính của underfitting là mô hình không đủ khả năng biểu diễn bài toán.

Ví dụ:

\begin{itemize}
    \item cần đường cong nhưng lại dùng đường thẳng;
    \item cần mạng sâu nhưng lại dùng mạng quá nhỏ;
    \item cần mô hình có nhiều tham số nhưng lại dùng mô hình quá ít tham số.
\end{itemize}

\begin{vidu}
Nếu yêu cầu một mạng nơ-ron rất nhỏ nhận dạng ảnh chó, mèo, xe hơi, máy bay từ ảnh thực tế, mạng đó có thể không đủ khả năng học bài toán dù huấn luyện rất lâu.
\end{vidu}

\subsection{Huấn luyện chưa đủ}

Một nguyên nhân khác là mô hình đủ mạnh nhưng chưa được huấn luyện đủ lâu.

Ví dụ, một bài toán cần huấn luyện khoảng một tuần mới đạt kết quả tốt, nhưng ta dừng sau 2--3 giờ. Khi đó mô hình có thể vẫn chưa học xong và còn underfit.

\begin{ghinho}
Underfitting có thể do mô hình yếu, nhưng cũng có thể do mô hình chưa được huấn luyện đủ.
\end{ghinho}

\subsection{Cách xử lý underfitting}

Một số cách xử lý:

\begin{enumerate}
    \item Dùng mô hình mạnh hơn.
    \item Tăng số lớp hoặc số neuron.
    \item Dùng kiến trúc phù hợp hơn với bài toán.
    \item Huấn luyện lâu hơn.
    \item Bổ sung đặc trưng đầu vào có ý nghĩa.
\end{enumerate}

\begin{vidu}
Với bài toán phân loại ảnh 1000 class, một mạng nhỏ tự thiết kế đơn giản có thể không đủ. Cần dùng các kiến trúc mạnh hơn, ví dụ các mạng tích chập sâu như AlexNet hoặc các kiến trúc hiện đại hơn.
\end{vidu}

% ==========================================================
\section{Causes of Overfitting -- Nguyên nhân gây Overfitting}

\subsection{Mô hình quá phức tạp}

Nếu mô hình quá mạnh so với bài toán, nó có thể học cả những chi tiết nhỏ, nhiễu hoặc điểm bất thường trong tập huấn luyện.

Ví dụ:

\begin{itemize}
    \item bài toán chỉ cần đa thức bậc 4 nhưng ta dùng đa thức bậc 15;
    \item bài toán đơn giản nhưng dùng mạng quá sâu;
    \item số tham số lớn trong khi dữ liệu ít.
\end{itemize}

\subsection{Dữ liệu quá ít}

Nếu dữ liệu ít, mô hình dễ ghi nhớ dữ liệu huấn luyện thay vì học quy luật tổng quát.

\begin{vidu}
Một bài toán phân loại ảnh có 1000 class nhưng chỉ có khoảng 1000 ảnh cho mỗi class. Với bài toán ảnh phức tạp, số lượng này có thể vẫn còn ít, đặc biệt khi ảnh có nhiều biến thể về góc chụp, ánh sáng, vị trí đối tượng và nền ảnh.
\end{vidu}

\subsection{Dữ liệu thiếu đa dạng}

Không chỉ số lượng, tính đa dạng của dữ liệu cũng rất quan trọng.

Nếu tất cả ảnh chó trong tập huấn luyện đều có con chó nằm chính giữa ảnh, mô hình có thể học nhầm rằng ``chó phải nằm ở giữa ảnh''. Khi test với ảnh chó lệch sang bên trái hoặc bên phải, mô hình có thể dự đoán sai.

\begin{ghinho}
Dữ liệu đủ nhiều chưa chắc đã đủ tốt. Dữ liệu còn cần đủ đa dạng để mô hình học được quy luật tổng quát.
\end{ghinho}

% ==========================================================
\section{Giải pháp cho Overfitting: Đơn giản hóa mô hình}

\subsection{Giảm độ phức tạp mô hình}

Nếu mô hình quá mạnh, một hướng xử lý là làm mô hình đơn giản hơn.

Ví dụ:

\begin{itemize}
    \item giảm số lớp;
    \item giảm số neuron;
    \item giảm số tham số;
    \item chọn kiến trúc phù hợp hơn với độ khó của bài toán.
\end{itemize}

\subsection{Loại bỏ bớt đặc trưng}

Nếu input có quá nhiều đặc trưng, mô hình có thể học các quan hệ phụ không cần thiết. Khi đó có thể cân nhắc loại bỏ một số đặc trưng.

\begin{vidu}
Trong bài toán dự đoán sức khỏe, dữ liệu có thể gồm chiều cao, cân nặng, huyết áp, nhịp tim, độ tuổi, đường huyết và nhiều chỉ số khác. Nếu đưa quá nhiều đặc trưng không liên quan, mô hình có thể học các chi tiết linh tinh và overfit.
\end{vidu}

\begin{luuy}
Không nên bỏ đặc trưng một cách tùy tiện. Cần dựa vào hiểu biết bài toán, phân tích dữ liệu hoặc thực nghiệm để quyết định đặc trưng nào nên giữ, đặc trưng nào nên bỏ.
\end{luuy}

% ==========================================================
\section{Giải pháp cho Overfitting: Tăng dữ liệu}

\subsection{Tại sao cần thêm dữ liệu?}

Một nguyên nhân phổ biến gây overfitting là thiếu dữ liệu. Khi dữ liệu quá ít, mô hình dễ học thuộc tập huấn luyện.

Cách trực tiếp nhất là thu thập thêm dữ liệu mới, thật, độc lập và đa dạng.

\subsection{Ví dụ ImageNet}

Trong bài học, giảng viên nhắc đến một bài toán nhận dạng ảnh quy mô lớn với khoảng:

\[
    1000 \text{ class}
\]

và khoảng:

\[
    1.2 \text{ triệu ảnh}
\]

Nghe có vẻ rất nhiều, nhưng nếu chia đều cho 1000 class thì trung bình mỗi class chỉ có khoảng:

\[
    \frac{1,200,000}{1000}=1200 \text{ ảnh/class}
\]

Với bài toán thị giác máy tính phức tạp, con số này vẫn có thể chưa đủ.

\begin{ghinho}
Trong deep learning, dữ liệu ảnh thường cần số lượng rất lớn và độ đa dạng cao. Một vài nghìn ảnh cho mỗi class có thể vẫn chưa đủ nếu bài toán phức tạp.
\end{ghinho}

\subsection{Ví dụ nhận diện khuôn mặt}

Bài toán nhận diện khuôn mặt còn có thể cần lượng dữ liệu lớn hơn rất nhiều. Có những hệ thống được huấn luyện với hàng chục triệu hoặc hàng trăm triệu ảnh khuôn mặt.

Điều này giúp mô hình học được nhiều biến thể:

\begin{itemize}
    \item góc mặt;
    \item ánh sáng;
    \item tuổi tác;
    \item biểu cảm;
    \item kiểu tóc;
    \item chất lượng ảnh;
    \item môi trường chụp.
\end{itemize}

% ==========================================================
\section{Data Augmentation}

\subsection{Data augmentation là gì?}

Data augmentation là kỹ thuật làm tăng số lượng và độ đa dạng của dữ liệu huấn luyện bằng cách biến đổi dữ liệu gốc.

Với ảnh, ta có thể tạo thêm ảnh mới từ ảnh ban đầu bằng các phép biến đổi như:

\begin{itemize}
    \item cắt ảnh;
    \item tịnh tiến vùng nhìn;
    \item lật đối xứng;
    \item xoay ảnh;
    \item thay đổi độ sáng;
    \item thay đổi màu sắc;
    \item thay đổi kênh màu.
\end{itemize}

\begin{ghinho}
Mục tiêu của data augmentation không phải là tạo ảnh kỳ quặc, mà là tạo thêm các tình huống hợp lý có thể xảy ra trong thực tế.
\end{ghinho}

\subsection{Cắt ảnh và dịch cửa sổ}

Giả sử ta có một ảnh con chó. Nếu trong mọi ảnh huấn luyện, con chó luôn nằm chính giữa, mô hình có thể học sai rằng con chó phải nằm ở giữa ảnh.

Ta có thể cắt các vùng khác nhau của ảnh để tạo thêm nhiều ảnh huấn luyện.

\begin{center}
\begin{tikzpicture}[scale=0.9]
    % original image rectangle
    \draw[thick] (0,0) rectangle (5,3);
    \node at (2.5,1.5) {Ảnh gốc};

    % dog
    \draw[blue,thick] (2.5,1.5) ellipse (0.7 and 0.4);
    \node[blue] at (2.5,1.5) {chó};

    % crop windows
    \draw[red,thick] (1.2,0.8) rectangle (3.8,2.4);
    \draw[orange,thick] (0.5,0.5) rectangle (3.1,2.1);
    \draw[green!50!black,thick] (1.9,0.6) rectangle (4.5,2.2);

    \node[red] at (3.7,2.65) {crop 1};
    \node[orange] at (1.4,0.25) {crop 2};
    \node[green!50!black] at (4.2,0.35) {crop 3};
\end{tikzpicture}
\end{center}

\begin{vidu}
Từ một ảnh gốc, ta có thể lấy nhiều vùng crop khác nhau. Nhờ đó, mô hình học rằng đối tượng có thể xuất hiện ở nhiều vị trí khác nhau, không nhất thiết nằm ở trung tâm.
\end{vidu}

\subsection{Lật đối xứng}

Nếu tất cả ảnh voi trong tập huấn luyện đều quay mặt sang phải, mô hình có thể không nhận ra voi quay mặt sang trái. Khi đó, ta có thể dùng phép lật đối xứng.

\[
    \text{ảnh voi quay phải}
    \quad \rightarrow \quad
    \text{ảnh voi quay trái}
\]

\begin{vidu}
Với ảnh động vật, lật ngang thường hợp lý. Một con voi quay mặt sang trái hay sang phải vẫn là con voi.
\end{vidu}

\subsection{Xoay ảnh}

Với một số đối tượng, ta có thể xoay ảnh một góc nhỏ để tăng độ đa dạng.

Ví dụ:

\[
    -30^\circ \leq \theta \leq 30^\circ
\]

Tuy nhiên, không phải bài toán nào cũng cho phép xoay tùy ý.

\begin{luuy}
Cần bảo đảm ảnh sau khi augmentation vẫn có ý nghĩa. Ví dụ, xoay ảnh con voi 180 độ làm voi chổng chân lên trời có thể không phù hợp nếu bài toán thực tế không bao giờ gặp tình huống đó.
\end{luuy}

\subsection{Thay đổi màu sắc và độ sáng}

Nếu tất cả ảnh hoa hồng trong tập huấn luyện có cùng màu sắc và độ sáng, mô hình có thể phụ thuộc quá nhiều vào màu đó.

Ta có thể thay đổi:

\begin{itemize}
    \item độ sáng;
    \item độ tương phản;
    \item sắc độ màu;
    \item các kênh màu RGB.
\end{itemize}

\begin{vidu}
Hoa hồng có thể xuất hiện dưới nhiều điều kiện ánh sáng khác nhau. Nếu chỉ học ảnh hoa hồng có màu đỏ tươi và ánh sáng chuẩn, mô hình có thể nhận sai hoa hồng trong ảnh tối hơn, sáng hơn hoặc có màu lệch.
\end{vidu}

\subsection{Augmentation phải phù hợp với bài toán}

Không phải phép biến đổi nào cũng hợp lý.

\begin{center}
\begin{tabular}{p{0.30\textwidth}p{0.30\textwidth}p{0.30\textwidth}}
\toprule
Đối tượng & Biến đổi thường hợp lý & Biến đổi cần cẩn thận \\
\midrule
Con chó & Crop, lật ngang, đổi sáng & Xoay 180 độ \\
Con voi & Crop, lật ngang, xoay nhẹ & Lộn ngược hoàn toàn \\
Cây búa & Xoay nhiều góc có thể hợp lý & Tùy bài toán cụ thể \\
Ngôi nhà & Đổi sáng, crop, xoay nhẹ & Lật ngược đầu \\
Hoa hồng & Đổi sáng, màu, crop & Biến màu quá giả tạo \\
\bottomrule
\end{tabular}
\end{center}

\begin{canhbao}
Data augmentation không phải là biến đổi dữ liệu vô tội vạ. Nếu tạo ra ảnh quá giả hoặc không có khả năng xảy ra trong thực tế, mô hình có thể học sai.
\end{canhbao}

\subsection{Tác dụng về số lượng dữ liệu}

Trong transcript, giảng viên nêu ví dụ: nếu một ảnh có thể biến thành 2048 ảnh khác bằng các kỹ thuật augmentation, thì tập dữ liệu tăng rất mạnh.

Nếu có:

\[
    1.2 \text{ triệu ảnh}
\]

và mỗi ảnh tạo thành:

\[
    2048 \text{ ảnh}
\]

thì tổng số ảnh xấp xỉ:

\[
    1,200,000 \times 2048
    =
    2,457,600,000
\]

Tức khoảng:

\[
    2.46 \text{ tỷ ảnh}
\]

\begin{ghinho}
Data augmentation vừa tăng số lượng ảnh, vừa tăng độ đa dạng của dữ liệu. Đây là một kỹ thuật rất quan trọng để giảm overfitting trong bài toán ảnh.
\end{ghinho}

% ==========================================================
\section{Train, Validation và Test Set}

\subsection{Vì sao phải chia dữ liệu?}

Nếu chỉ huấn luyện và đánh giá trên cùng một tập dữ liệu, ta không biết mô hình có thực sự tổng quát hóa tốt hay chỉ đang học thuộc dữ liệu.

Do đó, dữ liệu thường được chia thành nhiều tập.

\subsection{Chia thành hai tập}

Cách đơn giản:

\begin{itemize}
    \item training set: dùng để huấn luyện;
    \item test set: dùng để đánh giá sau cùng.
\end{itemize}

Ví dụ:

\[
    80\% \text{ training}
    \quad
    20\% \text{ test}
\]

\subsection{Chia thành ba tập}

Trong thực tế, thường nên chia thành ba tập:

\begin{itemize}
    \item training set;
    \item validation set;
    \item test set.
\end{itemize}

Ví dụ:

\[
    70\% \text{ training}
    \quad
    15\% \text{ validation}
    \quad
    15\% \text{ test}
\]

\begin{center}
\begin{tikzpicture}[scale=1.0]
    \draw[thick] (0,0) rectangle (10,1);
    \fill[blue!20] (0,0) rectangle (7,1);
    \fill[orange!25] (7,0) rectangle (8.5,1);
    \fill[red!20] (8.5,0) rectangle (10,1);

    \node at (3.5,0.5) {Training 70\%};
    \node at (7.75,0.5) {Validation 15\%};
    \node at (9.25,0.5) {Test 15\%};
\end{tikzpicture}
\end{center}

\subsection{Vai trò của training set}

Training set là tập được dùng để cập nhật tham số của mô hình.

\[
    \text{Training set}
    \rightarrow
    \text{tính loss}
    \rightarrow
    \text{backpropagation}
    \rightarrow
    \text{cập nhật trọng số}
\]

\subsection{Vai trò của validation set}

Validation set không dùng để cập nhật trọng số trực tiếp. Nó được dùng để đánh giá mô hình trong quá trình huấn luyện.

Ta dùng validation set để:

\begin{itemize}
    \item theo dõi mô hình có overfit không;
    \item chọn mô hình tốt;
    \item điều chỉnh hyperparameter;
    \item quyết định thời điểm dừng huấn luyện.
\end{itemize}

\begin{ghinho}
Validation set được dùng trong quá trình huấn luyện để đánh giá và lựa chọn mô hình, nhưng không được dùng để cập nhật trọng số như training set.
\end{ghinho}

\subsection{Vai trò của test set}

Test set là tập dữ liệu chỉ dùng để đánh giá cuối cùng.

\begin{canhbao}
Test set không được dùng trong quá trình huấn luyện, không được dùng để chỉnh mô hình, không được dùng để chọn hyperparameter. Test set chỉ được dùng khi ta đã quyết định mô hình cuối cùng.
\end{canhbao}

\subsection{Ví dụ chia dữ liệu}

Giả sử có 1,000 ảnh.

Có thể chia:

\[
    700 \text{ ảnh training}
\]

\[
    150 \text{ ảnh validation}
\]

\[
    150 \text{ ảnh test}
\]

Trong quá trình làm mô hình:

\begin{itemize}
    \item dùng 700 ảnh training để học;
    \item dùng 150 ảnh validation để theo dõi và điều chỉnh;
    \item cất riêng 150 ảnh test, chỉ dùng cuối cùng.
\end{itemize}

% ==========================================================
\section{Nguyên tắc sử dụng Test Set}

\subsection{Không được dùng test set nhiều lần để chỉnh mô hình}

Một lỗi nghiêm trọng là:

\begin{enumerate}
    \item huấn luyện mô hình;
    \item đem test set ra đánh giá;
    \item thấy kết quả chưa tốt;
    \item quay lại chỉnh mô hình;
    \item lại đem cùng test set ra đánh giá;
    \item lặp lại nhiều lần.
\end{enumerate}

Cách làm này làm mô hình hoặc quá trình lựa chọn mô hình dần dần khớp với test set.

\begin{canhbao}
Nếu dùng test set nhiều lần để quyết định chỉnh mô hình, test set không còn là tập kiểm tra độc lập nữa. Kết quả báo cáo sau đó có thể đánh giá sai năng lực thật của mô hình.
\end{canhbao}

\subsection{Vì sao đây là lỗi nghiêm trọng?}

Test set có vai trò mô phỏng dữ liệu mới trong thực tế. Nếu ta liên tục nhìn vào test set để chỉnh mô hình, ta đã làm mất tính độc lập của nó.

Khi đó, kết quả trên test set không còn phản ánh khả năng tổng quát hóa thật.

\subsection{Liên hệ với nghiên cứu khoa học}

Trong nghiên cứu khoa học, nếu báo cáo kết quả trên một test set đã bị dùng nhiều lần để chỉnh mô hình, kết quả có thể gây hiểu lầm rằng mô hình tốt hơn thực tế.

Điều này có thể bị xem là sai nguyên tắc nghiên cứu, thậm chí là gian lận nếu cố ý.

\begin{ghinho}
Test set là bài kiểm tra cuối cùng. Nếu đã dùng test set để chỉnh mô hình, cần có một test set độc lập mới để đánh giá cuối cùng.
\end{ghinho}

\subsection{Nếu lỡ dùng test set thì làm sao?}

Nếu đã lỡ dùng test set trong quá trình phát triển mô hình, cách xử lý đúng là:

\begin{itemize}
    \item thu thập thêm dữ liệu mới để làm test set độc lập;
    \item hoặc chia lại dữ liệu và tạo một test set mới chưa từng dùng;
    \item báo cáo rõ quy trình đánh giá.
\end{itemize}

% ==========================================================
\section{Cross-validation}

\subsection{Ý tưởng}

Khi dữ liệu ít, việc tách riêng một test set cố định có thể làm lãng phí dữ liệu. Một hướng xử lý là kiểm tra chéo, tiếng Anh là \textit{cross-validation}.

Ý tưởng chung:

\begin{itemize}
    \item chia dữ liệu thành nhiều phần;
    \item mỗi lần dùng một phần khác nhau làm tập kiểm tra;
    \item huấn luyện và đánh giá nhiều lần;
    \item tổng hợp kết quả.
\end{itemize}

\subsection{Ví dụ K-fold cross-validation}

Giả sử chia dữ liệu thành 5 phần:

\[
    D_1,D_2,D_3,D_4,D_5
\]

Ta thực hiện 5 lần:

\begin{center}
\begin{tabular}{ccc}
\toprule
Lần & Dùng để train & Dùng để kiểm tra \\
\midrule
1 & $D_2,D_3,D_4,D_5$ & $D_1$ \\
2 & $D_1,D_3,D_4,D_5$ & $D_2$ \\
3 & $D_1,D_2,D_4,D_5$ & $D_3$ \\
4 & $D_1,D_2,D_3,D_5$ & $D_4$ \\
5 & $D_1,D_2,D_3,D_4$ & $D_5$ \\
\bottomrule
\end{tabular}
\end{center}

Sau đó lấy trung bình kết quả đánh giá.

\begin{luuy}
Trong transcript, cross-validation chỉ được nhắc như một hướng liên quan khi cần chia lại và đánh giá dữ liệu nhiều lần. Nội dung chi tiết về cross-validation có thể học sâu hơn ở phần machine learning.
\end{luuy}

% ==========================================================
\section{Early Stopping -- Early Stopping}

\subsection{Early stopping là gì?}

Early stopping, hay dừng sớm, là kỹ thuật dừng quá trình huấn luyện khi mô hình bắt đầu có dấu hiệu overfitting.

Ta theo dõi loss trên:

\begin{itemize}
    \item training set;
    \item validation set.
\end{itemize}

Thông thường, training loss sẽ giảm dần. Nhưng nếu validation loss giảm đến một mức nào đó rồi bắt đầu tăng trở lại, đó là dấu hiệu mô hình đang overfit.

\subsection{Minh họa}

\begin{center}
\begin{tikzpicture}[scale=1.0,>=Stealth]
    \draw[->] (0,0) -- (8,0) node[right] {Thời gian huấn luyện};
    \draw[->] (0,0) -- (0,5) node[above] {Loss};

    % training loss
    \draw[thick,blue,smooth] plot coordinates {
        (0.5,4.4) (1.2,3.6) (2.0,2.8) (3.0,2.1)
        (4.0,1.6) (5.0,1.25) (6.0,1.05) (7.2,0.9)
    };
    \node[blue] at (6.7,1.3) {training loss};

    % validation loss
    \draw[thick,red,smooth] plot coordinates {
        (0.5,4.5) (1.2,3.7) (2.0,2.9) (3.0,2.2)
        (4.0,1.8) (5.0,1.65) (6.0,1.85) (7.2,2.3)
    };
    \node[red] at (6.6,2.6) {validation loss};

    % early stop point
    \draw[dashed] (5.0,0) -- (5.0,1.65);
    \filldraw[black] (5.0,1.65) circle (2pt);
    \node[below] at (5.0,0) {dừng};

\end{tikzpicture}
\end{center}

\subsection{Ý nghĩa}

Giả sử ban đầu dự định huấn luyện 100,000 bước. Nhưng sau 70,000 bước, validation loss bắt đầu tăng, trong khi training loss vẫn giảm. Khi đó nên dừng huấn luyện và giữ mô hình tại thời điểm validation loss tốt nhất.

\begin{ghinho}
Early stopping giúp tránh tình trạng mô hình tiếp tục học quá sát training set sau khi đã đạt khả năng tổng quát hóa tốt nhất trên validation set.
\end{ghinho}

\subsection{Không phải cứ early stopping là mô hình tốt}

Dừng sớm chỉ giúp giảm overfitting. Nó không đảm bảo mô hình cuối cùng chắc chắn tốt.

Nếu mô hình quá yếu, dữ liệu sai, hoặc bài toán chưa được thiết kế đúng, early stopping vẫn không cứu được mô hình.

% ==========================================================
\section{Dropout -- Dropout}

\subsection{Dropout là gì?}

Dropout là kỹ thuật trong quá trình huấn luyện, tại mỗi thời điểm, ta ngẫu nhiên loại bỏ một số neuron khỏi mạng với một xác suất nào đó.

Ví dụ, nếu xác suất giữ lại neuron là:

\[
    p = 0.5
\]

thì trung bình khoảng một nửa số neuron được giữ lại, một nửa bị tạm thời bỏ qua trong bước huấn luyện đó.

\begin{ghinho}
Dropout chỉ dùng trong quá trình training. Khi test hoặc sử dụng mô hình thật, ta dùng toàn bộ mạng.
\end{ghinho}

\subsection{Một thời điểm trong dropout là gì?}

Trong bài học, một ``thời điểm'' có thể hiểu là một lần đưa một mẫu dữ liệu vào mạng.

Ví dụ:

\begin{itemize}
    \item đưa ảnh thứ nhất vào, dropout một số neuron, tính xuôi, tính ngược, cập nhật;
    \item đưa ảnh thứ hai vào, dropout một nhóm neuron khác, tính xuôi, tính ngược, cập nhật;
    \item tiếp tục như vậy cho các mẫu tiếp theo.
\end{itemize}

\subsection{Minh họa dropout}

\begin{center}
\begin{tikzpicture}[x=2cm,y=1cm,>=Stealth,
    neuron/.style={circle,draw,minimum size=8mm},
    dropped/.style={circle,draw,dashed,gray,minimum size=8mm},
    every node/.style={font=\small}
]

% Input
\node[neuron] (i1) at (0,1) {$x_1$};
\node[neuron] (i2) at (0,0) {$x_2$};
\node[neuron] (i3) at (0,-1) {$x_3$};

% Hidden
\node[neuron] (h1) at (1.8,1.5) {$h_1$};
\node[dropped] (h2) at (1.8,0.5) {$h_2$};
\node[neuron] (h3) at (1.8,-0.5) {$h_3$};
\node[dropped] (h4) at (1.8,-1.5) {$h_4$};

% Output
\node[neuron] (o1) at (3.6,0.5) {$o_1$};
\node[neuron] (o2) at (3.6,-0.5) {$o_2$};

% Connections only active
\foreach \i in {i1,i2,i3}
    \foreach \h in {h1,h3}
        \draw[->] (\i) -- (\h);

\foreach \h in {h1,h3}
    \foreach \o in {o1,o2}
        \draw[->] (\h) -- (\o);

\node[gray] at (1.8,-2.1) {Neuron nét đứt bị dropout};

\end{tikzpicture}
\end{center}

\subsection{Vì sao dropout giúp giảm overfitting?}

Có hai cách hiểu trực quan.

\subsection{Góc nhìn ensemble}

Nếu mỗi lần dropout tạo ra một cấu hình mạng khác nhau, thì trong quá trình huấn luyện, ta giống như đang huấn luyện rất nhiều mạng con khác nhau.

Khi test, ta dùng toàn bộ mạng, có thể xem như sự kết hợp của nhiều mạng con đã được huấn luyện.

\begin{vidu}
Thay vì huấn luyện một mạng duy nhất, dropout khiến mô hình trải qua rất nhiều phiên bản mạng con. Điều này gần giống tư tưởng ensemble: nhiều mô hình cùng đóng góp để tạo ra kết quả ổn định hơn.
\end{vidu}

\subsection{Góc nhìn giảm phụ thuộc}

Nếu không có dropout, một neuron có thể phụ thuộc quá mạnh vào một vài neuron khác. Khi dropout ngẫu nhiên loại bỏ neuron, mạng buộc phải học cách sử dụng thông tin từ nhiều nguồn khác nhau.

Điều này giúp mạng mạnh mẽ hơn và ít phụ thuộc vào một đường truyền thông tin cụ thể.

\begin{ghinho}
Dropout buộc các neuron phải ``học cách làm việc với nhiều đồng đội khác nhau'', thay vì phụ thuộc quá mức vào một vài neuron cố định.
\end{ghinho}

\subsection{Dropout và robust model}

Một mô hình robust là mô hình vẫn hoạt động tốt trong nhiều tình huống khác nhau.

Dropout giúp tăng tính robust vì trong quá trình train, mạng thường xuyên phải làm việc trong điều kiện thiếu một số neuron. Nhờ đó, mạng học được biểu diễn phân tán hơn và ít nhạy với từng neuron riêng lẻ.

\subsection{Vì sao không dùng dropout khi test?}

Trong training, mỗi lần chỉ một phần neuron hoạt động. Nhưng khi test, ta muốn dùng toàn bộ năng lực của mạng đã học.

Vì vậy:

\[
    \text{Training: dùng dropout}
\]

\[
    \text{Testing: không dùng dropout}
\]

\subsection{Vấn đề scale tín hiệu}

Giả sử trong training, mỗi neuron được giữ lại với xác suất:

\[
    p=0.5
\]

Khi đó, trung bình một neuron ở lớp sau chỉ nhận được khoảng một nửa số tín hiệu so với khi dùng toàn bộ mạng.

Nếu khi test ta bật toàn bộ neuron mà không điều chỉnh gì, tổng tín hiệu đầu vào có thể tăng lên so với lúc training.

Theo cách trình bày trong bài học, để làm cho mức tín hiệu giữa training và testing tương thích hơn, output của neuron có thể được nhân với hệ số $p$ khi dùng toàn bộ mạng.

\[
    out_{\text{test}} = p \cdot out
\]

Ví dụ nếu:

\[
    p=0.5
\]

thì khi test:

\[
    out_{\text{test}}=0.5\cdot out
\]

\begin{luuy}
Trong nhiều thư viện hiện đại, người ta thường dùng biến thể gọi là inverted dropout: scale ngay trong lúc training, nên khi test không cần nhân thêm. Ý chính cần nhớ trong bài này là phải giữ mức tín hiệu giữa training và testing nhất quán.
\end{luuy}

\subsection{Xác suất dropout}

Giá trị $p$ là một hyperparameter. Ví dụ phổ biến:

\[
    p=0.5
\]

Nhưng không có một giá trị tốt nhất cho mọi bài toán. Trong thực tế, người ta thường dùng các giá trị đã được nghiên cứu trước hoặc thử nghiệm để chọn giá trị phù hợp.

\begin{luuy}
Không nên hiểu rằng $p=0.5$ luôn là tốt nhất. Nó chỉ là một lựa chọn phổ biến trong nhiều tình huống.
\end{luuy}

% ==========================================================
\section{Tổng hợp các kỹ thuật xử lý Overfitting}

\subsection{Các kỹ thuật chính}

\begin{center}
\begin{tabular}{p{0.28\textwidth}p{0.60\textwidth}}
\toprule
Kỹ thuật & Ý nghĩa \\
\midrule
Đơn giản hóa mô hình & Giảm số lớp, số neuron hoặc độ phức tạp để tránh học nhiễu \\
Feature selection & Bỏ bớt đặc trưng không cần thiết hoặc gây nhiễu \\
Thêm dữ liệu & Thu thập thêm dữ liệu thật, độc lập và đa dạng \\
Data augmentation & Tạo thêm dữ liệu hợp lý từ dữ liệu gốc \\
Train/validation/test split & Đánh giá đúng khả năng tổng quát hóa \\
Early stopping & Dừng khi validation loss bắt đầu xấu đi \\
Dropout & Ngẫu nhiên loại neuron trong training để giảm phụ thuộc và tăng robust \\
Cross-validation & Đánh giá ổn định hơn khi dữ liệu ít \\
\bottomrule
\end{tabular}
\end{center}

\subsection{Kỹ thuật nào quan trọng nhất?}

Không có một kỹ thuật duy nhất giải quyết mọi trường hợp. Trong thực tế thường phải kết hợp nhiều kỹ thuật.

Ví dụ với bài toán ảnh:

\begin{itemize}
    \item dùng mô hình đủ mạnh;
    \item dùng data augmentation;
    \item chia train/validation/test đúng;
    \item theo dõi validation loss;
    \item dùng early stopping;
    \item dùng dropout hoặc các kỹ thuật regularization khác.
\end{itemize}

% ==========================================================
\section{Technical Glossary -- Bảng thuật ngữ chuyên ngành}

\small
\begin{longtable}{p{0.26\textwidth}p{0.48\textwidth}p{0.19\textwidth}}
\toprule
\textbf{Thuật ngữ} & \textbf{Giải thích} & \textbf{Ví dụ} \\
\midrule
\endfirsthead

\toprule
\textbf{Thuật ngữ} & \textbf{Giải thích} & \textbf{Ví dụ} \\
\midrule
\endhead

Underfitting & Mô hình chưa khớp dữ liệu, chưa học được xu hướng tổng quát. & Dùng đường thẳng cho dữ liệu cong. \\

Overfitting & Mô hình quá khớp training data, học cả nhiễu. & Đa thức bậc cao đi qua mọi điểm. \\

Training set & Tập dữ liệu dùng để cập nhật tham số mô hình. & 70\% dữ liệu. \\

Validation set & Tập dùng để đánh giá trong quá trình huấn luyện và chọn mô hình. & 15\% dữ liệu. \\

Test set & Tập chỉ dùng để đánh giá cuối cùng. & 15\% dữ liệu còn lại. \\

Data augmentation & Tạo thêm dữ liệu từ dữ liệu gốc bằng biến đổi hợp lý. & Crop, flip, rotate ảnh. \\

Crop & Cắt một vùng trong ảnh gốc để tạo ảnh mới. & Cắt vùng chứa con chó. \\

Reflection/Flip & Lật ảnh theo chiều ngang hoặc chiều dọc. & Lật ảnh voi quay phải thành quay trái. \\

Rotation & Xoay ảnh một góc nào đó. & Xoay ảnh búa 30 độ. \\

Color augmentation & Thay đổi màu sắc, độ sáng hoặc kênh màu. & Làm ảnh hoa sáng hơn. \\

Early stopping & Dừng sớm khi validation loss bắt đầu tăng. & Dừng sau 5 ngày thay vì 7 ngày. \\

Dropout & Ngẫu nhiên loại bỏ một số neuron trong quá trình training. & Giữ lại neuron với xác suất 0.5. \\

Ensemble & Kết hợp nhiều mô hình để có kết quả ổn định hơn. & Trung bình kết quả nhiều mạng. \\

Robust & Mạnh mẽ, ít bị ảnh hưởng bởi thay đổi nhỏ hoặc thiếu một phần thông tin. & Mạng không phụ thuộc một neuron. \\

Hyperparameter & Tham số do người thiết kế chọn trước hoặc điều chỉnh. & Dropout rate, số lớp. \\

Cross-validation & Kiểm tra chéo, chia dữ liệu nhiều lần để đánh giá ổn định hơn. & 5-fold cross-validation. \\

Feature selection & Chọn hoặc loại bỏ đặc trưng đầu vào. & Bỏ chỉ số không liên quan. \\

Generalization & Khả năng tổng quát hóa trên dữ liệu mới. & Test accuracy cao. \\

\bottomrule
\end{longtable}
\normalsize

% ==========================================================
\section{Key Concepts Summary -- Tóm tắt kiến thức trọng tâm}

\subsection{Các ý chính cần nhớ}

\begin{enumerate}
    \item Underfitting xảy ra khi mô hình quá đơn giản hoặc chưa được huấn luyện đủ.
    \item Overfitting xảy ra khi mô hình quá phức tạp, dữ liệu ít hoặc thiếu đa dạng.
    \item Underfitting dễ nhận ra hơn vì loss thường còn cao.
    \item Overfitting nguy hiểm hơn vì training loss thấp khiến ta dễ nhầm rằng mô hình tốt.
    \item Để giảm underfitting, có thể dùng mô hình mạnh hơn hoặc huấn luyện lâu hơn.
    \item Để giảm overfitting, có thể đơn giản hóa mô hình, bỏ bớt đặc trưng, thêm dữ liệu hoặc dùng data augmentation.
    \item Data augmentation giúp tăng cả số lượng và độ đa dạng dữ liệu.
    \item Augmentation phải tạo ra dữ liệu hợp lý với bài toán thực tế.
    \item Dữ liệu nên được chia thành training, validation và test set.
    \item Test set chỉ được dùng để đánh giá cuối cùng.
    \item Dùng test set nhiều lần để chỉnh mô hình là sai nguyên tắc.
    \item Early stopping dùng validation loss để dừng trước khi mô hình overfit.
    \item Dropout ngẫu nhiên loại neuron trong training để giảm phụ thuộc và tăng khả năng tổng quát hóa.
    \item Dropout chỉ dùng trong training, không dùng trong testing.
\end{enumerate}

\subsection{Công thức và con số cần nhớ}

\subsection*{Tỷ lệ chia dữ liệu ví dụ}

\[
    70\% \text{ training}
    \quad
    15\% \text{ validation}
    \quad
    15\% \text{ test}
\]

\subsection*{Số ảnh trung bình mỗi class}

\[
    \frac{1,200,000}{1000}=1200 \text{ ảnh/class}
\]

\subsection*{Data augmentation 2048 lần}

\[
    1,200,000 \times 2048
    =
    2,457,600,000
\]

\[
    \approx 2.46 \text{ tỷ ảnh}
\]

\subsection*{Dropout scale theo xác suất giữ neuron}

\[
    out_{\text{test}} = p \cdot out
\]

% ==========================================================
\section{Review Questions -- Câu hỏi ôn tập}

\begin{enumerate}
    \item Underfitting là gì?
    \item Overfitting là gì?
    \item Vì sao overfitting nguy hiểm hơn underfitting trong quá trình đánh giá mô hình?
    \item Nêu hai nguyên nhân gây underfitting.
    \item Nêu ba nguyên nhân gây overfitting.
    \item Vì sao dữ liệu nhiều nhưng thiếu đa dạng vẫn có thể gây overfitting?
    \item Data augmentation là gì?
    \item Nêu ba kỹ thuật data augmentation cho ảnh.
    \item Vì sao không nên xoay ảnh một cách tùy tiện?
    \item Training set dùng để làm gì?
    \item Validation set dùng để làm gì?
    \item Test set dùng để làm gì?
    \item Vì sao không được dùng test set nhiều lần để chỉnh mô hình?
    \item Early stopping hoạt động dựa trên dấu hiệu nào?
    \item Dropout là gì?
    \item Dropout được dùng trong training hay testing?
    \item Vì sao dropout có thể giúp giảm overfitting?
    \item Xác suất giữ neuron $p=0.5$ có phải luôn là tốt nhất không?
\end{enumerate}

% ==========================================================
\section{Practice Exercises -- Bài tập vận dụng}

\subsection*{Bài tập 1: Nhận diện lỗi mô hình}
\addcontentsline{toc}{section}{Bài tập 1: Nhận diện lỗi mô hình}

Xác định trường hợp sau là underfitting hay overfitting:

\begin{enumerate}
    \item Training loss cao, validation loss cao.
    \item Training loss rất thấp, validation loss cao.
    \item Mô hình dùng đường thẳng cho dữ liệu có dạng cong.
    \item Mô hình dùng đa thức bậc rất cao và đi qua mọi điểm training.
\end{enumerate}

\subsection*{Lời giải gợi ý}

\begin{enumerate}
    \item Underfitting.
    \item Overfitting.
    \item Underfitting.
    \item Overfitting.
\end{enumerate}

\subsection*{Bài tập 2: Chia dữ liệu}
\addcontentsline{toc}{section}{Bài tập 2: Chia dữ liệu}

Bạn có 10,000 ảnh. Hãy chia theo tỷ lệ:

\[
    70\% \text{ training}, \quad 15\% \text{ validation}, \quad 15\% \text{ test}
\]

Tính số ảnh trong mỗi tập.

\subsection*{Lời giải gợi ý}

Training:

\[
    10,000 \times 0.70 = 7,000
\]

Validation:

\[
    10,000 \times 0.15 = 1,500
\]

Test:

\[
    10,000 \times 0.15 = 1,500
\]

\subsection*{Bài tập 3: Data augmentation}
\addcontentsline{toc}{section}{Bài tập 3: Data augmentation}

Một tập dữ liệu có 5,000 ảnh. Nếu mỗi ảnh được augmentation thành 100 ảnh, tổng số ảnh sau augmentation là bao nhiêu?

\subsection*{Lời giải gợi ý}

\[
    5,000 \times 100 = 500,000
\]

Vậy sau augmentation có 500,000 ảnh.

\subsection*{Bài tập 4: Chọn augmentation phù hợp}
\addcontentsline{toc}{section}{Bài tập 4: Chọn augmentation phù hợp}

Với từng bài toán, hãy cho biết phép augmentation nào hợp lý hơn:

\begin{enumerate}
    \item Nhận dạng cây búa trong ảnh.
    \item Nhận dạng ngôi nhà bình thường.
    \item Nhận dạng hoa hồng trong nhiều điều kiện ánh sáng.
    \item Nhận dạng con voi trong ảnh tự nhiên.
\end{enumerate}

\subsection*{Lời giải gợi ý}

\begin{enumerate}
    \item Có thể xoay ảnh nhiều góc nếu cây búa có thể xuất hiện ở nhiều hướng.
    \item Có thể crop, đổi sáng, xoay nhẹ; không nên lật ngược nhà 180 độ nếu thực tế không có nhà lộn ngược.
    \item Nên đổi sáng, màu sắc, tương phản.
    \item Có thể crop, lật ngang, xoay nhẹ; không nên lật ngược voi nếu bài toán không có tình huống đó.
\end{enumerate}

\subsection*{Bài tập 5: Dropout}
\addcontentsline{toc}{section}{Bài tập 5: Dropout}

Một layer có 100 neuron. Nếu dùng dropout với xác suất giữ lại:

\[
    p=0.5
\]

Trung bình mỗi lần training có bao nhiêu neuron được giữ lại?

\subsection*{Lời giải gợi ý}

\[
    100 \times 0.5 = 50
\]

Trung bình có 50 neuron được giữ lại.

% ==========================================================
\section*{Kết luận}
\addcontentsline{toc}{section}{Kết luận}

ANN8 giải thích sâu hơn nguyên nhân và cách xử lý underfitting, overfitting. Underfitting thường xảy ra khi mô hình quá đơn giản hoặc chưa được huấn luyện đủ. Overfitting thường xảy ra khi mô hình quá phức tạp, dữ liệu quá ít hoặc thiếu đa dạng.

Để xử lý underfitting, ta có thể dùng mô hình mạnh hơn hoặc huấn luyện lâu hơn. Để xử lý overfitting, ta có nhiều kỹ thuật như đơn giản hóa mô hình, bỏ bớt đặc trưng, tăng dữ liệu, data augmentation, chia đúng train/validation/test, early stopping và dropout.

Một điểm rất quan trọng của bài học là nguyên tắc sử dụng test set. Test set chỉ dùng để đánh giá cuối cùng, không dùng để chỉnh mô hình. Nếu dùng test set nhiều lần trong quá trình phát triển, kết quả đánh giá cuối cùng sẽ không còn đáng tin cậy.

Cuối cùng, dropout được trình bày như một kỹ thuật rất phổ biến để giảm overfitting. Nó ngẫu nhiên loại bỏ neuron trong quá trình training, buộc mạng học biểu diễn mạnh mẽ hơn và ít phụ thuộc vào một số neuron cụ thể.



\chapter{ANN9: Optimization Algorithms -- Thuật toán tối ưu}
\input{chapters/ann9.tex}

\chapter{Softmax Function -- Hàm Softmax}
\input{chapters/softmax.tex}

% ==========================================================
% PHẦN II: MẠNG NƠ-RON TÍCH CHẬP
% ==========================================================
\part{Mạng Nơ-ron Tích chập}

\chapter{CNN1: Convolutional Neural Networks -- Mạng tích chập}
% ==========================================================
\section*{Learning Objectives -- Mục tiêu bài học}
\addcontentsline{toc}{section}{Mục tiêu bài học}

Bài học này mở đầu cho phần \textit{Convolutional Neural Network}, viết tắt là CNN, hay còn gọi là mạng nơ-ron tích chập.

Sau khi học xong, người học cần nắm được:

\begin{enumerate}[label=\arabic*.]
    \item CNN là gì và vì sao CNN quan trọng trong deep learning.
    \item Vì sao CNN được thiết kế đặc biệt phù hợp cho dữ liệu hình ảnh.
    \item Nhận biết cấu trúc tổng quát của một mạng CNN.
    \item Phân biệt CNN với mạng nơ-ron fully connected thông thường.
    \item Hiểu vì sao mạng fully connected gặp vấn đề khi xử lý ảnh lớn.
    \item Hiểu hai ý tưởng quan trọng của CNN:
    \begin{enumerate}
        \item \textit{local connectivity} -- kết nối cục bộ;
        \item \textit{weight sharing} -- trọng số dùng chung.
    \end{enumerate}
    \item Hiểu kernel là gì.
    \item Biết cách thực hiện phép convolution bằng cách trượt kernel trên ảnh.
    \item Hiểu activation map là gì.
    \item Biết cách nối phần convolution với fully connected layer trong ví dụ nhận dạng chữ X/O.
\end{enumerate}

% ==========================================================
\section{Convolutional Neural Network -- Mạng nơ-ron tích chập}

\subsection{CNN là gì?}

CNN là viết tắt của:

\[
    \text{Convolutional Neural Network}
\]

Trong tiếng Việt, CNN thường được gọi là:

\[
    \text{mạng nơ-ron tích chập}
\]

Nếu các mạng nơ-ron trước đây ta học là mạng nơ-ron thông thường, thì CNN là một kiểu mạng nơ-ron đặc biệt hơn, được thiết kế để xử lý tốt dữ liệu có cấu trúc không gian, đặc biệt là ảnh.

\begin{ghinho}
CNN là một kiến trúc rất quan trọng trong deep learning, đặc biệt trong lĩnh vực thị giác máy tính.
\end{ghinho}

\subsection{CNN được thiết kế để giải quyết bài toán gì?}

CNN ban đầu được thiết kế chủ yếu để xử lý các bài toán liên quan đến hình ảnh.

Ví dụ:

\begin{itemize}
    \item phân loại ảnh;
    \item nhận dạng khuôn mặt;
    \item nhận dạng vật thể;
    \item phát hiện vật thể trong ảnh;
    \item phân tích ảnh y tế;
    \item nhận dạng hành động;
    \item hỗ trợ xe tự lái.
\end{itemize}

Trong bài học này, ta tập trung vào bài toán:

\[
    \text{ảnh đầu vào}
    \longrightarrow
    \text{class dự đoán}
\]

Ví dụ:

\[
    \text{ảnh}
    \longrightarrow
    \text{chó, mèo, xe hơi, máy bay, ...}
\]

% ==========================================================
\section{CNN Applications -- Ứng dụng của CNN}

\subsection{Thị giác máy tính}

CNN được sử dụng rất nhiều trong \textit{computer vision} -- thị giác máy tính.

Một số bài toán tiêu biểu:

\begin{itemize}
    \item nhận dạng khuôn mặt;
    \item nhận dạng đối tượng;
    \item phân loại ảnh;
    \item mô tả nội dung ảnh;
    \item nhận dạng hành động trong video.
\end{itemize}

\begin{vidu}
Nếu đưa vào một tấm ảnh, hệ thống có thể gán nhãn rằng đó là ``một con chó đang chơi trên bãi cỏ'' hoặc ``một người đang lướt sóng ở bãi biển''.
\end{vidu}

\subsection{Một số lĩnh vực khác}

Ngoài ảnh, CNN cũng có thể được dùng trong các bài toán khác như:

\begin{itemize}
    \item nhận dạng giọng nói;
    \item phân loại văn bản;
    \item xử lý tín hiệu;
    \item phân tích chuỗi thời gian.
\end{itemize}

Tuy nhiên, trong phần CNN đầu tiên này, ta chủ yếu khảo sát CNN trong bối cảnh xử lý ảnh.

% ==========================================================
\section{General CNN Architecture -- Cấu trúc tổng quát của CNN}

\subsection{Các thành phần chính}

Một mạng CNN dùng cho phân loại ảnh thường có cấu trúc tổng quát như sau:

\[
    \text{Input image}
    \rightarrow
    \text{Convolution layers}
    \rightarrow
    \text{Pooling layers}
    \rightarrow
    \text{Fully connected layers}
    \rightarrow
    \text{Output}
\]

Trong đó:

\begin{itemize}
    \item \textbf{Convolution layer}: lớp tích chập, dùng để trích xuất đặc trưng từ ảnh.
    \item \textbf{Pooling layer}: lớp giảm kích thước dữ liệu, giúp giữ lại thông tin quan trọng.
    \item \textbf{Fully connected layer}: lớp kết nối đầy đủ, dùng để phân loại dựa trên các đặc trưng đã học.
\end{itemize}

\begin{center}
\begin{tikzpicture}[node distance=1cm,>=Stealth,
    box/.style={rectangle,draw,rounded corners,align=center,
    minimum width=2.7cm,minimum height=0.9cm,font=\small}
]

\node[box] (input) {Ảnh đầu vào};
\node[box,right=of input] (conv) {Convolution\\Pooling};
\node[box,right=of conv] (fc) {Fully connected};
\node[box,right=of fc] (output) {Output\\class};

\draw[->] (input) -- (conv);
\draw[->] (conv) -- (fc);
\draw[->] (fc) -- (output);

\end{tikzpicture}
\end{center}

\begin{ghinho}
Trong CNN, phần đầu thường dùng convolution và pooling để học đặc trưng ảnh. Phần sau dùng fully connected layer để đưa ra quyết định phân loại.
\end{ghinho}

\subsection{Số lớp trong CNN}

Trong hình minh họa, ta có thể thấy một vài convolution layer và pooling layer. Tuy nhiên, trong thực tế, số lớp có thể rất nhiều.

Một CNN có thể có:

\begin{itemize}
    \item vài lớp convolution;
    \item vài chục lớp convolution;
    \item thậm chí hàng trăm lớp trong các mạng rất sâu.
\end{itemize}

\begin{luuy}
Hình vẽ trong slide thường chỉ minh họa ý tưởng. Một mạng CNN thực tế có thể phức tạp hơn nhiều.
\end{luuy}

% ==========================================================
\section{Image Representation -- Biểu diễn ảnh trong máy tính}

\subsection{Ảnh xám}

Một ảnh xám có thể được biểu diễn bằng một ma trận.

Ví dụ, ảnh kích thước $H \times W$:

\[
    X =
    \begin{bmatrix}
    x_{11} & x_{12} & \cdots & x_{1W}\\
    x_{21} & x_{22} & \cdots & x_{2W}\\
    \vdots & \vdots & \ddots & \vdots\\
    x_{H1} & x_{H2} & \cdots & x_{HW}
    \end{bmatrix}
\]

Mỗi phần tử $x_{ij}$ là một pixel.

Với ảnh 8-bit, giá trị pixel thường nằm trong khoảng:

\[
    0 \leq x_{ij} \leq 255
\]

\subsection{Ảnh màu}

Ảnh màu thường có 3 kênh:

\[
    R,\quad G,\quad B
\]

Do đó, ảnh màu kích thước $H \times W$ được biểu diễn dưới dạng:

\[
    H \times W \times 3
\]

\begin{vidu}
Một ảnh màu kích thước $32 \times 32$ có số giá trị là:

\[
    32 \cdot 32 \cdot 3 = 3072
\]

Tức là ảnh này có 3072 giá trị đầu vào.
\end{vidu}

% ==========================================================
\section{Problem with Fully Connected Networks -- Vấn đề của mạng fully connected khi xử lý ảnh}

\subsection{Flatten ảnh thành vector}

Với mạng nơ-ron thông thường, khi xử lý ảnh, ta thường biến ảnh thành một vector dài.

Ví dụ, ảnh màu kích thước:

\[
    32 \times 32 \times 3
\]

sẽ được biến thành vector có:

\[
    32 \cdot 32 \cdot 3 = 3072
\]

phần tử.

Sau đó, mỗi neuron ở lớp kế tiếp sẽ nhận toàn bộ 3072 giá trị này.

\subsection{Số trọng số tăng rất nhanh}

Nếu một neuron nhận toàn bộ 3072 input, thì riêng neuron đó đã cần 3072 trọng số.

Nếu ảnh lớn hơn, vấn đề nghiêm trọng hơn.

Ví dụ, ảnh màu kích thước:

\[
    200 \times 200 \times 3
\]

có số input là:

\[
    200 \cdot 200 \cdot 3 = 120000
\]

Như vậy, một neuron fully connected cần:

\[
    120000
\]

trọng số.

Nếu lớp đó có 1000 neuron, số trọng số là:

\[
    120000 \cdot 1000 = 120000000
\]

Đó mới chỉ là một lớp.

\begin{canhbao}
Khi dùng mạng fully connected cho ảnh lớn, số lượng tham số có thể trở nên khổng lồ. Điều này làm tăng chi phí tính toán và dễ dẫn đến overfitting.
\end{canhbao}

\subsection{Liên hệ với overfitting}

Khi mô hình có quá nhiều tham số, nó có khả năng học quá sát dữ liệu huấn luyện.

Điều này dễ dẫn đến:

\[
    \text{overfitting}
\]

Tức là mô hình làm tốt trên tập huấn luyện nhưng làm kém trên dữ liệu mới.

% ==========================================================
\section{Key Ideas of CNN -- Hai ý tưởng quan trọng của CNN}

\subsection{Local connectivity -- Kết nối cục bộ}

Trong mạng fully connected, một neuron có thể nhận toàn bộ input.

Trong CNN, một neuron chỉ nhìn một vùng nhỏ của ảnh.

Ý tưởng này gọi là:

\[
    \text{local connectivity}
\]

hay:

\[
    \text{kết nối cục bộ}
\]

\begin{vidu}
Thay vì nhìn toàn bộ ảnh $200 \times 200 \times 3$, một kernel có thể chỉ nhìn một vùng nhỏ $3 \times 3 \times 3$.
\end{vidu}

Số trọng số khi đó giảm rất mạnh:

\[
    3 \cdot 3 \cdot 3 = 27
\]

so với:

\[
    200 \cdot 200 \cdot 3 = 120000
\]

\begin{ghinho}
Local connectivity giúp CNN giảm số tham số bằng cách chỉ xử lý từng vùng nhỏ của ảnh.
\end{ghinho}

\subsection{Weight sharing -- Trọng số dùng chung}

Ý tưởng thứ hai là:

\[
    \text{weight sharing}
\]

hay:

\[
    \text{trọng số dùng chung}
\]

Thay vì mỗi vị trí trong ảnh có một bộ trọng số riêng, CNN dùng cùng một kernel để quét qua nhiều vị trí khác nhau.

Nếu kernel là:

\[
    K =
    \begin{bmatrix}
    k_{11} & k_{12}\\
    k_{21} & k_{22}
    \end{bmatrix}
\]

thì cùng bốn trọng số này được dùng ở tất cả các vị trí mà kernel trượt qua.

\begin{ghinho}
Weight sharing giúp CNN tiếp tục giảm số lượng tham số vì cùng một bộ trọng số được tái sử dụng trên toàn ảnh.
\end{ghinho}

\subsection{Tác dụng của hai ý tưởng này}

Hai ý tưởng này giúp CNN:

\begin{itemize}
    \item giảm rất mạnh số lượng tham số;
    \item giảm nguy cơ overfitting;
    \item tận dụng cấu trúc không gian của ảnh;
    \item phát hiện cùng một đặc trưng ở nhiều vị trí khác nhau.
\end{itemize}

% ==========================================================
\section{Kernel and Filter -- Kernel và bộ lọc}

\subsection{Kernel là gì?}

Kernel là một ma trận nhỏ chứa các trọng số được dùng để quét trên ảnh.

Ví dụ kernel kích thước $2 \times 2$:

\[
    K =
    \begin{bmatrix}
    k_{11} & k_{12}\\
    k_{21} & k_{22}
    \end{bmatrix}
\]

Trong thực tế, kernel có thể có kích thước:

\[
    3 \times 3,\quad 5 \times 5,\quad 7 \times 7,\quad 11 \times 11
\]

\begin{luuy}
Trong bài CNN1, kernel $2 \times 2$ được dùng để minh họa cho dễ tính tay. Trong mạng thực tế, kernel $3 \times 3$ rất phổ biến.
\end{luuy}

\subsection{Kernel học được từ đâu?}

Trong ví dụ, ta có thể thấy các giá trị kernel được cho sẵn. Tuy nhiên, trong mạng CNN thật, các giá trị này không phải do ta tự đặt thủ công.

Chúng được học thông qua quá trình huấn luyện:

\[
    \text{forward}
    \rightarrow
    \text{loss}
    \rightarrow
    \text{backpropagation}
    \rightarrow
    \text{update}
\]

\begin{ghinho}
Kernel cũng là tham số của mạng. Nó được học từ dữ liệu giống như trọng số trong mạng nơ-ron thông thường.
\end{ghinho}

% ==========================================================
\section{Convolution Forward Pass -- Tính xuôi trong lớp tích chập}

\subsection{Ý tưởng trượt kernel}

Để thực hiện convolution, ta làm như sau:

\begin{enumerate}[label=\arabic*.]
    \item Đặt kernel lên một vùng nhỏ của ảnh.
    \item Nhân từng phần tử tương ứng giữa kernel và vùng ảnh.
    \item Cộng tất cả các tích lại.
    \item Cộng thêm bias.
    \item Đưa qua hàm kích hoạt.
    \item Ghi kết quả vào vị trí tương ứng trong activation map.
    \item Trượt kernel sang vị trí tiếp theo và lặp lại.
\end{enumerate}

\subsection{Công thức tại vị trí đầu tiên}

Giả sử input là ma trận $3 \times 3$:

\[
    X =
    \begin{bmatrix}
    x_{11} & x_{12} & x_{13}\\
    x_{21} & x_{22} & x_{23}\\
    x_{31} & x_{32} & x_{33}
    \end{bmatrix}
\]

Kernel là:

\[
    K =
    \begin{bmatrix}
    k_{11} & k_{12}\\
    k_{21} & k_{22}
    \end{bmatrix}
\]

Khi đặt kernel ở góc trên bên trái, vùng ảnh tương ứng là:

\[
    \begin{bmatrix}
    x_{11} & x_{12}\\
    x_{21} & x_{22}
    \end{bmatrix}
\]

Kết quả tại vị trí đầu tiên là:

\[
    a_{11}
    =
    f(
    k_{11}x_{11}
    +
    k_{12}x_{12}
    +
    k_{21}x_{21}
    +
    k_{22}x_{22}
    +
    b
    )
\]

Trong đó:

\begin{itemize}
    \item $b$ là bias;
    \item $f$ là hàm kích hoạt;
    \item $a_{11}$ là phần tử đầu tiên của activation map.
\end{itemize}

\subsection{Các vị trí còn lại}

Với input $3 \times 3$, kernel $2 \times 2$, stride bằng 1 và không padding, output sẽ có kích thước:

\[
    2 \times 2
\]

Activation map là:

\[
    A =
    \begin{bmatrix}
    a_{11} & a_{12}\\
    a_{21} & a_{22}
    \end{bmatrix}
\]

Trong đó:

\[
    a_{12}
    =
    f(
    k_{11}x_{12}
    +
    k_{12}x_{13}
    +
    k_{21}x_{22}
    +
    k_{22}x_{23}
    +
    b
    )
\]

\[
    a_{21}
    =
    f(
    k_{11}x_{21}
    +
    k_{12}x_{22}
    +
    k_{21}x_{31}
    +
    k_{22}x_{32}
    +
    b
    )
\]

\[
    a_{22}
    =
    f(
    k_{11}x_{22}
    +
    k_{12}x_{23}
    +
    k_{21}x_{32}
    +
    k_{22}x_{33}
    +
    b
    )
\]

\begin{ghinho}
Mỗi phần tử trong activation map tương ứng với một vị trí đặt kernel trên ảnh.
\end{ghinho}

% ==========================================================
\section{Activation Map -- Bản đồ kích hoạt}

\subsection{Activation map là gì?}

Activation map là kết quả thu được sau khi một kernel quét qua toàn bộ ảnh.

\[
    \text{Input image}
    \xrightarrow{\text{kernel}}
    \text{activation map}
\]

Nếu một kernel phát hiện một loại đặc trưng nào đó, activation map cho biết đặc trưng đó xuất hiện mạnh ở vị trí nào.

\subsection{Một kernel tạo ra một activation map}

Nếu ta dùng một kernel, ta thu được một activation map.

Nếu ta dùng hai kernel, ta thu được hai activation map.

Tổng quát:

\[
    F \text{ kernels}
    \rightarrow
    F \text{ activation maps}
\]

\begin{vidu}
Một kernel có thể học để phát hiện cạnh ngang. Một kernel khác có thể học để phát hiện cạnh dọc. Mỗi kernel tạo ra một activation map riêng.
\end{vidu}

\subsection{Kích thước activation map}

Khi một kernel trượt qua input, tại mỗi vị trí nó tạo ra một giá trị output. Tập hợp tất cả các giá trị output này tạo thành một \textbf{activation map}.

Nói cách khác:

\[
    \text{mỗi vị trí đặt kernel}
    \Rightarrow
    \text{một giá trị trong activation map}
\]

Vì vậy, kích thước của activation map phụ thuộc vào:

\begin{itemize}
    \item kích thước input;
    \item kích thước kernel;
    \item stride;
    \item padding.
\end{itemize}

\subsubsection{Trường hợp đơn giản: stride bằng 1 và không padding}

Nếu input có kích thước:

\[
    H \times W
\]

trong đó:

\[
    H = \text{chiều cao input}
\]

\[
    W = \text{chiều rộng input}
\]

Kernel có kích thước:

\[
    K_h \times K_w
\]

trong đó:

\[
    K_h = \text{chiều cao kernel}
\]

\[
    K_w = \text{chiều rộng kernel}
\]

Nếu stride bằng 1 và không padding, thì output có kích thước:

\[
    (H-K_h+1) \times (W-K_w+1)
\]

\subsubsection{Vì sao có công thức này?}

Giả sử input có kích thước $3 \times 3$:

\[
I=
\begin{bmatrix}
1 & 2 & 3\\
4 & 5 & 6\\
7 & 8 & 9
\end{bmatrix}
\]

và kernel có kích thước $2 \times 2$.

Kernel $2 \times 2$ có thể đặt vào bốn vị trí khác nhau trên input.

Vị trí thứ nhất, góc trên bên trái:

\[
\begin{bmatrix}
1 & 2\\
4 & 5
\end{bmatrix}
\]

Vị trí thứ hai, góc trên bên phải:

\[
\begin{bmatrix}
2 & 3\\
5 & 6
\end{bmatrix}
\]

Vị trí thứ ba, góc dưới bên trái:

\[
\begin{bmatrix}
4 & 5\\
7 & 8
\end{bmatrix}
\]

Vị trí thứ tư, góc dưới bên phải:

\[
\begin{bmatrix}
5 & 6\\
8 & 9
\end{bmatrix}
\]

Như vậy, theo chiều cao, kernel đặt được:

\[
    3-2+1=2
\]

vị trí.

Theo chiều rộng, kernel cũng đặt được:

\[
    3-2+1=2
\]

vị trí.

Do đó activation map có kích thước:

\[
    2 \times 2
\]

Áp dụng công thức:

\[
    (H-K_h+1)\times(W-K_w+1)
\]

với:

\[
    H=3,\quad W=3,\quad K_h=2,\quad K_w=2
\]

ta có:

\[
    (3-2+1)\times(3-2+1)=2\times2
\]

\subsubsection{Ví dụ tính activation map cụ thể}

Giả sử input là:

\[
I=
\begin{bmatrix}
1 & 2 & 3\\
4 & 5 & 6\\
7 & 8 & 9
\end{bmatrix}
\]

Kernel là:

\[
K=
\begin{bmatrix}
1 & 0\\
0 & -1
\end{bmatrix}
\]

Bias:

\[
b=0
\]

Ta không dùng hàm kích hoạt để tính cho đơn giản.

Tại vị trí thứ nhất:

\[
\begin{bmatrix}
1 & 2\\
4 & 5
\end{bmatrix}
\]

ta tính:

\[
1\cdot1+2\cdot0+4\cdot0+5\cdot(-1)
\]

\[
=1-5=-4
\]

Tại vị trí thứ hai:

\[
\begin{bmatrix}
2 & 3\\
5 & 6
\end{bmatrix}
\]

ta tính:

\[
2\cdot1+3\cdot0+5\cdot0+6\cdot(-1)
\]

\[
=2-6=-4
\]

Tại vị trí thứ ba:

\[
\begin{bmatrix}
4 & 5\\
7 & 8
\end{bmatrix}
\]

ta tính:

\[
4\cdot1+5\cdot0+7\cdot0+8\cdot(-1)
\]

\[
=4-8=-4
\]

Tại vị trí thứ tư:

\[
\begin{bmatrix}
5 & 6\\
8 & 9
\end{bmatrix}
\]

ta tính:

\[
5\cdot1+6\cdot0+8\cdot0+9\cdot(-1)
\]

\[
=5-9=-4
\]

Vậy activation map thu được là:

\[
\begin{bmatrix}
-4 & -4\\
-4 & -4
\end{bmatrix}
\]

Nếu sau đó dùng hàm kích hoạt ReLU:

\[
    \operatorname{ReLU}(x)=\max(0,x)
\]

thì activation map sau ReLU là:

\[
\begin{bmatrix}
0 & 0\\
0 & 0
\end{bmatrix}
\]

\subsubsection{Stride là gì?}

\textbf{Stride} là số bước mà kernel dịch chuyển sau mỗi lần tính.

Nếu:

\[
    S=1
\]

thì kernel trượt từng ô một.

Nếu:

\[
    S=2
\]

thì kernel nhảy cách hai ô một lần.

Ví dụ, với input $4 \times 4$:

\[
I=
\begin{bmatrix}
1 & 2 & 3 & 4\\
5 & 6 & 7 & 8\\
9 & 10 & 11 & 12\\
13 & 14 & 15 & 16
\end{bmatrix}
\]

và kernel $2 \times 2$.

Nếu stride bằng 1, kernel có thể đi qua nhiều vị trí hơn, nên activation map lớn hơn.

Nếu stride bằng 2, kernel nhảy xa hơn, số vị trí đặt kernel ít hơn, nên activation map nhỏ hơn.

\[
    \text{stride lớn}
    \Rightarrow
    \text{kernel nhảy xa}
    \Rightarrow
    \text{activation map nhỏ hơn}
\]

\subsubsection{Padding là gì?}

\textbf{Padding} là việc thêm các giá trị vào viền xung quanh input. Thông thường, giá trị được thêm vào là số 0, nên còn gọi là \textbf{zero padding}.

Ví dụ input $3 \times 3$:

\[
I=
\begin{bmatrix}
1 & 2 & 3\\
4 & 5 & 6\\
7 & 8 & 9
\end{bmatrix}
\]

Nếu padding bằng 1, ta thêm một lớp số 0 xung quanh input:

\[
I_{pad}=
\begin{bmatrix}
0 & 0 & 0 & 0 & 0\\
0 & 1 & 2 & 3 & 0\\
0 & 4 & 5 & 6 & 0\\
0 & 7 & 8 & 9 & 0\\
0 & 0 & 0 & 0 & 0
\end{bmatrix}
\]

Khi đó input từ kích thước:

\[
    3 \times 3
\]

trở thành:

\[
    5 \times 5
\]

Padding có hai tác dụng quan trọng.

Thứ nhất, padding giúp activation map không bị nhỏ quá nhanh sau mỗi lớp convolution.

Thứ hai, padding giúp các pixel ở rìa ảnh được kernel xét đến nhiều hơn. Nếu không padding, các pixel ở góc và cạnh ảnh thường được dùng ít hơn các pixel ở giữa ảnh.

\subsubsection{Công thức tổng quát có stride và padding}

Với input kích thước:

\[
    H \times W
\]

kernel kích thước:

\[
    K_h \times K_w
\]

padding:

\[
    P
\]

stride:

\[
    S
\]

thì kích thước activation map là:

\[
    H_{out}
    =
    \left\lfloor
    \frac{H+2P-K_h}{S}
    \right\rfloor
    +1
\]

\[
    W_{out}
    =
    \left\lfloor
    \frac{W+2P-K_w}{S}
    \right\rfloor
    +1
\]

Do đó:

\[
    \text{output size}
    =
    H_{out}
    \times
    W_{out}
\]

Trong đó ký hiệu:

\[
    \left\lfloor x \right\rfloor
\]

nghĩa là lấy phần nguyên xuống.

\subsubsection{Ví dụ có padding}

Giả sử input có kích thước:

\[
    5 \times 5
\]

kernel có kích thước:

\[
    3 \times 3
\]

stride:

\[
    S=1
\]

không padding:

\[
    P=0
\]

Khi đó:

\[
    H_{out}
    =
    \frac{5+2\cdot0-3}{1}+1
    =
    3
\]

\[
    W_{out}
    =
    \frac{5+2\cdot0-3}{1}+1
    =
    3
\]

Vậy output có kích thước:

\[
    3 \times 3
\]

Nếu vẫn dùng input $5 \times 5$, kernel $3 \times 3$, stride $1$, nhưng thêm padding:

\[
    P=1
\]

thì:

\[
    H_{out}
    =
    \frac{5+2\cdot1-3}{1}+1
    =
    5
\]

\[
    W_{out}
    =
    \frac{5+2\cdot1-3}{1}+1
    =
    5
\]

Vậy output có kích thước:

\[
    5 \times 5
\]

Như vậy, padding giúp giữ nguyên kích thước không gian của activation map so với input ban đầu.

\subsubsection{Tóm tắt}

Stride là bước nhảy của kernel:

\[
    \text{stride càng lớn}
    \Rightarrow
    \text{activation map càng nhỏ}
\]

Padding là phần viền được thêm vào input:

\[
    \text{padding càng lớn}
    \Rightarrow
    \text{activation map càng lớn}
\]

Trong trường hợp đơn giản nhất, stride bằng 1 và không padding:

\[
    H_{out}=H-K_h+1
\]

\[
    W_{out}=W-K_w+1
\]

Còn trong trường hợp tổng quát:

\[
    H_{out}
    =
    \left\lfloor
    \frac{H+2P-K_h}{S}
    \right\rfloor
    +1
\]

\[
    W_{out}
    =
    \left\lfloor
    \frac{W+2P-K_w}{S}
    \right\rfloor
    +1
\]

% ==========================================================
\section{Simple X/O Classification Example -- Ví dụ phân loại chữ X/O}

\subsection{Mô tả bài toán}

Trong video, ta khảo sát một ví dụ rất nhỏ:

\[
    \text{input là ảnh } 3\times3
\]

và chỉ có hai class:

\[
    \text{chữ X}
    \quad \text{hoặc} \quad
    \text{chữ O}
\]

Nhiệm vụ của mạng là xác định ảnh đầu vào thuộc class nào.

\subsection{Biểu diễn chữ X và chữ O}

Để dễ tính, ta có thể biểu diễn pixel trắng là 1 và pixel đen là 0.

Chữ X có thể biểu diễn như sau:

\[
    X_{\text{letter X}}
    =
    \begin{bmatrix}
    1 & 0 & 1\\
    0 & 1 & 0\\
    1 & 0 & 1
    \end{bmatrix}
\]

Chữ O có thể biểu diễn như sau:

\[
    X_{\text{letter O}}
    =
    \begin{bmatrix}
    1 & 1 & 1\\
    1 & 0 & 1\\
    1 & 1 & 1
    \end{bmatrix}
\]

\begin{luuy}
Trong slide, pixel trắng có thể được biểu diễn bằng 255 và pixel đen bằng 0. Ở đây ta dùng 1 và 0 để việc tính tay đơn giản hơn.
\end{luuy}

\subsection{One-hot output}

Vì có hai class, output có hai phần tử.

Ta quy ước:

\[
    \text{X}
    =
    \begin{bmatrix}
    1\\
    0
    \end{bmatrix}
\]

\[
    \text{O}
    =
    \begin{bmatrix}
    0\\
    1
    \end{bmatrix}
\]

Nếu đưa chữ X vào, ta mong muốn output là:

\[
    \begin{bmatrix}
    1\\
    0
    \end{bmatrix}
\]

Nếu đưa chữ O vào, ta mong muốn output là:

\[
    \begin{bmatrix}
    0\\
    1
    \end{bmatrix}
\]

% ==========================================================
\section{Convolution Example for Letter X -- Ví dụ tích chập cho chữ X}

\subsection{Hai kernel minh họa}

Giả sử ta dùng hai kernel đơn giản.

Kernel thứ nhất:

\[
    K_1 =
    \begin{bmatrix}
    1 & 0\\
    0 & 1
    \end{bmatrix}
\]

Kernel thứ hai:

\[
    K_2 =
    \begin{bmatrix}
    0 & 1\\
    1 & 0
    \end{bmatrix}
\]

Giả sử:

\[
    b=0
\]

và hàm kích hoạt là hàm đồng nhất:

\[
    f(u)=u
\]

\subsection{Input chữ X}

\[
    X =
    \begin{bmatrix}
    1 & 0 & 1\\
    0 & 1 & 0\\
    1 & 0 & 1
    \end{bmatrix}
\]

\subsection{Activation map với kernel thứ nhất}

Tại vị trí trên trái:

\[
    a_{11}^{(1)}
    =
    1\cdot1
    +
    0\cdot0
    +
    0\cdot0
    +
    1\cdot1
    =
    2
\]

Tại vị trí trên phải:

\[
    a_{12}^{(1)}
    =
    1\cdot0
    +
    0\cdot1
    +
    0\cdot1
    +
    1\cdot0
    =
    0
\]

Tại vị trí dưới trái:

\[
    a_{21}^{(1)}
    =
    1\cdot0
    +
    0\cdot1
    +
    0\cdot1
    +
    1\cdot0
    =
    0
\]

Tại vị trí dưới phải:

\[
    a_{22}^{(1)}
    =
    1\cdot1
    +
    0\cdot0
    +
    0\cdot0
    +
    1\cdot1
    =
    2
\]

Vậy activation map thứ nhất là:

\[
    A_1 =
    \begin{bmatrix}
    2 & 0\\
    0 & 2
    \end{bmatrix}
\]

\subsection{Activation map với kernel thứ hai}

Tương tự, với kernel:

\[
    K_2 =
    \begin{bmatrix}
    0 & 1\\
    1 & 0
    \end{bmatrix}
\]

ta thu được:

\[
    A_2 =
    \begin{bmatrix}
    0 & 2\\
    2 & 0
    \end{bmatrix}
\]

\begin{ghinho}
Với chữ X, hai kernel trên lần lượt phản ứng mạnh với hai hướng đường chéo khác nhau.
\end{ghinho}

% ==========================================================
\section{Flatten and Fully Connected Layer -- Flatten và lớp kết nối đầy đủ}

\subsection{Flatten là gì?}

Sau khi có các activation map, ta thường biến chúng thành một vector để đưa vào phần fully connected.

Thao tác này gọi là:

\[
    \text{flatten}
\]

Ví dụ:

\[
    A_1 =
    \begin{bmatrix}
    2 & 0\\
    0 & 2
    \end{bmatrix}
\]

\[
    A_2 =
    \begin{bmatrix}
    0 & 2\\
    2 & 0
    \end{bmatrix}
\]

Nếu đọc theo từng hàng, ta có vector:

\[
    v =
    \begin{bmatrix}
    2\\
    0\\
    0\\
    2\\
    0\\
    2\\
    2\\
    0
    \end{bmatrix}
\]

\subsection{Fully connected layer}

Vector sau flatten được đưa vào fully connected layer để phân loại.

Nếu có hai output neuron:

\[
    z_1 =
    \sum_{i=1}^{m} w_{1i}v_i+b_1
\]

\[
    z_2 =
    \sum_{i=1}^{m} w_{2i}v_i+b_2
\]

Trong bài toán X/O:

\[
    z_1
    \rightarrow
    \text{điểm cho chữ X}
\]

\[
    z_2
    \rightarrow
    \text{điểm cho chữ O}
\]

Sau đó ta có thể dùng hàm phù hợp, ví dụ softmax, để chuyển thành xác suất phân loại.

\begin{ghinho}
Phần convolution trích xuất đặc trưng. Phần fully connected dùng các đặc trưng đó để quyết định class.
\end{ghinho}

% ==========================================================
\section{Where Do CNN Parameters Come From? -- Các tham số trong CNN đến từ đâu?}

\subsection{Kernel, bias và trọng số}

Trong ví dụ, kernel, bias và trọng số fully connected có thể được cho sẵn để ta tính tay.

Tuy nhiên, trong mạng thật, các tham số này được học từ dữ liệu.

Các tham số cần học gồm:

\begin{itemize}
    \item các giá trị trong kernel;
    \item bias của convolution layer;
    \item trọng số fully connected;
    \item bias fully connected.
\end{itemize}

\subsection{Huấn luyện CNN}

Về nguyên tắc, CNN vẫn được huấn luyện theo quy trình quen thuộc:

\[
    \text{forward}
    \rightarrow
    \text{loss}
    \rightarrow
    \text{backpropagation}
    \rightarrow
    \text{update parameters}
\]

Tuy nhiên, vì CNN có weight sharing và convolution layer, backpropagation trong CNN sẽ hơi khác so với mạng fully connected thông thường.

\begin{luuy}
Bài CNN1 chủ yếu tập trung vào cách tính forward trong convolution layer. Phần backpropagation của CNN sẽ cần được học kỹ hơn ở các bài sau.
\end{luuy}

% ==========================================================
\section{Why CNN Works Well for Images -- Vì sao CNN phù hợp với ảnh?}

\subsection{Ảnh có tính cục bộ}

Trong ảnh, nhiều đặc trưng quan trọng xuất hiện trong những vùng nhỏ.

Ví dụ:

\begin{itemize}
    \item cạnh ngang;
    \item cạnh dọc;
    \item góc;
    \item đường cong;
    \item họa tiết;
    \item mắt, mũi, miệng trong ảnh khuôn mặt.
\end{itemize}

Vì vậy, một kernel nhỏ vẫn có thể phát hiện được đặc trưng quan trọng.

\subsection{Đặc trưng có thể xuất hiện ở nhiều vị trí}

Một đặc trưng như cạnh, góc hoặc đường cong có thể xuất hiện ở bất kỳ vị trí nào trong ảnh.

Nếu một kernel học được cách phát hiện cạnh dọc, ta có thể dùng kernel đó quét trên toàn ảnh để tìm cạnh dọc ở nhiều vị trí.

Đây là lý do weight sharing phù hợp với ảnh.

\subsection{CNN giảm overfitting}

Vì CNN giảm số lượng tham số so với fully connected network, nên CNN thường giảm nguy cơ overfitting khi xử lý ảnh.

\begin{ghinho}
CNN hiệu quả vì nó tận dụng hai đặc điểm quan trọng của ảnh: thông tin cục bộ và sự lặp lại của đặc trưng theo vị trí.
\end{ghinho}

% ==========================================================
\section{ANN vs CNN -- So sánh ANN thông thường và CNN}

\begin{center}
\begin{tabular}{p{0.27\textwidth}p{0.31\textwidth}p{0.31\textwidth}}
\toprule
Tiêu chí & ANN fully connected & CNN \\
\midrule
Cách xử lý ảnh & Flatten ảnh thành vector & Giữ cấu trúc không gian của ảnh \\
Kết nối & Mỗi neuron nhận toàn bộ input & Mỗi neuron nhìn một vùng cục bộ \\
Trọng số & Mỗi neuron có bộ trọng số riêng & Kernel được dùng chung nhiều vị trí \\
Số tham số & Rất lớn với ảnh lớn & Nhỏ hơn nhiều \\
Phù hợp với ảnh & Kém hơn khi ảnh lớn/phức tạp & Rất phù hợp \\
Nguy cơ overfitting & Cao nếu dữ liệu không đủ & Thấp hơn do ít tham số hơn \\
Đặc trưng học được & Ít khai thác cấu trúc không gian & Học cạnh, góc, pattern cục bộ \\
\bottomrule
\end{tabular}
\end{center}

% ==========================================================
\section{Technical Glossary -- Bảng thuật ngữ chuyên ngành}

\small
\begin{longtable}{p{0.27\textwidth}p{0.48\textwidth}p{0.18\textwidth}}
\toprule
\textbf{Thuật ngữ} & \textbf{Giải thích} & \textbf{Ví dụ} \\
\midrule
\endfirsthead

\toprule
\textbf{Thuật ngữ} & \textbf{Giải thích} & \textbf{Ví dụ} \\
\midrule
\endhead

CNN & Convolutional Neural Network, mạng nơ-ron tích chập. & Mạng phân loại ảnh. \\

Convolution & Phép tích chập, quét kernel trên ảnh để tạo đặc trưng. & Kernel $3\times3$. \\

Kernel & Ma trận trọng số nhỏ dùng trong convolution. & $2\times2$, $3\times3$. \\

Filter & Bộ lọc; thường được dùng gần nghĩa với kernel. & Bộ lọc cạnh. \\

Activation map & Ma trận kết quả sau khi kernel quét trên ảnh. & Output $2\times2$. \\

Feature map & Cách gọi khác của activation map. & Bản đồ đặc trưng. \\

Local connectivity & Kết nối cục bộ, neuron chỉ nhìn một vùng nhỏ. & Kernel nhìn vùng $3\times3$. \\

Weight sharing & Trọng số dùng chung ở nhiều vị trí. & Một kernel quét toàn ảnh. \\

Fully connected layer & Lớp kết nối đầy đủ như ANN truyền thống. & Lớp phân loại cuối. \\

Flatten & Biến ma trận hoặc tensor thành vector. & $2\times2$ thành vector 4 phần tử. \\

Pixel & Điểm ảnh. & Giá trị từ 0 đến 255. \\

Channel & Kênh của ảnh. & R, G, B. \\

RGB & Ảnh màu có 3 kênh đỏ, xanh lá, xanh dương. & $H\times W\times3$. \\

Pooling & Lớp giảm kích thước đặc trưng. & Max pooling. \\

Stride & Bước trượt của kernel. & Stride bằng 1. \\

Padding & Thêm viền quanh ảnh để điều chỉnh kích thước output. & Zero padding. \\

Overfitting & Mô hình quá khớp dữ liệu huấn luyện. & Quá nhiều tham số. \\

Backpropagation & Lan truyền ngược để cập nhật tham số. & Học kernel. \\

\bottomrule
\end{longtable}
\normalsize

% ==========================================================
\section{Key Concepts Summary -- Tóm tắt kiến thức trọng tâm}

\subsection{Các ý chính cần nhớ}

\begin{enumerate}
    \item CNN là mạng nơ-ron tích chập, rất quan trọng trong xử lý ảnh.
    \item CNN thường gồm phần convolution/pooling và phần fully connected.
    \item Mạng fully connected gặp vấn đề lớn khi xử lý ảnh vì số lượng tham số tăng rất nhanh.
    \item Ảnh $32\times32\times3$ có 3072 giá trị.
    \item Ảnh $200\times200\times3$ có 120000 giá trị.
    \item Một neuron fully connected nhận toàn bộ input nên cần rất nhiều trọng số.
    \item CNN giảm số tham số bằng local connectivity và weight sharing.
    \item Kernel là một ma trận trọng số nhỏ trượt trên ảnh.
    \item Mỗi lần đặt kernel lên một vùng ảnh, ta nhân tương ứng, cộng lại, cộng bias và đưa qua hàm kích hoạt.
    \item Kết quả khi một kernel quét toàn ảnh là một activation map.
    \item Nhiều kernel tạo ra nhiều activation map.
    \item Sau convolution, activation map có thể được flatten rồi đưa vào fully connected layer.
    \item Kernel và trọng số fully connected được học bằng backpropagation.
    \item CNN hiệu quả vì ảnh có tính cục bộ và đặc trưng có thể xuất hiện ở nhiều vị trí.
\end{enumerate}

\subsection{Công thức quan trọng}

\subsection*{Số giá trị của ảnh RGB}

\[
    H \cdot W \cdot 3
\]

\subsection*{Số trọng số của một neuron fully connected}

\[
    H \cdot W \cdot C
\]

\subsection*{Số trọng số của một kernel}

\[
    K_h \cdot K_w \cdot C
\]

\subsection*{Convolution với kernel $2\times2$}

\[
    a_{ij}
    =
    f(
    k_{11}x_{ij}
    +
    k_{12}x_{i,j+1}
    +
    k_{21}x_{i+1,j}
    +
    k_{22}x_{i+1,j+1}
    +
    b
    )
\]

\subsection*{Kích thước output không padding, stride 1}

\[
    (H-K_h+1)\times(W-K_w+1)
\]

% ==========================================================
\section{Review Questions -- Câu hỏi ôn tập}

\begin{enumerate}
    \item CNN là viết tắt của cụm từ gì?
    \item CNN được thiết kế ban đầu chủ yếu cho loại dữ liệu nào?
    \item Một ảnh màu $32\times32$ có bao nhiêu giá trị input?
    \item Một ảnh màu $200\times200$ có bao nhiêu giá trị input?
    \item Vì sao mạng fully connected dễ có quá nhiều tham số khi xử lý ảnh?
    \item Local connectivity là gì?
    \item Weight sharing là gì?
    \item Kernel là gì?
    \item Activation map là gì?
    \item Một kernel tạo ra bao nhiêu activation map?
    \item Nếu có 10 kernel thì có bao nhiêu activation map?
    \item Với input $3\times3$ và kernel $2\times2$, stride 1, không padding, output có kích thước bao nhiêu?
    \item Flatten là gì?
    \item Vì sao CNN thường giảm overfitting tốt hơn fully connected layer khi xử lý ảnh?
    \item Các kernel trong CNN thật được lấy từ đâu?
\end{enumerate}

% ==========================================================
\section{Practice Exercises -- Bài tập vận dụng}

\subsection*{Bài tập 1: Tính số input}
\addcontentsline{toc}{section}{Bài tập 1: Tính số input}

Một ảnh màu có kích thước:

\[
    64\times64\times3
\]

Hỏi ảnh này có bao nhiêu giá trị pixel?

\subsection*{Lời giải gợi ý}

\[
    64\cdot64\cdot3=12288
\]

Vậy ảnh có 12288 giá trị.

\subsection*{Bài tập 2: So sánh số trọng số}
\addcontentsline{toc}{section}{Bài tập 2: So sánh số trọng số}

Một ảnh màu kích thước:

\[
    100\times100\times3
\]

\begin{enumerate}
    \item Một neuron fully connected cần bao nhiêu trọng số?
    \item Một kernel $3\times3$ cần bao nhiêu trọng số?
\end{enumerate}

\subsection*{Lời giải gợi ý}

Neuron fully connected:

\[
    100\cdot100\cdot3=30000
\]

Kernel $3\times3$:

\[
    3\cdot3\cdot3=27
\]

\subsection*{Bài tập 3: Tính convolution}
\addcontentsline{toc}{section}{Bài tập 3: Tính convolution}

Cho input:

\[
    X =
    \begin{bmatrix}
    1 & 2 & 0\\
    0 & 1 & 3\\
    2 & 1 & 0
    \end{bmatrix}
\]

Kernel:

\[
    K =
    \begin{bmatrix}
    1 & 0\\
    0 & 1
    \end{bmatrix}
\]

Giả sử $b=0$ và $f(u)=u$. Tính activation map.

\subsection*{Lời giải gợi ý}

\[
    a_{11}=1\cdot1+0\cdot2+0\cdot0+1\cdot1=2
\]

\[
    a_{12}=1\cdot2+0\cdot0+0\cdot1+1\cdot3=5
\]

\[
    a_{21}=1\cdot0+0\cdot1+0\cdot2+1\cdot1=1
\]

\[
    a_{22}=1\cdot1+0\cdot3+0\cdot1+1\cdot0=1
\]

Vậy:

\[
    A =
    \begin{bmatrix}
    2 & 5\\
    1 & 1
    \end{bmatrix}
\]

\subsection*{Bài tập 4: Kích thước output}
\addcontentsline{toc}{section}{Bài tập 4: Kích thước output}

Input có kích thước:

\[
    10\times10
\]

Kernel có kích thước:

\[
    3\times3
\]

Stride bằng 1, không padding. Hỏi activation map có kích thước bao nhiêu?

\subsection*{Lời giải gợi ý}

\[
    (10-3+1)\times(10-3+1)=8\times8
\]

\subsection*{Bài tập 5: Giải thích trực giác}
\addcontentsline{toc}{section}{Bài tập 5: Giải thích trực giác}

Giải thích ngắn gọn vì sao CNN dùng local connectivity và weight sharing lại phù hợp với ảnh.

\subsection*{Lời giải gợi ý}

Ảnh có cấu trúc không gian. Các đặc trưng như cạnh, góc, đường cong thường xuất hiện trong những vùng nhỏ. Vì vậy, local connectivity giúp kernel tập trung vào vùng cục bộ. Ngoài ra, cùng một đặc trưng có thể xuất hiện ở nhiều vị trí khác nhau trong ảnh, nên weight sharing cho phép cùng một kernel phát hiện đặc trưng đó ở mọi nơi. Hai cơ chế này giúp giảm số tham số và giảm nguy cơ overfitting.

% ==========================================================
\section*{Conclusion -- Kết luận}
\addcontentsline{toc}{section}{Kết luận}

CNN1 giới thiệu mạng nơ-ron tích chập, một kiến trúc rất quan trọng trong deep learning, đặc biệt với các bài toán hình ảnh. Khác với mạng fully connected truyền thống, CNN không biến ảnh thành một vector rồi kết nối mọi pixel với mọi neuron ngay từ đầu. Thay vào đó, CNN dùng các kernel nhỏ trượt trên ảnh để trích xuất đặc trưng cục bộ.

Hai ý tưởng cốt lõi của CNN là \textit{local connectivity} và \textit{weight sharing}. Local connectivity giúp mỗi kernel chỉ nhìn một vùng nhỏ của ảnh, còn weight sharing giúp cùng một bộ trọng số được dùng lại ở nhiều vị trí. Nhờ đó, CNN giảm mạnh số lượng tham số, khai thác tốt cấu trúc không gian của ảnh và giảm nguy cơ overfitting.

Bài học cũng minh họa cách tính convolution forward bằng kernel $2\times2$ trên ảnh $3\times3$, cách tạo activation map, cách flatten kết quả và đưa vào fully connected layer để phân loại. Các tham số như kernel, bias và trọng số fully connected không phải tự chọn bằng tay trong mạng thật, mà được học thông qua quá trình huấn luyện và backpropagation.

\chapter{CNN2: Multichannel Convolution, Pooling, and Padding -- Tích chập nhiều kênh, pooling và padding}
% ==========================================================
\section*{Learning Objectives -- Mục tiêu bài học}
\addcontentsline{toc}{section}{Mục tiêu bài học}

Bài học này tiếp tục phần CNN, tập trung vào cách xử lý ảnh màu, convolution nhiều kênh, pooling layer, việc lặp lại nhiều lớp tích chập và kỹ thuật zero padding.

Sau khi học xong, người học cần nắm được:

\begin{enumerate}[label=\arabic*.]
    \item Hiểu sự khác nhau giữa convolution trên ảnh xám và ảnh màu.
    \item Biết vì sao kernel phải có số layer tương ứng với số channel của input.
    \item Hiểu cách tính convolution khi input có nhiều channel.
    \item Biết nhiều kernel tạo ra nhiều activation map.
    \item Hiểu pooling layer là gì.
    \item Phân biệt max pooling và average pooling.
    \item Hiểu pooling là một dạng downsampling.
    \item Biết mục tiêu chính của pooling là giảm kích thước dữ liệu và giảm số tham số.
    \item Hiểu cách các lớp convolution và pooling được lặp lại trong CNN.
    \item Hiểu khi input của convolution layer có nhiều activation map thì kernel cũng phải có cùng số layer.
    \item Hiểu ý nghĩa của flatten trước khi đưa vào fully connected layer.
    \item Hiểu zero padding là gì và vì sao cần dùng padding.
\end{enumerate}

% ==========================================================
\section{Review of Simple Convolution -- Ôn lại phép tích chập đơn giản}

\subsection{Trường hợp ảnh xám một kênh}

Trong bài trước, ta đã khảo sát trường hợp đơn giản nhất: input là một ảnh xám rất nhỏ.

Ví dụ input là ma trận $3\times3$:

\[
    X =
    \begin{bmatrix}
    x_{11} & x_{12} & x_{13}\\
    x_{21} & x_{22} & x_{23}\\
    x_{31} & x_{32} & x_{33}
    \end{bmatrix}
\]

Vì ảnh xám chỉ có một kênh, nên input chỉ là một ma trận 2D.

Kernel cũng là một ma trận 2D, ví dụ kernel $2\times2$:

\[
    K =
    \begin{bmatrix}
    k_{11} & k_{12}\\
    k_{21} & k_{22}
    \end{bmatrix}
\]

Khi đặt kernel lên một vùng ảnh, ta nhân từng phần tử tương ứng, cộng lại, cộng bias, rồi đưa qua hàm kích hoạt.

\[
    a_{11}
    =
    f(
    k_{11}x_{11}
    +
    k_{12}x_{12}
    +
    k_{21}x_{21}
    +
    k_{22}x_{22}
    +
    b
    )
\]

\begin{ghinho}
Với ảnh một kênh, kernel là một ma trận 2D. Với ảnh nhiều kênh, kernel không còn chỉ là một ma trận 2D nữa.
\end{ghinho}

\subsection{Kích thước kernel}

Trong ví dụ đơn giản, ta dùng kernel $2\times2$ để dễ tính tay.

Trong thực tế, kernel có thể có nhiều kích thước khác nhau:

\[
    2\times2,\quad 3\times3,\quad 5\times5,\quad 7\times7,\quad 11\times11
\]

Kernel $3\times3$ là một lựa chọn rất phổ biến trong nhiều mạng CNN hiện đại.

% ==========================================================
\section{Color Image Convolution -- Tích chập trên ảnh màu}

\subsection{Ảnh màu có nhiều channel}

Ảnh màu RGB có ba kênh màu:

\[
    R,\quad G,\quad B
\]

Có thể hiểu mỗi kênh màu là một ma trận riêng.

Vì vậy, một ảnh màu có kích thước:

\[
    H \times W \times 3
\]

Trong đó:

\begin{itemize}
    \item $H$ là chiều cao ảnh;
    \item $W$ là chiều rộng ảnh;
    \item $3$ là số channel màu RGB.
\end{itemize}

\begin{vidu}
Một ảnh màu $224\times224$ không chỉ là một ma trận $224\times224$, mà là một khối dữ liệu:

\[
    224 \times 224 \times 3
\]

gồm ba ma trận tương ứng với ba kênh màu.
\end{vidu}

\subsection{Kernel cũng phải có cùng số layer với input}

Nếu input có 3 channel, kernel cũng phải có 3 layer tương ứng.

Ví dụ, nếu input là ảnh RGB:

\[
    X \in \mathbb{R}^{H\times W\times 3}
\]

thì một kernel kích thước không gian $3\times3$ phải có dạng:

\[
    K \in \mathbb{R}^{3\times3\times3}
\]

Tức là kernel gồm 3 ma trận $3\times3$:

\[
    K_R,\quad K_G,\quad K_B
\]

Mỗi layer của kernel áp lên một channel tương ứng của ảnh.

\[
    K_R \text{ áp lên kênh } R
\]

\[
    K_G \text{ áp lên kênh } G
\]

\[
    K_B \text{ áp lên kênh } B
\]

\begin{ghinho}
Độ sâu của kernel phải bằng số channel của input. Nếu input có 3 channel thì kernel cũng phải có 3 layer.
\end{ghinho}

% ==========================================================
\section{Multi-channel Convolution Formula -- Công thức tích chập nhiều kênh}

\subsection{Ý tưởng tính toán}

Giả sử input có 3 channel:

\[
    X_R,\quad X_G,\quad X_B
\]

và kernel có 3 layer:

\[
    K_R,\quad K_G,\quad K_B
\]

Khi đặt kernel lên một vùng ảnh, ta làm như sau:

\begin{enumerate}[label=\arabic*.]
    \item Nhân tương ứng $K_R$ với vùng ảnh trên channel $R$.
    \item Nhân tương ứng $K_G$ với vùng ảnh trên channel $G$.
    \item Nhân tương ứng $K_B$ với vùng ảnh trên channel $B$.
    \item Cộng tất cả các tích của cả ba channel.
    \item Cộng bias.
    \item Đưa qua hàm kích hoạt.
\end{enumerate}

\subsection{Ví dụ kernel $3\times3\times3$}

Nếu kernel có kích thước:

\[
    3\times3\times3
\]

thì số phép nhân tại một vị trí là:

\[
    3 \cdot 3 \cdot 3 = 27
\]

Sau đó cộng 27 tích lại, cộng thêm bias và đưa qua hàm kích hoạt.

\[
    a_{ij}
    =
    f
    \left(
    \sum_{c=1}^{3}
    \sum_{u=1}^{3}
    \sum_{v=1}^{3}
    K_{u,v,c}
    X_{i+u-1,j+v-1,c}
    +
    b
    \right)
\]

Trong đó:

\begin{itemize}
    \item $c$ là chỉ số channel;
    \item $u,v$ là chỉ số trong kernel;
    \item $b$ là bias;
    \item $f$ là hàm kích hoạt;
    \item $a_{ij}$ là một giá trị trong activation map.
\end{itemize}

\begin{ghinho}
Trong convolution nhiều kênh, tại mỗi vị trí ta cộng tất cả các tích trên toàn bộ chiều rộng, chiều cao và chiều sâu của kernel.
\end{ghinho}

\subsection{Một kernel vẫn tạo ra một activation map}

Dù kernel có nhiều layer, ví dụ $3\times3\times3$, thì toàn bộ kernel đó vẫn tạo ra một activation map duy nhất.

\[
    \text{1 kernel } 3\times3\times3
    \rightarrow
    \text{1 activation map}
\]

Lý do là tại mỗi vị trí trượt, sau khi cộng tất cả các tích trên mọi channel, ta chỉ thu được một giá trị output.

% ==========================================================
\section{Multiple Kernels -- Nhiều kernel trong một lớp convolution}

\subsection{Một lớp convolution thường có nhiều kernel}

Trong thực tế, một convolution layer thường không chỉ có một kernel.

Ta có thể có:

\[
    K \text{ kernels}
\]

Mỗi kernel học một kiểu đặc trưng khác nhau.

Ví dụ:

\begin{itemize}
    \item một kernel học cạnh ngang;
    \item một kernel học cạnh dọc;
    \item một kernel học đường chéo;
    \item một kernel học vùng màu hoặc texture.
\end{itemize}

\subsection{Số kernel bằng số activation map đầu ra}

Nếu có $K$ kernel, ta thu được $K$ activation map.

\[
    K \text{ kernels}
    \rightarrow
    K \text{ activation maps}
\]

Nếu một lớp convolution có 64 kernel, output của lớp đó sẽ có 64 activation map.

\begin{ghinho}
Số activation map đầu ra của một convolution layer bằng số kernel của layer đó.
\end{ghinho}

\subsection{Output của convolution layer là một khối dữ liệu}

Nếu mỗi activation map có kích thước:

\[
    H' \times W'
\]

và có $K$ activation map, thì output của convolution layer có kích thước:

\[
    H' \times W' \times K
\]

Tức là output cũng là một tensor 3D.

% ==========================================================
\section{Pooling Layer -- Lớp pooling}

\subsection{Pooling layer là gì?}

Ngoài convolution layer, CNN thường có thêm một loại layer rất quan trọng là pooling layer.

Pooling layer nhận input là activation map và tạo ra output có kích thước nhỏ hơn.

Ý tưởng là chọn một cửa sổ nhỏ, ví dụ:

\[
    2\times2
\]

rồi trượt cửa sổ đó trên activation map.

\begin{ghinho}
Pooling là kỹ thuật giảm kích thước không gian của activation map.
\end{ghinho}

\subsection{Pooling window và stride}

Pooling layer cần hai thông số quan trọng:

\begin{itemize}
    \item \textbf{window size}: kích thước vùng pooling;
    \item \textbf{stride}: bước nhảy khi trượt cửa sổ.
\end{itemize}

Một lựa chọn phổ biến là:

\[
    \text{window size}=2\times2
\]

\[
    \text{stride}=2
\]

Khi đó, các vùng pooling thường không chồng lấn lên nhau.

\subsection{Ví dụ cửa sổ pooling $2\times2$}

Giả sử ta có activation map:

\[
    A =
    \begin{bmatrix}
    1 & 3 & 2 & 4\\
    5 & 6 & 1 & 2\\
    7 & 2 & 3 & 1\\
    0 & 4 & 5 & 8
    \end{bmatrix}
\]

Với pooling window $2\times2$ và stride $2$, ta chia thành bốn vùng:

\[
    \begin{bmatrix}
    1 & 3\\
    5 & 6
    \end{bmatrix},
    \quad
    \begin{bmatrix}
    2 & 4\\
    1 & 2
    \end{bmatrix}
\]

\[
    \begin{bmatrix}
    7 & 2\\
    0 & 4
    \end{bmatrix},
    \quad
    \begin{bmatrix}
    3 & 1\\
    5 & 8
    \end{bmatrix}
\]

% ==========================================================
\section{Max Pooling -- Max pooling}

\subsection{Ý tưởng max pooling}

Max pooling lấy giá trị lớn nhất trong mỗi vùng pooling.

Với vùng:

\[
    \begin{bmatrix}
    1 & 3\\
    5 & 6
    \end{bmatrix}
\]

giá trị max là:

\[
    6
\]

\subsection{Ví dụ max pooling}

Với activation map:

\[
    A =
    \begin{bmatrix}
    1 & 3 & 2 & 4\\
    5 & 6 & 1 & 2\\
    7 & 2 & 3 & 1\\
    0 & 4 & 5 & 8
    \end{bmatrix}
\]

Max pooling $2\times2$, stride $2$ cho kết quả:

\[
    P =
    \begin{bmatrix}
    6 & 4\\
    7 & 8
    \end{bmatrix}
\]

Vì:

\[
    \max
    \begin{bmatrix}
    1 & 3\\
    5 & 6
    \end{bmatrix}
    =
    6
\]

\[
    \max
    \begin{bmatrix}
    2 & 4\\
    1 & 2
    \end{bmatrix}
    =
    4
\]

\[
    \max
    \begin{bmatrix}
    7 & 2\\
    0 & 4
    \end{bmatrix}
    =
    7
\]

\[
    \max
    \begin{bmatrix}
    3 & 1\\
    5 & 8
    \end{bmatrix}
    =
    8
\]

\begin{ghinho}
Max pooling giữ lại tín hiệu mạnh nhất trong mỗi vùng. Đây là kỹ thuật pooling rất phổ biến.
\end{ghinho}

% ==========================================================
\section{Average Pooling -- Average pooling}

\subsection{Ý tưởng average pooling}

Average pooling lấy giá trị trung bình trong mỗi vùng pooling.

Với vùng:

\[
    \begin{bmatrix}
    1 & 3\\
    5 & 6
    \end{bmatrix}
\]

giá trị average pooling là:

\[
    \frac{1+3+5+6}{4}
    =
    3.75
\]

\subsection{Ví dụ average pooling}

Với activation map:

\[
    A =
    \begin{bmatrix}
    1 & 3 & 2 & 4\\
    5 & 6 & 1 & 2\\
    7 & 2 & 3 & 1\\
    0 & 4 & 5 & 8
    \end{bmatrix}
\]

Average pooling $2\times2$, stride $2$ cho kết quả:

\[
    P =
    \begin{bmatrix}
    3.75 & 2.25\\
    3.25 & 4.25
    \end{bmatrix}
\]

Vì:

\[
    \frac{1+3+5+6}{4}=3.75
\]

\[
    \frac{2+4+1+2}{4}=2.25
\]

\[
    \frac{7+2+0+4}{4}=3.25
\]

\[
    \frac{3+1+5+8}{4}=4.25
\]

\begin{luuy}
Max pooling thường được dùng phổ biến hơn trong nhiều CNN cổ điển, nhưng average pooling cũng xuất hiện trong một số kiến trúc.
\end{luuy}

% ==========================================================
\section{Pooling as Downsampling -- Pooling như một dạng downsampling}

\subsection{Downsampling là gì?}

Downsampling là quá trình giảm kích thước dữ liệu.

Trong ảnh, downsampling có thể hiểu là làm ảnh nhỏ lại.

Pooling cũng có tác dụng tương tự: nó làm activation map nhỏ lại.

Ví dụ:

\[
    10\times10
    \xrightarrow{\text{pooling }2\times2,\ stride\ 2}
    5\times5
\]

\subsection{Mục tiêu của pooling}

Pooling có các mục tiêu chính:

\begin{enumerate}[label=\arabic*.]
    \item Giảm kích thước activation map.
    \item Giảm số lượng tính toán.
    \item Giảm số tham số ở các lớp phía sau.
    \item Giúp mô hình bớt nhạy với thay đổi nhỏ về vị trí.
    \item Góp phần giảm overfitting.
\end{enumerate}

\begin{ghinho}
Mục tiêu lớn nhất của pooling là giảm dung lượng dữ liệu và giảm độ phức tạp của hệ thống.
\end{ghinho}

\subsection{Pooling có làm mất thông tin không?}

Có. Khi giảm kích thước, pooling có thể làm mất một phần thông tin.

Tuy nhiên, nếu hệ thống được thiết kế hợp lý, những thông tin quan trọng vẫn được giữ lại ở mức đủ dùng.

\begin{luuy}
Việc giảm kích thước ảnh hoặc activation map có thể làm mất thông tin. Tuy nhiên, trong nhiều bài toán, điều quan trọng là giữ đặc trưng lớn và tổng quát, không phải giữ mọi chi tiết nhỏ.
\end{luuy}

% ==========================================================
\section{Repeated Convolution and Pooling -- Lặp lại convolution và pooling}

\subsection{CNN thường có nhiều lớp}

Trong thực tế, CNN hiếm khi chỉ có một convolution layer duy nhất.

Thông thường, ta lặp lại nhiều lần các khối:

\[
    \text{convolution}
    \rightarrow
    \text{pooling}
\]

hoặc:

\[
    \text{convolution}
    \rightarrow
    \text{convolution}
    \rightarrow
    \text{pooling}
\]

Tùy kiến trúc, số lớp có thể là:

\begin{itemize}
    \item vài lớp;
    \item vài chục lớp;
    \item thậm chí hàng trăm lớp.
\end{itemize}

\subsection{Output của lớp trước là input của lớp sau}

Sau convolution và pooling, ta thu được một tensor mới.

Tensor này trở thành input cho convolution layer tiếp theo.

Ví dụ:

\[
    X
    \rightarrow
    \text{Conv}_1
    \rightarrow
    A^{(1)}
    \rightarrow
    \text{Pooling}
    \rightarrow
    P^{(1)}
    \rightarrow
    \text{Conv}_2
\]

\begin{ghinho}
Trong CNN, các activation map của lớp trước trở thành các channel đầu vào cho lớp convolution tiếp theo.
\end{ghinho}

% ==========================================================
\section{Convolution on Multiple Activation Maps -- Tích chập trên nhiều activation map}

\subsection{Input của lớp convolution sau có nhiều layer}

Giả sử sau lớp convolution và pooling đầu tiên, ta thu được 4 activation map.

Khi đó input cho lớp convolution tiếp theo có dạng:

\[
    H \times W \times 4
\]

Tức là input có 4 channel.

\subsection{Kernel của lớp sau cũng phải có 4 layer}

Vì input có 4 channel, kernel của lớp convolution tiếp theo cũng phải có 4 layer.

Nếu kernel có kích thước không gian $3\times3$, thì kernel đầy đủ có kích thước:

\[
    3\times3\times4
\]

Tức là nó gồm 4 ma trận $3\times3$.

\[
    K^{(1)}, K^{(2)}, K^{(3)}, K^{(4)}
\]

Mỗi layer của kernel áp lên một activation map tương ứng.

\subsection{Một kernel nhiều layer vẫn tạo ra một activation map}

Tương tự ảnh RGB, nếu input có 4 channel và kernel là $3\times3\times4$, thì một kernel vẫn chỉ tạo ra một activation map.

Nếu lớp đó có $K$ kernel thì sẽ tạo ra $K$ activation map.

\[
    K \text{ kernels}
    \rightarrow
    K \text{ output maps}
\]

\begin{ghinho}
Ở bất kỳ convolution layer nào, độ sâu của kernel phải bằng độ sâu của input.
\end{ghinho}

% ==========================================================
\section{Flatten and Fully Connected Layers -- Flatten và các lớp fully connected}

\subsection{Khi nào cần flatten?}

Sau nhiều lần convolution và pooling, ta thu được một tensor đặc trưng.

Ví dụ:

\[
    H' \times W' \times C'
\]

Để đưa tensor này vào mạng fully connected, ta cần biến nó thành một vector.

Thao tác này gọi là:

\[
    \text{flatten}
\]

\subsection{Ví dụ flatten}

Nếu tensor có kích thước:

\[
    5 \times 5 \times 16
\]

thì sau flatten, vector có số phần tử:

\[
    5\cdot5\cdot16=400
\]

Vector này được đưa vào fully connected layer.

\subsection{Fully connected layer ở cuối CNN}

Phần cuối của CNN thường gồm một hoặc vài fully connected layer.

Ví dụ:

\[
    \text{Flatten}
    \rightarrow
    \text{FC}_1
    \rightarrow
    \text{FC}_2
    \rightarrow
    \text{Output}
\]

Trong bài toán phân loại, số output cuối cùng thường bằng số class.

\begin{vidu}
Nếu bài toán có 1000 class, output cuối cùng thường có 1000 phần tử.
\end{vidu}

% ==========================================================
\section{Sliding Order -- Thứ tự trượt kernel}

\subsection{Trượt ngang trước hay dọc trước có khác không?}

Khi thực hiện convolution, ta thường minh họa bằng cách trượt kernel từ trái sang phải, rồi từ trên xuống dưới.

Tuy nhiên, về bản chất, ta có thể duyệt các vị trí theo thứ tự khác.

Ví dụ:

\begin{itemize}
    \item trượt ngang trước rồi xuống dòng;
    \item trượt dọc trước rồi sang cột;
    \item duyệt theo bất kỳ thứ tự nào.
\end{itemize}

Kết quả cuối cùng vẫn giống nhau nếu ta ghi giá trị output vào đúng vị trí tương ứng.

\begin{ghinho}
Thứ tự duyệt các vùng không quan trọng. Điều quan trọng là mỗi kết quả phải được ghi đúng vị trí trong activation map.
\end{ghinho}

\subsection{Sai lầm cần tránh}

Nếu trượt kernel ở vị trí bên phải nhưng lại ghi kết quả xuống dòng dưới, hoặc trượt xuống dưới nhưng lại ghi kết quả sang bên phải, thì activation map sẽ bị sai.

\begin{canhbao}
Khi tính convolution bằng tay, cần chú ý ánh xạ đúng giữa vị trí kernel trên input và vị trí output trong activation map.
\end{canhbao}

% ==========================================================
\section{Output Size Reduction -- Kích thước output giảm dần}

\subsection{Vì sao output có thể nhỏ dần?}

Nếu convolution không dùng padding, kích thước output thường nhỏ hơn input.

Ví dụ input:

\[
    3\times3
\]

kernel:

\[
    2\times2
\]

stride bằng 1, không padding, output là:

\[
    2\times2
\]

Vì kernel $2\times2$ chỉ có thể đặt được ở 4 vị trí hợp lệ trên ảnh $3\times3$.

\subsection{Pooling cũng làm kích thước nhỏ lại}

Pooling thường làm kích thước giảm mạnh hơn.

Ví dụ:

\[
    10\times10
    \xrightarrow{\text{pooling }2\times2,\ stride\ 2}
    5\times5
\]

Vì vậy, nếu lặp quá nhiều convolution và pooling, kích thước không gian của feature map có thể trở nên rất nhỏ.

\begin{luuy}
Trong thực tế, kiến trúc CNN phải được thiết kế hợp lý để kích thước feature map không giảm quá nhanh hoặc quá mức cần thiết.
\end{luuy}

% ==========================================================
\section{Zero Padding -- Đệm số 0}

\subsection{Zero padding là gì?}

Zero padding là kỹ thuật thêm một hoặc nhiều lớp số 0 xung quanh input.

Ví dụ, input $3\times3$:

\[
    \begin{bmatrix}
    x_{11} & x_{12} & x_{13}\\
    x_{21} & x_{22} & x_{23}\\
    x_{31} & x_{32} & x_{33}
    \end{bmatrix}
\]

Nếu thêm padding một lớp số 0 xung quanh, ta được:

\[
    \begin{bmatrix}
    0 & 0 & 0 & 0 & 0\\
    0 & x_{11} & x_{12} & x_{13} & 0\\
    0 & x_{21} & x_{22} & x_{23} & 0\\
    0 & x_{31} & x_{32} & x_{33} & 0\\
    0 & 0 & 0 & 0 & 0
    \end{bmatrix}
\]

\subsection{Mục đích của padding}

Zero padding được dùng để:

\begin{enumerate}[label=\arabic*.]
    \item kiểm soát kích thước output;
    \item giúp feature map không bị nhỏ quá nhanh;
    \item cho phép kernel xử lý tốt hơn các pixel ở biên ảnh;
    \item giữ lại thông tin ở vùng rìa ảnh.
\end{enumerate}

\begin{ghinho}
Zero padding là cách thêm số 0 quanh input để điều chỉnh kích thước output sau convolution.
\end{ghinho}

\subsection{Công thức kích thước output}

Với input kích thước $H\times W$, kernel $F\times F$, padding $P$, stride $S$, kích thước output là:

\[
    H_{\text{out}}
    =
    \left\lfloor
    \frac{H+2P-F}{S}
    \right\rfloor
    +1
\]

\[
    W_{\text{out}}
    =
    \left\lfloor
    \frac{W+2P-F}{S}
    \right\rfloor
    +1
\]

Nếu muốn giữ kích thước output bằng input, ta cần chọn padding phù hợp.

\subsection{Lưu ý với kernel chẵn}

Với kernel có kích thước chẵn, ví dụ $2\times2$, việc giữ nguyên chính xác kích thước input bằng padding đối xứng không phải lúc nào cũng đơn giản.

Ví dụ, input $3\times3$, kernel $2\times2$, stride $1$:

\[
    H_{\text{out}} = H
\]

sẽ cần:

\[
    \frac{H+2P-2}{1}+1=H
\]

Suy ra:

\[
    2P=1
\]

Tức là padding đối xứng nguyên không thực hiện được. Khi đó, một số hệ thống có thể dùng padding bất đối xứng hoặc quy ước riêng.

\begin{luuy}
Ý chính cần nhớ là padding giúp kiểm soát kích thước output. Với kernel lẻ như $3\times3$, padding đối xứng dễ dùng hơn, ví dụ padding $P=1$ giúp giữ kích thước khi stride bằng 1.
\end{luuy}

% ==========================================================
\section{Technical Glossary -- Bảng thuật ngữ chuyên ngành}

\small
\begin{longtable}{p{0.27\textwidth}p{0.48\textwidth}p{0.18\textwidth}}
\toprule
\textbf{Thuật ngữ} & \textbf{Giải thích} & \textbf{Ví dụ} \\
\midrule
\endfirsthead

\toprule
\textbf{Thuật ngữ} & \textbf{Giải thích} & \textbf{Ví dụ} \\
\midrule
\endhead

Channel & Kênh dữ liệu của ảnh hoặc feature map. & RGB có 3 channel. \\

RGB & Hệ màu gồm ba kênh đỏ, xanh lá, xanh dương. & Ảnh $224\times224\times3$. \\

Kernel & Bộ trọng số nhỏ dùng để tích chập trên input. & $3\times3\times3$. \\

Kernel depth & Độ sâu của kernel, phải bằng số channel input. & Input 4 channel thì kernel depth 4. \\

Activation map & Ma trận output do một kernel tạo ra. & Output $10\times10$. \\

Feature map & Cách gọi khác của activation map. & Bản đồ đặc trưng. \\

Convolution layer & Lớp tích chập dùng kernel để trích xuất đặc trưng. & Conv $3\times3$. \\

Pooling layer & Lớp giảm kích thước activation map. & Max pooling. \\

Max pooling & Lấy giá trị lớn nhất trong mỗi vùng pooling. & Max của vùng $2\times2$. \\

Average pooling & Lấy giá trị trung bình trong mỗi vùng pooling. & Average của vùng $2\times2$. \\

Window size & Kích thước cửa sổ pooling. & $2\times2$. \\

Stride & Bước nhảy khi trượt kernel hoặc pooling window. & Stride $2$. \\

Downsampling & Giảm kích thước dữ liệu. & $10\times10$ thành $5\times5$. \\

Flatten & Biến tensor nhiều chiều thành vector. & $5\times5\times16$ thành vector 400. \\

Fully connected layer & Lớp kết nối đầy đủ ở cuối CNN. & Lớp phân loại. \\

Zero padding & Thêm viền số 0 quanh input. & Padding $P=1$. \\

Padding & Đệm thêm quanh input để điều chỉnh output size. & Same padding. \\

Output size & Kích thước đầu ra sau convolution hoặc pooling. & $H_{\text{out}}\times W_{\text{out}}$. \\

\bottomrule
\end{longtable}
\normalsize

% ==========================================================
\section{Key Concepts Summary -- Tóm tắt kiến thức trọng tâm}

\subsection{Các ý chính cần nhớ}

\begin{enumerate}
    \item Ảnh màu RGB có 3 channel.
    \item Khi input có nhiều channel, kernel cũng phải có cùng số layer với input.
    \item Với input RGB, một kernel $3\times3$ thực chất có kích thước $3\times3\times3$.
    \item Một kernel nhiều layer vẫn chỉ tạo ra một activation map.
    \item Nếu một convolution layer có $K$ kernel, nó tạo ra $K$ activation map.
    \item Pooling layer dùng để giảm kích thước activation map.
    \item Max pooling lấy giá trị lớn nhất trong mỗi vùng.
    \item Average pooling lấy giá trị trung bình trong mỗi vùng.
    \item Pooling là một dạng downsampling.
    \item Pooling giúp giảm dung lượng dữ liệu, giảm tính toán và giảm số tham số.
    \item CNN thường lặp lại nhiều lần convolution và pooling.
    \item Output của layer trước trở thành input nhiều channel cho layer sau.
    \item Trước khi đưa vào fully connected layer, tensor thường được flatten thành vector.
    \item Zero padding thêm viền số 0 quanh input để kiểm soát kích thước output.
    \item Thứ tự trượt kernel không quan trọng nếu ghi output đúng vị trí.
\end{enumerate}

\subsection{Công thức quan trọng}

\subsection*{Số tham số của một kernel nhiều kênh}

\[
    K_h \cdot K_w \cdot C
\]

Trong đó $C$ là số channel của input.

\subsection*{Convolution nhiều kênh}

\[
    a_{ij}
    =
    f
    \left(
    \sum_{c=1}^{C}
    \sum_{u=1}^{K_h}
    \sum_{v=1}^{K_w}
    K_{u,v,c}
    X_{i+u-1,j+v-1,c}
    +
    b
    \right)
\]

\subsection*{Pooling $2\times2$, stride 2}

\[
    H\times W
    \rightarrow
    \frac{H}{2}\times \frac{W}{2}
\]

nếu $H,W$ chia hết cho 2.

\subsection*{Kích thước output sau convolution}

\[
    H_{\text{out}}
    =
    \left\lfloor
    \frac{H+2P-F}{S}
    \right\rfloor
    +1
\]

\[
    W_{\text{out}}
    =
    \left\lfloor
    \frac{W+2P-F}{S}
    \right\rfloor
    +1
\]

% ==========================================================
\section{Review Questions -- Câu hỏi ôn tập}

\begin{enumerate}
    \item Ảnh màu RGB có bao nhiêu channel?
    \item Nếu input có 3 channel, kernel phải có bao nhiêu layer?
    \item Một kernel $3\times3\times3$ có bao nhiêu trọng số?
    \item Một kernel tạo ra bao nhiêu activation map?
    \item Nếu một convolution layer có 32 kernel, output có bao nhiêu activation map?
    \item Pooling layer dùng để làm gì?
    \item Max pooling khác average pooling như thế nào?
    \item Với activation map $10\times10$, pooling $2\times2$, stride 2 tạo ra output kích thước bao nhiêu?
    \item Downsampling là gì?
    \item Vì sao pooling có thể làm mất thông tin?
    \item Vì sao vẫn dùng pooling dù nó có thể làm mất thông tin?
    \item Khi output của layer trước có 4 activation map, kernel của layer sau phải có depth bằng bao nhiêu?
    \item Flatten là gì?
    \item Zero padding là gì?
    \item Vì sao cần zero padding?
    \item Thứ tự trượt kernel có làm thay đổi kết quả không nếu ghi output đúng vị trí?
\end{enumerate}

% ==========================================================
\section{Practice Exercises -- Bài tập vận dụng}

\subsection*{Bài tập 1: Kích thước kernel nhiều kênh}
\addcontentsline{toc}{section}{Bài tập 1: Kích thước kernel nhiều kênh}

Input là ảnh RGB. Ta dùng kernel có kích thước không gian $5\times5$.

Hỏi kernel đầy đủ có kích thước bao nhiêu và có bao nhiêu trọng số?

\subsection*{Lời giải gợi ý}

Ảnh RGB có 3 channel, nên kernel phải có depth bằng 3.

Kích thước kernel là:

\[
    5\times5\times3
\]

Số trọng số:

\[
    5\cdot5\cdot3=75
\]

\subsection*{Bài tập 2: Số activation map}
\addcontentsline{toc}{section}{Bài tập 2: Số activation map}

Một convolution layer có 64 kernel. Hỏi output của layer này có bao nhiêu activation map?

\subsection*{Lời giải gợi ý}

Mỗi kernel tạo ra một activation map.

Do đó 64 kernel tạo ra:

\[
    64
\]

activation map.

\subsection*{Bài tập 3: Max pooling}
\addcontentsline{toc}{section}{Bài tập 3: Max pooling}

Cho activation map:

\[
    A =
    \begin{bmatrix}
    1 & 4 & 2 & 3\\
    5 & 6 & 1 & 0\\
    2 & 8 & 7 & 1\\
    3 & 4 & 5 & 9
    \end{bmatrix}
\]

Dùng max pooling $2\times2$, stride 2. Tính output.

\subsection*{Lời giải gợi ý}

Các vùng pooling là:

\[
    \begin{bmatrix}
    1 & 4\\
    5 & 6
    \end{bmatrix}
    \Rightarrow 6
\]

\[
    \begin{bmatrix}
    2 & 3\\
    1 & 0
    \end{bmatrix}
    \Rightarrow 3
\]

\[
    \begin{bmatrix}
    2 & 8\\
    3 & 4
    \end{bmatrix}
    \Rightarrow 8
\]

\[
    \begin{bmatrix}
    7 & 1\\
    5 & 9
    \end{bmatrix}
    \Rightarrow 9
\]

Vậy output là:

\[
    P =
    \begin{bmatrix}
    6 & 3\\
    8 & 9
    \end{bmatrix}
\]

\subsection*{Bài tập 4: Average pooling}
\addcontentsline{toc}{section}{Bài tập 4: Average pooling}

Dùng lại activation map ở bài tập 3. Tính average pooling $2\times2$, stride 2.

\subsection*{Lời giải gợi ý}

\[
    \frac{1+4+5+6}{4}=4
\]

\[
    \frac{2+3+1+0}{4}=1.5
\]

\[
    \frac{2+8+3+4}{4}=4.25
\]

\[
    \frac{7+1+5+9}{4}=5.5
\]

Vậy output là:

\[
    P =
    \begin{bmatrix}
    4 & 1.5\\
    4.25 & 5.5
    \end{bmatrix}
\]

\subsection*{Bài tập 5: Kích thước output sau convolution}
\addcontentsline{toc}{section}{Bài tập 5: Kích thước output sau convolution}

Input có kích thước $32\times32$, kernel $3\times3$, stride $1$, padding $1$.

Tính kích thước output.

\subsection*{Lời giải gợi ý}

\[
    H_{\text{out}}
    =
    \left\lfloor
    \frac{32+2\cdot1-3}{1}
    \right\rfloor
    +1
\]

\[
    H_{\text{out}}
    =
    32
\]

Tương tự:

\[
    W_{\text{out}}=32
\]

Vậy output có kích thước:

\[
    32\times32
\]

\subsection*{Bài tập 6: Flatten}
\addcontentsline{toc}{section}{Bài tập 6: Flatten}

Sau nhiều lớp CNN, ta thu được tensor kích thước:

\[
    7\times7\times128
\]

Hỏi sau flatten, vector có bao nhiêu phần tử?

\subsection*{Lời giải gợi ý}

\[
    7\cdot7\cdot128=6272
\]

Vector sau flatten có 6272 phần tử.

% ==========================================================
\section*{Conclusion -- Kết luận}
\addcontentsline{toc}{section}{Kết luận}

CNN2 mở rộng phép convolution từ trường hợp ảnh xám một kênh sang ảnh màu nhiều kênh. Khi input có nhiều channel, kernel cũng phải có cùng số layer tương ứng. Ví dụ, ảnh RGB có 3 channel thì kernel $3\times3$ thực chất là kernel $3\times3\times3$. Tại mỗi vị trí, ta nhân tương ứng trên tất cả các channel, cộng toàn bộ các tích, cộng bias và đưa qua hàm kích hoạt để tạo ra một giá trị trong activation map.

Bài học cũng giới thiệu pooling layer, một thành phần quan trọng trong CNN. Max pooling lấy giá trị lớn nhất trong mỗi vùng, còn average pooling lấy giá trị trung bình. Pooling là một dạng downsampling, giúp giảm kích thước dữ liệu, giảm số lượng tính toán và giảm số tham số ở các lớp phía sau.

Trong CNN thực tế, convolution và pooling thường được lặp lại nhiều lần. Output của lớp trước trở thành input nhiều channel cho lớp sau, nên kernel ở lớp sau cũng phải có depth bằng số channel của input. Cuối cùng, tensor đặc trưng được flatten thành vector rồi đưa vào fully connected layer để phân loại.

Ngoài ra, bài học cũng nhắc đến zero padding, tức là thêm viền số 0 quanh input để kiểm soát kích thước output và giúp feature map không bị nhỏ quá nhanh sau nhiều lớp convolution.

\end{document}
